% Part1_Qubit_era.md — sanitized presentation layer (generated; see source/CONVERSION_LOG.md)
\chapter{\texorpdfstring{$\langle qu(b(i(t|t))ick)\rangle$ — \textbf{Part 1 · Qubit\_era} — Axioms · Conventions · Qubit-Era Definitions · Pair Canon · Ledger Dynamics}{Qubitick - **Part 1 - Qubit _era** - Axioms - Conventions - Qubit-Era Definitions - Pair Canon - Ledger Dynamics}}

\textbf{0.1 Corpus map.} The canon is the era chain

%dsp-level:0
\[
\textbf{Part 1 · Qubit\_era}\ \longrightarrow\ \textbf{Part 2 · Bosonic\_era}\ \longrightarrow\ \textbf{Part 3 · Fermion\_era}\ \longrightarrow\ \textbf{Part 4 · Baryomesonic\_era},
\]

each part citing only its priors, itself (earlier tags), and the Census annex (the computed data cells, cited row-wise by tag). Part 1 (this chapter): §A axioms A1–A9 · §B conventions C1–C13 · §C the definition layer D0.1–D4.3 (the corpus-wide vocabulary root: Tiers 0–1 qubit-era, Tiers 2–4 naming the structures the bosonic era derives) · §D master symbol table · §E calibration policy P1–P5 · §F foundations (CT0.1–CT0.14, CT1.0–CT1.11b, L0–L5). Part 2 (section letters preserve the classical canon's): §C CT1.12–CT2.11 (geometry), §D EW.1–EW.36a (electroweak), §E EM.1–EM.26 (electromagnetic), §F CT3.1–CT3.22a (colour), §G CT4.1–CT4.13 (saturation), §Cal-2 (registry note). Part 3: Def-H.1–Def-H.43, FT.1–FT.114, FT.116–FT.131, FT.133–FT.139, FT.141–FT.145 (structural tags), FT.153–FT.160. Part 4: Def-BM.1–Def-BM.20, BM.1–BM.152. Census: the 25 data-cell tags (Category B of the separation charter).

\textbf{0.2 Citation closure.} $\forall$ claims in Parts 1–4 cite only: A1–A9, C1–C13, D0.1–D4.3, P1–P5 (Parts 1–2), previously proved tags, and Census cells (row-wise). $\neg\exists$ imported external physical law, constant, or measured value. Physics vocabulary $:=$ naming of derived structures only. External comparison data enter only quarantined audit rows (Census FT.114, FT.120; P1, P5).

\textbf{0.2a Math \& hypothesis policy.} Standard mathematics (linear algebra, representation theory, tensor calculus, classification theorems) is cited as $[\text{math}]$-class tags in \textbf{Req} and is not part of the physical axioms. Conditional theorems carry their hypotheses explicitly in \textbf{Thm} and \textbf{Req}. The complete list of extra-axiomatic inputs is empty: CT2.6 (ii), CT2.6 (iii) — the two continuum-limit conditions — are derived theorems of Part 4 (Tier BM.24, BM.151): (ii) is witnessed by the ledger's own unbounded ladder, (iii) by the density-matrix spectral bound on the D2.1 coefficient fields; no others exist. The declarations at use stand as scope-inheritance markers over the same loci. CT2.6 (i) is a definitional precondition: the valid-universe condition of D2.11 in application (“there is something”), $\neg$ an extra-axiomatic hypothesis. CT2.8 carries no hypotheses: its reflection and factorization clauses are derived (D2.7, D1.9, D1.3, D2.8). The edge phase is registration data: its alphabet is derived (CT0.9) and its per-edge value is determined at crystallization (Def-EM.20, EM.20). \textbf{Integration note.} The integer-edge hypothesis of CT1.29a is \textit{not} independent: it is realized on materialized closed chains of the fully antialigned class (CT1.19b(3)). The ambient-richness condition in the CT3.22a Rmk is a witness scope for that Rmk's non-locality construction only; no theorem clause of Parts 1–2 rests on it. The materialization domain of CT1.19b is CT2.5’s own pole condition (no new input); its domain-generalization step from EW.17’s wrapping chains is flagged in its Pf. The per-edge independence hypothesis of CT4.5a is an ensemble scope for that theorem alone — its dynamics realization (a crystallization-rate datum) is NOT declared in Parts 1–2. The crystallization–tilt coupling of L0 (a seal-time state update: one Gibbs tilt step per A4 event) and the finite-incidence scope of L4 are theorem-local scopes of the same species, declared at use. The interface definitions of Part 3 (Def-H.35–Def-H.43) carry their own non-derived clauses explicitly.

\textbf{0.3 Free data.} $\exists!\,$ free dimensionful scale: $\mathfrak h$ (D2.9, P1). Every other quantity: dimensionless-derived, or a power of $\mathfrak h$. Free internal configuration parameters (P4-class, $\neg$constants of the theory): the P4 core (§E) plus the era registries Cal-2 (Part 2), Cal-F (Part 3), Cal-BM (Part 4).

\textbf{0.4 Theorem format (all parts).} Per tag: \textbf{Thm} — statement, purely symbolic · \textbf{Req} — \{axioms, conventions, definitions, prior theorems\} used · \textbf{Pf} — derivation, purely symbolic, justifications in $[\cdot]$, end $\blacksquare$ · \textbf{Phys} — short English summary of the theorem and its physics. Corollaries: \textbf{Cor}. Definitions: \textbf{Def.} blocks; scope conditions and constitutive clauses are carried inside the \textbf{Def.}/\textbf{Thm.} text. Forward pointers (scope conditions cited at use) are marked “forward pointer” at the citing site.

\textbf{0.5 Symbol reconciliations.} $P_M\equiv\mathcal P_M$ (C7); $F_Z\equiv F_{W^3}$ (C8); $\tau(v)\in\Xi_n$ (type set) vs $v\in V_n$ (vertex set), D0.8.

\textbf{0.6 Machine verification.} \texttt{checks/\allowbreak{}}: dependency-free exact-arithmetic verification suite (\texttt{python3\ checks/\allowbreak{}run\_\allowbreak{}all.py}). \texttt{verify\_\allowbreak{}era\_\allowbreak{}structure.py}: era-chain law, citation closure, census separation, tag ownership, declaration hygiene, \textbf{Thm·Req·Pf·Phys} completeness, index regeneration. \texttt{verify\_\allowbreak{}B.py}–\texttt{verify\_\allowbreak{}G.py} and the \texttt{verify\_\allowbreak{}*} content suite: algebraic/geometry/arithmetical identities, exact, no floating point. See \texttt{checks/\allowbreak{}README.md}.

\phantomsection
\section[§A — Axioms]{\texorpdfstring{§A — Axioms}{SA - Axioms}}

\begin{axiom}[{\texorpdfstring{A1}{A1}}]\label{tag:A1}
%dsp-level:0
\[
\mathcal Q=\{q_n\}_{n\ge1};
\]

%dsp-level:0
\[
\qquad \mathcal H_q\cong\mathbb C^2;
\]

%dsp-level:0
\[
\qquad \forall(p,q)\in\mathcal Q\times\mathcal Q,\ p\ne q:\quad
\rho_{pq}\in\mathfrak H(\mathcal H_p\otimes\mathcal H_q),\ \ \rho_{pq}\succeq0,\ \ \mathfrak T(\rho_{pq})=1 .
\]

\end{axiom}

\begin{axiom}[{\texorpdfstring{A2}{A2}}]\label{tag:A2}
%dsp-level:0
\[
\forall(p,q):\qquad C(p,q)=1\,\mathfrak q .
\]

\end{axiom}

\begin{axiom}[{\texorpdfstring{A3}{A3}}]\label{tag:A3}
%dsp-level:0
\[
\mathcal H_p\otimes\mathcal H_q=\vee^2\oplus\wedge^2;
\]

%dsp-level:0
\[
S(\sigma_\mu\otimes\sigma_\nu)=\sigma_\nu\otimes\sigma_\mu,\quad S^2=\mathbb1;\qquad
T(\sigma_\mu\otimes\sigma_\nu)=\eta^\mu\eta^\nu\,(\sigma_\mu\otimes\sigma_\nu),\quad
T^2=\mathbb1;
\]

%dsp-level:0
\[
\eta\in\{\pm1\}^4;\qquad
\exists!\,\mathfrak o\in\{0,1,2,3\}:\ \ \eta^{\mathfrak o}=-\eta^{\mu}\ \ \forall\mu\ne\mathfrak o .
\]

\end{axiom}

\begin{axiom}[{\texorpdfstring{A4}{A4}}]\label{tag:A4}
%dsp-level:0
\[
(p,q)\in\mathcal R\iff\bigl|\Pi_{\mathcal V}\rho_{pq}\bigr|\le\varepsilon_q\ \ \vee\ \ \bigl|\Pi_{\mathcal V}\rho_{pq}\bigr|\ge1-\varepsilon_q;\qquad
n_{\mathcal R}(p)=\bigl|\{q:(p,q)\in\mathcal R\}\bigr|
\]

%dsp-level:0
\[
\alpha_{pq}:=\begin{cases}0&\bigl|\Pi_{\mathcal V}\rho_{pq}\bigr|\le\varepsilon_q\\[1mm] \pi&\bigl|\Pi_{\mathcal V}\rho_{pq}\bigr|\ge1-\varepsilon_q\end{cases};\qquad
\alpha_{pq}\in\{0,\pi\}\quad\forall(p,q)\in\mathcal R;
\]

%dsp-level:0
\[
\exists\,\rho_{pq}:\ \bigl|\Pi_{\mathcal V}\rho_{pq}\bigr|\le\varepsilon_q;\qquad
\exists\,\rho_{p'q'}:\ \bigl|\Pi_{\mathcal V}\rho_{p'q'}\bigr|\ge1-\varepsilon_q .
\]

\end{axiom}

\begin{axiom}[{\texorpdfstring{A5}{A5}}]\label{tag:A5}
%dsp-level:0
\[
\rho_\circ=\tfrac14 I;\qquad
\rho\mapsto U\rho U^\dagger,\quad U\sim\mu_\star\bigl(\mathrm{SU}(2)\bigr),\quad U\notin\mathfrak L .
\]

\end{axiom}

\begin{axiom}[{\texorpdfstring{A6}{A6}}]\label{tag:A6}
%dsp-level:0
\[
\mu_\star:=\mu_{\mathrm{Haar}}\bigl(\mathrm{SU}(2)\bigr);\qquad
U\sim\mu_\star;\qquad
U\notin\mathfrak L;\qquad
\mathcal T(\rho):=U\rho U^\dagger;\qquad
\mathcal T^{-1}(\rho)=U^\dagger\rho U
\]

%dsp-level:0
\[
\mathfrak T\bigl(\mathcal T(\rho)\bigr)=\mathfrak T(\rho);\qquad
\mathcal T^{2}\ne\mathbb1;\qquad
\tau_n\to\tau_{n+1}:\ \ \rho\mapsto\mathcal T(\rho);\qquad
\neg\exists\,\tau^{-1}
\]

%dsp-level:0
\[
\bigl|\mathfrak L(\tau_{n+1})\bigr|\ \ge\ \bigl|\mathfrak L(\tau_n)\bigr|;\qquad
\bigl|\mathfrak L(\tau_{n+1})\bigr|\ =\ \bigl|\mathfrak L(\tau_n)\bigr|\,+\,n_{\mathrm{real\_out}}(\tau_n) .
\]

\end{axiom}

\begin{axiom}[{\texorpdfstring{A7}{A7}}]\label{tag:A7}
%dsp-level:0
\[
\forall n\ge0:\quad\Gamma=(V,E),\qquad E\subseteq V\times V
\]

%dsp-level:0
\[
\tau(v):=\bigl(\partial^{\leftarrow}(v),\partial^{\rightarrow}(v)\bigr)\in\Xi_n;\qquad
e\in E_{\mathcal R}\mapsto\phi_e;\qquad
e\in E_{\mathcal R}\mapsto\alpha_e\in\{0,\pi\}
\]

%dsp-level:0
\[
\mathfrak L=\bigcup_{n\ge0}\mathfrak L(\tau_n);
\]

%dsp-level:0
\[
\qquad \mathfrak L(\tau_n)\subseteq\mathfrak L(\tau_{n+1});
\]

%dsp-level:0
\[
\qquad \bigl|\mathfrak L(\tau_0)\bigr|=0;
\]

%dsp-level:0
\[
\qquad \bigl|\mathfrak L(\tau_n)\bigr|=\sum_{k<n}n_{\mathrm{real\_out}}(\tau_k) .
\]

\end{axiom}

\begin{axiom}[{\texorpdfstring{A8}{A8}}]\label{tag:A8}
%dsp-level:0
\[
\mathbf F=\mathbf F_{(+,\mathcal R)}\oplus\mathbf F_{(-,\mathcal R)}\oplus\mathbf F_{(+,\mathcal V)}\oplus\mathbf F_{(-,\mathcal V)}\oplus\mathbf F_{\mathrm{st}};
\]

%dsp-level:0
\[
\mathbf F_{(+,\mathcal R)}=F_G\oplus F_H\oplus F_{D^+};\qquad
\mathbf F_{(-,\mathcal R)}=F_{W^+}\oplus F_{W^-}\oplus F_Z;\qquad
\mathbf F_{(+,\mathcal V)}=F_{J^+};
\]

%dsp-level:0
\[
\mathbf F_{(-,\mathcal V)}=F_{J^-};\qquad
\mathbf F_{\mathrm{st}}=F_{\mathrm{EM}}\oplus F_{\mathrm{strong}}\oplus F_{D^-};
\]

%dsp-level:0
\[
F_{D^+}=\Lambda\,\mathbb I;\qquad
F_{D^-}:=\bigoplus_{n\ge4}F_n;\qquad
\mathbf F_M=P_M\,\mathbf F\,P_M;
\]

%dsp-level:0
\[
w(C)=\mathcal J(C)\!\!\prod_{e\in E(C)}\!\!G(e)\ \cdot\ (-1)^{\aleph(\sigma,C)};\qquad
Z(\mathcal R)=\sum_{C\in\mathcal R}\ \sum_\sigma w(C) .
\]

\end{axiom}

\begin{axiom}[{\texorpdfstring{A9}{A9}}]\label{tag:A9}
%dsp-level:0
\[
\mathcal U_\odot\in\mathfrak S\bigl(Z[\mathbf F_\oplus,M_\oplus]\bigr);\qquad
\mathcal U_\odot\notin\mathfrak S\implies\neg\bigl(\mathrm{A1}\wedge\cdots\wedge\mathrm{A9}\bigr) .
\]

$\blacksquare$

\end{axiom}

\phantomsection
\section[§B — Conventions]{\texorpdfstring{§B — Conventions}{SB - Conventions}}

\begin{convention}[{\texorpdfstring{C1}{C1}}]\label{tag:C1}
$\mu,\nu,\rho,\ldots\in\{0,1,2,3\}$: Pauli basis $\sigma_0=I_2,\sigma_1,\sigma_2,\sigma_3$ / emergent spacetime coordinates where a metric exists; $i,j,k,\ldots\in\{1,2,3\}$: spatial/Pauli-spatial; $a,b,c,\ldots\in\{1,\dots,8\}$: colour adjoint. Repeated indices: summed.

\end{convention}

\begin{convention}[{\texorpdfstring{C2}{C2}}]\label{tag:C2}
$T_A^\dagger=T_A$; $\ [T_A,T_B]=i\,f_{AB}{}^{C}\,T_C$; $\ f_{AB}{}^{C}\in\mathbb R$.

\end{convention}

\begin{convention}[{\texorpdfstring{C3}{C3}}]\label{tag:C3}
$\mathcal E_{st}:=\{\,X\in\mathfrak H(\mathbb C^4):\ SX=sX\ \wedge\ TX=tX\,\}$, $s,t\in\{+,-\}$. Symbol $\wedge$: antisymmetrization only ($\wedge^2$, $\varepsilon$-contractions); $\neg$ a sector-label conjunct.

\end{convention}

\begin{convention}[{\texorpdfstring{C4}{C4}}]\label{tag:C4}
$\varkappa$: sector-to-metric embedding scale (D2.1); $\kappa$: gravitational coupling (D2.9, CT2.9); $\varkappa\ne\kappa$; $\neg$ interchangeable.

\end{convention}

\begin{convention}[{\texorpdfstring{C5}{C5}}]\label{tag:C5}
$\Sigma:=\sum_i\sigma_i\otimes\sigma_i$ (operator); $\Sigma_{\mathrm{tr}}:=\delta^{ij}q_{ij}\in\mathbb R$ (scalar, D2.1); $\neg$ interchangeable.

\end{convention}

\begin{convention}[{\texorpdfstring{C6}{C6}}]\label{tag:C6}
$\mathsf E_{\mathrm{tr}}$: transit entropy; $\mathsf E_{\mathrm{bd}}$: bound entropy (D1.5); $\neg$ squared-field notation for either.

\end{convention}

\begin{convention}[{\texorpdfstring{C7}{C7}}]\label{tag:C7}
$P_M\equiv\mathcal P_M$ (A8). Rung-resolved: $\mathcal P_M^{(n)}$ (D2.3); $\ \mathbf F_M^{(n)}=\mathcal P_M^{(n)}\,\mathbf F\,\mathcal P_M^{(n)}$.

\end{convention}

\begin{convention}[{\texorpdfstring{C8}{C8}}]\label{tag:C8}
$F_Z\equiv F_{W^3}$: neutral-weak aspect of A8; $\ F_W:=F_{W^+}\oplus F_{W^-}\oplus F_Z$; pre-mixing seed generator: $W_3$. Bare $Z$ in a spectrum: physical neutral eigenstate (D3.10, EW.24); $Z(\cdot)$ with argument: partition function (A8) / masked resolvent (D2.4).

\end{convention}

\begin{convention}[{\texorpdfstring{C9}{C9}}]\label{tag:C9}
Entropy in nats: $S=\ln\Omega$; \_bits\_ only where explicit. One real edge $=1$ bit $=\ln2$ nats (CT0.9, CT1.29).

\end{convention}

\begin{convention}[{\texorpdfstring{C10}{C10}}]\label{tag:C10}
Register accounting (reality principle): registrations convert quantum data to classical alternatives at the ledger. Registering a classical alternative set $X$ for single-shot linear readout ($\mathfrak L$ must distinguish which value was carried, D0.7) assigns the alternatives mutually orthogonal register directions: register cost $=|X|$. Shannon content $=\log_2|X|$ bits is a distinct quantity: source-availability register count $\ne$ channel bit throughput. One channel quantum supplies $C_{\mathrm{cap}}=\dim_{\mathbb R}\mathfrak H(\mathcal H_p\otimes\mathcal H_q)=16$ register directions (A2, D0.9). Ledger-level registration data are real: edge transport acts on the reality line by $\{\pm1\}$ (D1.1, CT0.9).

\end{convention}

\begin{convention}[{\texorpdfstring{C11}{C11}}]\label{tag:C11}
A5/A6 reading: the conjugation clause is node-level — $\rho\mapsto U\rho U^\dagger$ [A6: $\mathcal T$] with $\rho\in\mathfrak H(\mathcal H_q)\cong\mathfrak H(\mathbb C^2)$ and $U\sim\mu_\star(\mathrm{SU}(2))$ ($2\times2$: well-typed); $\rho_\circ=\tfrac14I$ is the vacuum pair state ($4\times4$). Node tick lifecycle: wait $\to$ data in $\to$ Haar draw/tick [A6: $\tau_n\to\tau_{n+1}$] (virtual states as basis) $\to$ data out/hop $\to$ wait. Draws are not retained: $U\notin\mathfrak L$ [A5, A6].

\end{convention}

\begin{convention}[{\texorpdfstring{C12}{C12}}]\label{tag:C12}
In A8's $Z(\mathcal R)=\sum_{C\in\mathcal R}\sum_\sigma w(C)$ the summation domain reads $\mathfrak C(\mathcal R)$ (D0.10): the configuration set on the real-edge set. $\mathcal R$ itself (A4, D0.1) is a set of pairs and is never a summation domain.

\end{convention}

\begin{convention}[{\texorpdfstring{C13}{C13}}]\label{tag:C13}
Signature convention. A3 fixes $T$ only through the products $\eta^\mu\eta^\nu$, and fixes $\eta$ only up to its odd slot $\mathfrak o$ and its overall sign. The corpus takes $\mathfrak o:=0$ (the identity slot $\sigma_0=I_2$, C1) and $\eta^{0}:=-1$, i.e. $\eta:=(-,+,+,+)$ (mostly-plus; D2.1: $\bar g_{00}:=-1$). The overall sign does not change $T$ [CT1.0(2)]; the three choices $\mathfrak o\in\{1,2,3\}$ are conjugate under $\mathrm{SU}(2)$ relabelling of the spatial Pauli axes [CT1.0(4)]; $\mathfrak o=0$ versus $\mathfrak o\ge1$ is the choice of the $\mathrm{SU}(2)$-covariant representative [CT1.0(5)], on which CT1.2's generator table rests [CT1.0 Rmk]. Uses of the representative carry the tag C13. $\neg$ an axiom.

\end{convention}

\phantomsection
\section[§C — The definition layer]{\texorpdfstring{§C — The definition layer}{SC - The definition layer}}

Definitions $:=$ vocabulary stipulations. $\forall$ content-bearing claims: theorems of Parts 1–4. Tiers 0–1 are qubit-era vocabulary; Tiers 2–4 name the structures the bosonic era derives — the definition layer is the corpus-wide vocabulary root (the classical canon's A §Part 3).

\phantomsection
\subsection[Tier 0 — bookkeeping]{\texorpdfstring{Tier 0 — bookkeeping}{Tier 0 - bookkeeping}}

\begin{definition}[{\texorpdfstring{D0.1}{D0.1}}]\label{tag:D0-1}
$\mathcal R:=\{(p,q):\ \mathrm{A4}\ \text{holds}\}$; $\ \mathcal V:=(\mathcal Q\times\mathcal Q)\setminus\mathcal R$; $\ n_{\mathcal R}:=\sum_p n_{\mathcal R}(p)$; global vacuum $\iff n_{\mathcal R}=0$.

\end{definition}

\begin{definition}[{\texorpdfstring{D0.2}{D0.2}}]\label{tag:D0-2}
$\bigl|\Pi X\bigr|:=\mathfrak T(\Pi X\Pi)=\mathfrak T(\Pi X)$; $\ \ \rho\succeq0\wedge\mathfrak T\rho=1\implies\bigl|\Pi\rho\bigr|\in[0,1]$.

\end{definition}

\begin{definition}[{\texorpdfstring{D0.3}{D0.3}}]\label{tag:D0-3}
$S_{\mathrm{raw}}(\rho):=-\mathfrak T(\rho\log\rho)=\mathfrak T(\rho K)$; $\ K:=-\log\rho$ (faithful $\rho$). $\ \gamma:=1-4\varepsilon_q$ [A4, D0.4, CT0.8]; $\ \gamma\in(0,1]$ for $\varepsilon_q\in[0,\tfrac14)$. Gibbs tilt (functional calculus): $\tilde{\mathcal T}_\gamma(X):=Xe^{-\gamma X}/\mathfrak T(Xe^{-\gamma X})$ where $\mathfrak T(Xe^{-\gamma X})\ne0$; on states: $\tilde{\mathcal T}_\gamma(\rho)=\rho\,e^{-\gamma\rho}/\mathfrak T(\rho\,e^{-\gamma\rho})$; $\ \gamma=0:$ the identity on trace-$1$ states; for faithful $\rho$: $\tilde{\mathcal T}_\gamma(K)=K\rho^{\gamma}/\mathfrak T(K\rho^{\gamma})$ $\ (\gamma=1:\ =\rho(-\log\rho)/S_{\mathrm{raw}}(\rho))$. Variance functional (reference state explicit): $\mathbb V_\rho(X):=\mathfrak T(\rho X^2)-\bigl(\mathfrak T(\rho X)\bigr)^2$. Lyapunov functional of the tilt: $V(\tilde K):=\mathfrak T(\tilde K^2)$, $\ \tilde K:=K-\mathfrak T(K\rho)\,I$ — monotonicity: L1 (§F, forward pointer). $\tilde{\mathcal T}_\gamma\ne\mathcal T$ [A6]; $\neg$ interchangeable.

\end{definition}

\begin{definition}[{\texorpdfstring{D0.4}{D0.4}}]\label{tag:D0-4}
$\varepsilon_q\in[0,1]$: threshold of A4; admissible range: CT0.8; sharp choice: $\varepsilon_q=0$.

\end{definition}

\begin{definition}[{\texorpdfstring{D0.5}{D0.5}}]\label{tag:D0-5}
$\Phi$: phase alphabet of a real edge (CT0.9: $\Phi=\{0,\pi\}$); $\ \sigma_e:=e^{i\phi_e}=(-1)^{\phi_e/\pi}\in\{+1,-1\}$.

\end{definition}

\begin{definition}[{\texorpdfstring{D0.6}{D0.6}}]\label{tag:D0-6}
hop $:=$ one traversal of one $e\in E_{\mathcal R}$; $\ \delta h(\gamma):=\#\{e\in\gamma\}$; $\ \mathfrak h:=$ length assigned to one hop.

\end{definition}

\begin{definition}[{\texorpdfstring{D0.7}{D0.7}}]\label{tag:D0-7}
$\mathfrak L:=$ ordered record of retained events [A6, A7: stratification $\mathfrak L=\bigcup_{n\ge0}\mathfrak L(\tau_n)$, $|\mathfrak L(\tau_0)|=0$, monotone under ticks]; $\ \delta\tau:=$ ledger ticks (retained updates along real edges); $\ U\notin\mathfrak L$ [A5, A6].

\end{definition}

\begin{definition}[{\texorpdfstring{D0.8}{D0.8}}]\label{tag:D0-8}
$n(v):=\partial^{\leftarrow}(v)+\partial^{\rightarrow}(v)$; $\ V_n:=\{v\in V:\ n(v)=n\}$ (vertex set); $\ \Xi_n:=\tau(V_n)=\{(a,b)\in\mathbb N^2:\ a+b=n\}$ (type set); “at rung $n$” $:=$ on $V_n$.

\end{definition}

\begin{definition}[{\texorpdfstring{D0.9}{D0.9}}]\label{tag:D0-9}
$D_{\mathrm{tick}}(n):=\#\{$independent real registers a vertex of $V_n$ pushes across one pair channel per tick$\}$ — counted by the register accounting of \textbf{C10} (one orthogonal direction per admissible classical alternative: source availability, single-shot linear readout; $\neg$ Shannon $\log_2$ counting). $\ C_{\mathrm{cap}}:=\dim_{\mathbb R}\mathfrak H(\mathcal H_p\otimes\mathcal H_q)$ (register directions supplied by $1\mathfrak q$, A2); $\ O(n):=\max\bigl(0,\ D_{\mathrm{tick}}(n)-C_{\mathrm{cap}}\bigr)$.

\end{definition}

\begin{definition}[{\texorpdfstring{D0.10}{D0.10}}]\label{tag:D0-10}
$\mathfrak C(\mathcal R):=$ configuration set on $E_{\mathcal R}$: finite oriented edge sets $E(C)\subseteq E_{\mathcal R}$ with sign assignments $\sigma$, seal data (D1.2) and vertex types (A7). $w(C)$ (A8) is defined on $\mathfrak C(\mathcal R)$; valid universes (D2.11) form a subset; $Z(\mathcal R):=\sum_{C\in\mathfrak C(\mathcal R)}\sum_\sigma w(C)$ (C12).

\end{definition}

\phantomsection
\subsection[Tier 1 — graph, chains, ledger dynamics]{\texorpdfstring{Tier 1 — graph, chains, ledger dynamics}{Tier 1 - graph, chains, ledger dynamics}}

\begin{definition}[{\texorpdfstring{D1.1}{D1.1}}]\label{tag:D1-1}
$e=(p,q)\in\mathcal R_{\rightarrow}(p)\ \longmapsto\ (\mathcal G_n,\phi_e)$; $\ \mathcal G_n:=$ rung class of $q$; $\ \phi_e:$ D0.5; $\ \neg$registration $\implies e\in\mathcal V$. Reality principle of registration (C10): $\phi_e:=$ phase of the registration transport's action on the reality line $\wedge^2$ (1-dimensional, since $\dim\mathcal H_q=2$ [A1]); the transport preserves the $(S,T)$-grading [A3, A4] and the real Pauli-frame structure — registrations are classical data. On a minimal real edge ($\widehat{ds}^{\,2}=1$ [D1.6]) the geometric factor of the transfer weight $G(e)$ [D3.5] is the edge's alignment holonomy $e^{i\alpha_e}$ — the binary registration read of the A4 branch index at the seal ($\alpha_e=\alpha_{pq}$ for $e=(p,q)$ [A4]; per-edge determination read in \textbf{E}: Def-EM.20, EM.20, forward pointer).

\end{definition}

\begin{definition}[{\texorpdfstring{D1.2}{D1.2}}]\label{tag:D1-2}
seal $:=$ transition of one pair across the A4 criterion (entry/exit of either branch): $\ J^+:\ n_{\mathcal R}:0\to1$ (crystallization); $\ J^-:\ n_{\mathcal R}:1\to0$ (dissolution).

\end{definition}

\begin{definition}[{\texorpdfstring{D1.3}{D1.3}}]\label{tag:D1-3}
Janus-balanced $\iff\ \#J^+=\#J^-$.

\end{definition}

\begin{definition}[{\texorpdfstring{D1.4}{D1.4}}]\label{tag:D1-4}
chain $c=(e_1,\dots,e_k)$: $e_i\in E_{\mathcal R}$, $\ t(e_i)=s(e_{i+1})$; $\ s(c):=s(e_1)$; $\ t(c):=t(e_k)$; $\ \partial c:=t(c)-s(c)$ (formal difference of vertices); closed $\iff\ \partial c=0\ \wedge\ \#J^+=\#J^-$.

\end{definition}

\begin{definition}[{\texorpdfstring{D1.5}{D1.5}}]\label{tag:D1-5}
bound $:=$ closed self-contained chain (content: fermionic); transit $:=$ propagating open chain (content: bosonic); $\ \Delta S=\mathsf E_{\mathrm{bd}}+\mathsf E_{\mathrm{tr}}$; $\ \mathsf E_{\mathrm{bd}}:=\Delta S_{\mathrm{bound}}$; $\ \mathsf E_{\mathrm{tr}}:=\Delta S_{\mathrm{transit}}$.

\end{definition}

\begin{definition}[{\texorpdfstring{D1.6}{D1.6}}]\label{tag:D1-6}
$ds^2(e):=g_{\mu\nu}(e)\,\Delta x^\mu_e\,\Delta x^\nu_e$ [$g_{\mu\nu}:$ D2.1]; $\ \widehat{ds}^{\,2}(e):=ds^2(e)/\mathfrak h^2$.

\end{definition}

\begin{definition}[{\texorpdfstring{D1.7}{D1.7}}]\label{tag:D1-7}
$\Theta(\gamma):=\sum_{e\in\gamma}\widehat{ds}^{\,2}(e)$; $\ \Phi(\gamma):=\sum_{e\in\gamma}\alpha_e\,\widehat{ds}^{\,2}(e)$ [$\alpha_e\in\Phi:$ per-edge phase, D3.5; uniform-phase case: $\Phi=\alpha\,\Theta$]; $\ \operatorname{Hol}_{\mathrm{spiral}}(\gamma):=e^{i\Phi(\gamma)}$; $\ \gamma\in\mathsf{Trivial}\iff\forall e\in\gamma:\ \widehat{ds}^{\,2}(e)=0$; wrapping $\iff$ spatial hop projection has winding $\ne0$ around a closed ledger cycle.

\end{definition}

\begin{definition}[{\texorpdfstring{D1.8}{D1.8}}]\label{tag:D1-8}
$\tau_n:=\sum\delta\tau$ accumulated along a causal chain of real edges; emergent coordinate $x^0:=\tau_n$.

\end{definition}

\begin{definition}[{\texorpdfstring{D1.9}{D1.9}}]\label{tag:D1-9}
tick locality: $\ \operatorname{dom}\bigl(\mathrm{tick}(v)\bigr)\subseteq\{v\}\cup\{e\in E:\ e\ni v\}$, pre- and post-tick data only; $\neg\exists$ read at edge-distance $>1$.

\end{definition}

\begin{definition}[{\texorpdfstring{D1.10}{D1.10}}]\label{tag:D1-10}
node force: $\ \mathbf F\!\restriction_v$ resolves local information content into channels \{bound, transit, holonomy, charge, flux\}; \{bound, transit\} $:=$ entropy channels [D1.5]; \{holonomy, charge, flux\} $:=$ registration data [D1.1] $\cup$ transition amplitude [D3.3].

\end{definition}

\begin{definition}[{\texorpdfstring{D1.11}{D1.11}}]\label{tag:D1-11}
accessible information content at $v$ $\leftrightarrow n(v)$; content beyond the capacity bound (CT0.11): $\neg$accessible [D4.1].

\end{definition}

\begin{definition}[{\texorpdfstring{D1.12}{D1.12}}]\label{tag:D1-12}
$\mathcal A(n):=\{$aspects of $\mathbf F$ acting nontrivially at rung $n\}$; value derived (CT0.13), $\neg$ stipulated.

\end{definition}

\phantomsection
\subsection[Tier 2 — masks, geometry, units]{\texorpdfstring{Tier 2 — masks, geometry, units}{Tier 2 - masks, geometry, units}}

\begin{definition}[{\texorpdfstring{D2.1}{D2.1}}]\label{tag:D2-1}
sector coefficients [basis: CT1.2]: $\ X_{(+)}=a_{00}E_0+q_{ij}E^{ij}$, $q_{ij}=q_{ji}$; $\ X_{(-)}=w_iW_i$;

%dsp-level:0
\[
\Sigma_{\mathrm{tr}}:=\delta^{ij}q_{ij};
\]

%dsp-level:0
\[
\quad \hat h_{ij}:=q_{ij}-\tfrac13\delta_{ij}\Sigma_{\mathrm{tr}};
\]

%dsp-level:0
\[
\quad h:=\tfrac13\Sigma_{\mathrm{tr}};
\]

%dsp-level:0
\[
\quad N_i:=\beta w_i;
\]

%dsp-level:0
\[
\quad \Omega:=e^{\alpha_\Omega a_{00}};
\]

%dsp-level:0
\[
\quad \alpha_\Omega,\beta,\varkappa\in\mathbb R;
\]

%dsp-level:0
\[
\bar g_{00}:=-1;\quad
\bar g_{0i}:=\varkappa N_i;\quad
\bar g_{ij}:=\delta_{ij}+\varkappa\bigl(\hat h_{ij}+\tfrac13\delta_{ij}\Sigma_{\mathrm{tr}}\bigr);\quad
g_{\mu\nu}:=\Omega^2\,\bar g_{\mu\nu}.
\]

\end{definition}

\begin{definition}[{\texorpdfstring{D2.2}{D2.2}}]\label{tag:D2-2}
$\Upsilon:=\ln\Omega$; $\ \Upsilon_\mu:=\bar\nabla_\mu\Upsilon$.

\end{definition}

\begin{definition}[{\texorpdfstring{D2.3}{D2.3}}]\label{tag:D2-3}
$\mathcal K:$ aspect space; $\ \mathcal K^{(A)}:$ block of aspect $A$; $\ \mathcal K_n:=\bigoplus_{A\in\mathcal A(n)}\mathcal K^{(A)}$; $\ \mathcal P_M^{(n)}:\mathcal K\to\mathcal K_n:$ orthogonal projection.

\end{definition}

\begin{definition}[{\texorpdfstring{D2.4}{D2.4}}]\label{tag:D2-4}
$Z(\mathcal R,n):=\bigl\langle\mathrm{out}\bigl|(1-\mathbf F_M^{(n)})^{-1}\bigr|\mathrm{in}\bigr\rangle$; poles of $Z(\mathcal R,n):=$ rung-$n$ spectrum; residues $:=$ amplitudes.

\end{definition}

\begin{definition}[{\texorpdfstring{D2.5}{D2.5}}]\label{tag:D2-5}
massless $\iff$ resolvent pole at $k^2=0$; massive $\iff$ pole at $k^2=m^2>0$; $m^2:=$ pole displacement.

\end{definition}

\begin{definition}[{\texorpdfstring{D2.6}{D2.6}}]\label{tag:D2-6}
$\omega:$ orientation-reversal involution: $s(c)\leftrightarrow t(c)$; chirality $\mapsto-$chirality; charge $\mapsto-$charge.

\end{definition}

\begin{definition}[{\texorpdfstring{D2.7}{D2.7}}]\label{tag:D2-7}
$\Theta_{\mathfrak L}:$ ledger reflection: reverses ledger order; $J^+\leftrightarrow J^-$; $G(e)\mapsto G(e)^\ast$.

\end{definition}

\begin{definition}[{\texorpdfstring{D2.8}{D2.8}}]\label{tag:D2-8}
convergence sheet: $1\mapsto1+i\varepsilon$, $\ \varepsilon\to0^+$, for $\forall$ resolvents and configuration sums.

\end{definition}

\begin{definition}[{\texorpdfstring{D2.9}{D2.9}}]\label{tag:D2-9}
units $:=$ powers of $\mathfrak q$; $\ c=\hbar=k_B=1$; $\ [\mathfrak h]=\mathfrak q^{-1/2}$; $\ \mathfrak h:=$ single free scale; $\ \kappa,\ G_N:$ names for combinations of $\mathfrak h$ (CT2.9), $\neg$ independent inputs.

\end{definition}

\begin{definition}[{\texorpdfstring{D2.10}{D2.10}}]\label{tag:D2-10}
$\mathfrak M=\begin{pmatrix}B&Q\\ R&F\end{pmatrix}$: even blocks $B,F$; odd blocks $Q,R$; Hermitian $\iff B^\dagger=B\wedge F^\dagger=F\wedge R=Q^\dagger$; $\ \operatorname{Str}\mathfrak M:=\mathfrak T B-\mathfrak T F$; $\ [X,Y\}:=XY-(-1)^{|X||Y|}YX$ ($X,Y$ homogeneous).

\end{definition}

\begin{definition}[{\texorpdfstring{D2.11}{D2.11}}]\label{tag:D2-11}
valid universe $\iff$ real-edge network closed $\ \wedge\ \#J^+=\#J^-$ [D1.3] $\ \wedge\ \in\operatorname{Im}($active mask$)$ $\ \wedge\ \neg\exists$ dangling ledger boundary beyond the sealed Janus boundary.

\end{definition}

\begin{definition}[{\texorpdfstring{D2.12}{D2.12}}]\label{tag:D2-12}
Aspect realization (slot$\to$aspect bridge): $\iota_n:\mathfrak u(n)\to\mathcal K$ — linear realization of the rung-$n$ slot algebra in aspect space, with $\iota_n\bigl(Z(\mathfrak u(2))\bigr)$ spanning $\mathcal K^{(\mathrm{EM})}$ ($n=2$) and $\iota_n\bigl(\mathfrak{su}(3)\bigr)$ spanning $\mathcal K^{(\mathrm{strong})}$ ($n=3$), per the derivation of CT0.13. Slot algebras and pair algebras remain distinct objects (F header); D3.2 “already realized at rung $m$” $:=$ the abstract isomorphism class lies in $\iota_m(\mathfrak u(m))$ or in the pair blocks of CT0.12. \textbf{Realization domain:} exactly $n\le3$ — for $n\ge4$ no aspect realization exists [CT4.4a]; the silence of this definition at $n\ge4$ is a non-existence consequence, not an underdetermination.

\end{definition}

\begin{definition}[{\texorpdfstring{D2.13}{D2.13}}]\label{tag:D2-13}
Subnetwork and microstate space \textit{(scope clauses)}. $\Sigma\subseteq U$ [D2.11] $:=$ a finite connected vertex set together with its interior edges; $\partial\Sigma:=$ edges with exactly one endpoint in $\Sigma$. $\mathcal M(\Sigma):=$ valid microscopic configurations on $\Sigma$: restrictions $C|_\Sigma$ of valid ambient configurations $C$ [D2.11], carrying edge signs $\sigma_e$ [D0.5, CT0.9], geometric data $\widehat{ds}^{\,2}(e)$ [D1.6], seal data [D1.2] and vertex types [A7, D0.8] — subject to the rung masks [CT0.14, CT2.2] and Janus balance [D1.3], $\neg$ to colour sealing (sealing is ambient-global [CT3.18]; the completeness filter acts downstream: CT3.22a Rmk). $|\Sigma|:=\#\{\text{real edges touching}\ \Sigma\}$ (interior $+$ boundary).

\end{definition}

\begin{definition}[{\texorpdfstring{D2.14}{D2.14}}]\label{tag:D2-14}
Colour boundary and the contraction-completeness projector. $\mathrm{Colour}(U):=\{u\in\mathcal Q:\ u\ \text{carries a colour attachment}\ [\mathrm{D3.8}]\}$; for valid $C$, $\pi_C(v)\subseteq\mathrm{Colour}(U):=$ the other members of the elementary invariant contraction group of $v$ under $C$ (one for a meson pair via $\delta^i{}_j$; two for a (anti)baryon triple via $\varepsilon_{ijk}/\varepsilon^{ijk}$ [CT3.17]); $\pi_C$ is total on $\mathrm{Colour}(U)$ for valid $C$ [CT3.18 $+$ D2.11 applied to the colour ledger]. For $\Sigma$ [D2.13]: $\mathrm{Colour}(\Sigma):=\mathrm{Colour}(U)\cap\Sigma$;

%dsp-level:0
\[
\partial_{\mathrm{col}}(\Sigma;C):=\{v\in\mathrm{Colour}(\Sigma):\ \pi_C(v)\not\subseteq\Sigma\};\qquad
\mathcal C(\Sigma;C):=\begin{cases}1&\partial_{\mathrm{col}}(\Sigma;C)=\varnothing\\ 0&\text{otherwise}\end{cases};\qquad
\Pi_{\mathcal C}(\Sigma;C):=\mathcal C(\Sigma;C)\cdot\mathbb I^{(\mathrm{strong})}_\Sigma
\]

(Hermitian scalar idempotent on the colour slot space of $\Sigma$ [D3.1]; matches the mask pattern $\mathcal P_M^{(n)}$ [CT0.14]).

\end{definition}

\phantomsection
\subsection[Tier 3 — rung-2 and rung-3 structures]{\texorpdfstring{Tier 3 — rung-2 and rung-3 structures}{Tier 3 - rung-2 and rung-3 structures}}

\begin{definition}[{\texorpdfstring{D3.1}{D3.1}}]\label{tag:D3-1}
$\mathcal H^{(n)}_{\mathrm{slot}}:=\operatorname{span}_{\mathbb C}\{|1\rangle,\dots,|n\rangle\}\cong\mathbb C^{n}$; $n=3$: slot index $i\in\{1,2,3\}$ $:=$ colour.

\end{definition}

\begin{definition}[{\texorpdfstring{D3.2}{D3.2}}]\label{tag:D3-2}
subset principle: strictly new at rung $n$ $:=$ part of the rung-$n$ structure $\neg$realized at any $m<n$.

\end{definition}

\begin{definition}[{\texorpdfstring{D3.3}{D3.3}}]\label{tag:D3-3}
at $v$ of type $(1,1)$, edge signs $\sigma_{\leftarrow}(v),\sigma_{\rightarrow}(v)$:

%dsp-level:0
\[
f_v:=\frac{\sigma_{\rightarrow}(v)-\sigma_{\leftarrow}(v)}{2}\in\{-1,0,+1\};\qquad
\aleph_v:=|f_v|;\qquad
\aleph(\sigma,C):=\sum_{v\in C}\aleph_v .
\]

\end{definition}

\begin{definition}[{\texorpdfstring{D3.4}{D3.4}}]\label{tag:D3-4}
$q_{\mathrm{EM}}(c):=\dfrac{\sigma_{s(c)}-\sigma_{t(c)}}{2}$ (oriented chain $c$). Closed chains: the basepoint core read $q_{\mathrm{core}}$ is defined, ternary, and charge-conjugation odd [EM.9a].

\end{definition}

\begin{definition}[{\texorpdfstring{D3.5}{D3.5}}]\label{tag:D3-5}
$G^{(n)}(e):=\sigma_e\,\exp\!\bigl(i\,\alpha_e\,\widehat{ds}^{\,2}(e)\bigr)$ [D0.5, D1.6]; $\ \alpha_e\in\Phi:$ per-edge dynamic phase $=\alpha_{pq}$ for $e=(p,q)$ — the A4 branch index (read in Part 2 §E: Def-EM.20, EM.20, forward pointer); dimensionless; the same alphabet at every rung. $\ \alpha_{\mathrm{EM}}^{(n)}:$ EM-block spectral read at rung $n$ (EM.17, forward pointer).

\end{definition}

\begin{definition}[{\texorpdfstring{D3.6}{D3.6}}]\label{tag:D3-6}
$\mathcal J(C):=\prod_{J^+\in C}(+1)\cdot\prod_{J^-\in C}(+1)^\ast$ $\ \implies\ \mathcal J(\Theta_{\mathfrak L}C)=\mathcal J(C)^\ast$.

\end{definition}

\begin{definition}[{\texorpdfstring{D3.7}{D3.7}}]\label{tag:D3-7}
photon $:=$ trivial chain [D1.7] $\subseteq V_2$; $\ $gluon $:=$ trivial chain $\cap\,V_3\ne\varnothing$.

\end{definition}

\begin{definition}[{\texorpdfstring{D3.8}{D3.8}}]\label{tag:D3-8}
colour triplet $q^i$: one fundamental slot index; colour antitriplet $r_i$: one dual slot index; quark $:=$ bound (fermionic) loop with a triplet attachment.

\end{definition}

\begin{definition}[{\texorpdfstring{D3.9}{D3.9}}]\label{tag:D3-9}
baryon $B:=\varepsilon_{ijk}\,q^iq^jq^k$; $\ $meson $M:=\delta^i{}_j\,q^jr_i$. [Disambiguation: $M$ here $\ne$ total mask $M_\oplus$ of A9.]

\end{definition}

\begin{definition}[{\texorpdfstring{D3.10}{D3.10}}]\label{tag:D3-10}
$W^\pm:=\tfrac1{\sqrt2}\bigl(W_1\mp iW_2\bigr)$; $\ Z:=\cos\theta_W\,W_3-\sin\theta_W\,Y$; $\ A:=\sin\theta_W\,W_3+\cos\theta_W\,Y$ [$\theta_W:$ EW.24].

\end{definition}

\phantomsection
\subsection[Tier 4 — saturation]{\texorpdfstring{Tier 4 — saturation}{Tier 4 - saturation}}

\begin{definition}[{\texorpdfstring{D4.1}{D4.1}}]\label{tag:D4-1}
$F_{D^-}=\bigoplus_{n\ge4}F_n$ [A8]: single possibility space receiving $\forall$ configurations with rung $n\ge4$.

\end{definition}

\begin{definition}[{\texorpdfstring{D4.2}{D4.2}}]\label{tag:D4-2}
$D^-$ absorbs information at rate $\propto O(n)$ [D0.9]; absorbed $\implies$ locked $:=$ retained in $\mathfrak L$, $\neg$accessible [D1.11].

\end{definition}

\begin{definition}[{\texorpdfstring{D4.3}{D4.3}}]\label{tag:D4-3}
$F_{D^+}=\Lambda\,\mathbb I$ [A8]: uniform scalar aspect; $\neg$configuration-dependent.

\end{definition}

\phantomsection
\section[§D — Master symbol table]{\texorpdfstring{§D — Master symbol table}{SD - Master symbol table}}

%tbl-level:2
{\footnotesize\setlength{\tabcolsep}{4pt}
\begin{longtable}{>{\raggedright\arraybackslash}p{\dimexpr0.2460\textwidth-2\tabcolsep\relax}|>{\raggedright\arraybackslash}p{\dimexpr0.6032\textwidth-2\tabcolsep\relax}|>{\raggedright\arraybackslash}p{\dimexpr0.1508\textwidth-2\tabcolsep\relax}}
\toprule
Symbol & Meaning & Declared \\
\addlinespace[2pt]
\midrule
\endfirsthead
\toprule
Symbol & Meaning & Declared \\
\addlinespace[2pt]
\midrule
\endhead
\bottomrule
\endlastfoot
$\mathcal Q$, $q_n$ & qubit set; $n$-th qubit & A1 \\
$\mathcal H_q$ & carrier space of $q$, $\cong\mathbb C^2$ & A1 \\
$\rho_{pq}$ & state of ordered pair $(p,\allowbreak q)$ & A1 \\
$\mathfrak H(X)$ & real space of Hermitian operators on $X$ & A1 \\
$\mathfrak T(A)$ & trace of $A$ & A1 \\
$\succeq0$ & positive semidefinite & A1 \\
$C(p,\allowbreak q)$, $\mathfrak q$ & channel capacity of a pair; channel quantum & A2 \\
$\vee^2$, $\wedge^2$ & symmetric, antisymmetric square & A3 \\
$S$, $T$ & swap involution; sign-grading involution ($T=\allowbreak \mathsf D_\zeta\otimes\mathsf D_\zeta$; metric-sign reading: C13) & A3, CT1.0 \\
$\sigma_\mu$ & Pauli basis, $\mu\in\allowbreak \{0,1,2,3\}$, $\sigma_0=\allowbreak I_2$ & A3, C1 \\
$\mathbb1$, $I$ & identity operator; identity matrix & A3, A5 \\
$\eta$ & sign vector of the $T$-grading, $\eta\in\allowbreak \{\pm1\}^4$, one odd entry; representative $(-,\allowbreak +\allowbreak ,\allowbreak +\allowbreak ,\allowbreak +\allowbreak )$ & A3, C13 \\
$\mathfrak o$, $\eta^{(\mathfrak o)}$ & odd-sign slot of $\eta$; sign vector with a single $-1$ at slot $\mathfrak o$ & A3, CT1.0 \\
$\zeta_i$, $\mathsf D_\zeta$, $\mathrm{sf}$ & $\eta^0\eta^i$; Bloch sign map $\sigma_\mu\mapsto\allowbreak \zeta_\mu\sigma_\mu$; spin flip $X\mapsto\allowbreak \sigma_2X^{\mathrm T}\sigma_2$ & CT1.0 \\
$\mathcal R$, $\mathcal V$ & real (crystallized) set; virtual set & A4, D0.1 \\
$\Pi_{\mathcal V}$, $\Pi_{\mathcal R}$ & projections onto $\wedge^2$, $\vee^2$ & A4, CT0.7 \\
$\varepsilon_q$ & crystallization threshold & A4, D0.4 \\
$n_{\mathcal R}(p)$, $n_{\mathcal R}$ & real-edge count at $p$; global count & A4, D0.1 \\
$\rho_\circ$ & vacuum pair state $I/4$ & A5 \\
$U$, $\mu_\star$, $\dagger$ & Haar draw on $\mathrm{SU}(2)$; Haar measure; conjugate transpose & A5, A6 \\
$\mathfrak L$, $\mathfrak L(\tau_n)$, $n_{\mathrm{real\_out}}$ & ledger; ledger stratum at tick $\tau_n$; real-out count per tick & A5, A6, A7, D0.7 \\
$\mathcal T$, $\mathcal T^{-1}$ & tick map (Haar conjugation); its inverse & A6 \\
$\tilde{\mathcal T}_\gamma$, $\gamma$, $K$, $S_{\mathrm{raw}}$, $\mathbb V_\rho$ & $\varepsilon_q$-dependent Gibbs tilt; tilt strength $1-4\varepsilon_q$; surprisal; raw entropy; variance & D0.3 \\
$\tilde K$, $V(\tilde K)$ & centered surprisal; Lyapunov functional & D0.3, L1 \\
$\Gamma=\allowbreak (V,\allowbreak E)$ & directed graph & A7 \\
$\tau(v)$, $\partial^{\leftarrow}$, $\partial^{\rightarrow}$ & vertex type; in-degree; out-degree & A7 \\
$n(v)$, $V_n$, $\Xi_n$ & rung of $v$; rung-$n$ vertex set; type set & D0.8 \\
$E_{\mathcal R}$, $\phi_e$, $\sigma_e$, $\Phi$ & real edges; edge phase; edge sign; phase alphabet & A7, D0.5 \\
$\delta h$, $\mathfrak h$ & hop count; hop length & D0.6 \\
$\delta\tau$, $\tau_n$ & ledger tick count; ledger time & D0.7, D1.8 \\
$D_{\mathrm{tick}}(n)$, $C_{\mathrm{cap}}$, $O(n)$ & tick load; channel capacity; oversubscription & D0.9 \\
$\mathbf F$, $F_A$ & force operator; aspect $A$ & A8 \\
$F_{W^+}$, $F_{W^-}$, $F_Z$, $F_W$ & charged weak; neutral weak ($F_Z\equiv F_{W^3}$); weak block & A8, C8 \\
$F_{J^+}$, $F_{J^-}$ & crystallization; dissolution aspects & A8, D1.2 \\
$F_{\mathrm{EM}}$, $F_{\mathrm{strong}}$ & electromagnetic; colour aspects & A8 \\
$\Lambda$, $\mathbb I$ & stretch scalar; identity on force space & A8 \\
$\mathcal U_\odot$, $\mathfrak S$, $\mathbf F_\oplus$, $M_\oplus$ & observed universe; solution set; total force; total mask & A9 \\
$\mathfrak C(\mathcal R)$ & configuration set on $E_{\mathcal R}$; summation domain of $Z$ & D0.10, C12 \\
$\tilde{\mathcal T}_\gamma(X)$ & operator Gibbs tilt (functional calculus) & D0.3 \\
$\mathcal E_{st}$ & joint $(S,\allowbreak T)$ eigenspace & C3 \\
$f_{AB}{}^{C}$ & structure constants (Hermitian convention) & C2 \\
$\mathsf E_{\mathrm{bd}}$, $\mathsf E_{\mathrm{tr}}$ & bound, transit entropy & C6, D1.5 \\
$E_0$, $E^{ij}$, $W_i$, $J_i$, $K_i$ & generator basis of $\mathfrak H(\mathbb C^4)$ & CT1.2 \\
$\mathcal A(n)$ & active aspect set at rung $n$ & D1.12, CT0.13 \\
$\mathcal G_n$ & target rung class of an edge & D1.1 \\
$J^+$, $J^-$ & crystallization seal; dissolution seal & D1.2 \\
$c$, $\gamma$, $s(c)$, $t(c)$, $\partial c$ & chain; source; target; boundary & D1.4 \\
$ds^2(e)$, $\widehat{ds}^{\,2}(e)$ & discrete line element; hop-normalized & D1.6 \\
$\Theta(\gamma)$, $\Phi(\gamma)$, $\operatorname{Hol}_{\mathrm{spiral}}$, $\mathsf{Trivial}$ & spiral sum; phase; holonomy; trivial class & D1.7 \\
$x^0$ & emergent time coordinate & D1.8 \\
$\ne\allowbreak g$, $\wedge$, $\vee$, $\implies$, $\iff$ & negation; conjunction; disjunction; implication; equivalence & — \\
$\oplus$, $\bigoplus$, $\otimes$ & direct sum; indexed direct sum; tensor product & A1, A3, A8 \\
$F_G$, $F_H$, $F_{D^+}$, $F_{D^-}$ & spin-2; scalar (Higgs); uniform dark; saturated dark & A8, D4.1, D4.3 \\
$P_M$, $\mathcal P_M^{(n)}$, $\mathbf F_M^{(n)}$ & mask projection; rung mask; masked force & A8, C7, D2.3 \\
$w(C)$, $\mathcal J(C)$, $G(e)$ & configuration weight; Janus factor; edge transfer weight & A8, D3.5, D3.6 \\
$f_v$, $\aleph_v$, $\aleph(\sigma,\allowbreak C)$ & vertex transition amplitude; indicator; count & D3.3 \\
$Z(\mathcal R)$, $Z(\mathcal R,\allowbreak n)$ & partition function; masked resolvent & A8, D2.4 \\
$\iota_n$ & aspect realization $\mathfrak u(n)\to\mathcal K$ (slot$\to$aspect bridge) & D2.12 \\
$\mathcal P^{(\mathrm{EM})}$, $Y_{\mathrm{asp}}$, $Y_{\mathrm{slot}}$, $Y_{\mathrm{car}}$ & EM block projection; the three realizations of $Y$ & EM.2, EW.8, D2.12 \\
$\varkappa$, $\kappa$, $G_N$ & embedding scale; gravitational coupling; Newton constant & C4, D2.1, D2.9, CT2.9 \\
$\Sigma$, $\Sigma_{\mathrm{tr}}$ & operator $\sum_i\sigma_i\otimes\sigma_i$; scalar $\delta^{ij}q_{ij}$ & C5, D2.1 \\
$a_{00}$, $q_{ij}$, $w_i$ & sector coefficients & D2.1 \\
$\hat h_{ij}$, $h$, $N_i$, $\Omega$, $\Upsilon$ & traceless shear; trace mode; shift; conformal factor; potential & D2.1, D2.2 \\
$\bar g_{\mu\nu}$, $g_{\mu\nu}$ & auxiliary metric; emergent metric & D2.1 \\
$\alpha_\Omega$, $\beta$ & metric embedding constants & D2.1 \\
$\mathcal K$, $\mathcal K^{(A)}$, $\mathcal K_n$ & aspect space; aspect block; rung-$n$ block & D2.3 \\
$m^2$ & mass squared (pole displacement) & D2.5 \\
$\omega$, $\Theta_{\mathfrak L}$ & orientation reversal; ledger reflection & D2.6, D2.7 \\
$\mathfrak M$, $\operatorname{Str}$, $[\cdot,\allowbreak \cdot\}$ & supermatrix; supertrace; graded bracket & D2.10 \\
$\mathcal H^{(n)}_{\mathrm{slot}}$ & slot space $\cong\mathbb C^n$ & D3.1 \\
$q_{\mathrm{EM}}(c)$ & chain end-charge & D3.4 \\
$\alpha_e$, $\alpha_{pq}$ & per-edge dynamic phase $\in\allowbreak \Phi$; branch phase of a crystallizing pair & A4, A7, D3.5 \\
$\alpha^{(n)}_{\mathrm{EM}}$ & EM-block spectral read at rung $n$ & EM.17 \\
$\mathcal P(\Sigma)$, $M(\mathcal P)$ & alignment pattern; mirror image of the pattern & Def-EM.21, Def-EM.22 \\
$q^i$, $r_i$ & colour triplet; antitriplet & D3.8 \\
$B$, $M$ & baryon; meson & D3.9 \\
$W^\pm$, $Z$, $A$, $Y$, $\theta_W$ & weak eigenstates; photon direction; hypercharge (realizations: EW.8); mixing angle & D3.10, EW.8, EW.24 \\
$T_a$, $f_{abc}$ & colour generators; colour structure constants & CT3.9, CT3.21 \\
$\varepsilon_{ijk}$, $\delta^i{}_j$ & colour invariant tensors & CT3.17 \\
$g$, $g'$, $e$, $h_0$ & weak, hypercharge, EM couplings; Higgs background & EW.33, P4 \\
$F_n$ & rung-$n$ saturation block, $n\ge\allowbreak 4$ & A8, D4.1 \\
$\Sigma$, $\mathcal M(\Sigma)$, $C|_\Sigma$ & subnetwork, its microstate space, configuration restriction & D2.13 \\
$\partial_{\mathrm{col}}$, $\mathcal C(\Sigma;C)$, $\Pi_{\mathcal C}$ & colour boundary, contraction completeness, its projector & D2.14 \\
$\kappa_e$, $\Theta(c)$, $\bar\kappa$ & per-edge line-element quantum, loop total, mean quantum & CT1.29a \\
\end{longtable}
}

\phantomsection
\section[§E — Calibration policy]{\texorpdfstring{§E — Calibration policy}{SE - Calibration policy}}

\phantomsection\label{tag:P1}
\textbf{P1.} $\forall$ quantities of the corpus: derived from A1–A9 $\wedge$ C1–C13 $\wedge$ D0.1–D4.3, except the hop length $\mathfrak h$; $[\mathfrak h]=\mathfrak q^{-1/2}$; $\mathfrak h$ fixes the absolute scale.

\phantomsection\label{tag:P2}
\textbf{P2.} Derived, zero freedom:

%dsp-level:0
\[
(\dim\mathcal E_{++},\dim\mathcal E_{+-},\dim\mathcal E_{-+},\dim\mathcal E_{--})=(7,3,3,3);\quad
C_{\mathrm{cap}}=16;\quad
D_{\mathrm{tick}}(n)=2^{n+1};\quad
n_{\max}=3;
\]

%dsp-level:0
\[
N_{\text{colour}}=3;\quad
N_{\text{gluon}}=8;\quad
N_{\gamma\text{-pol}}=2;\quad
q_{\mathrm{EM}}\in\{-1,0,+1\};\quad
\alpha_e\in\{0,\pi\}\ \text{per edge};
\]

%dsp-level:0
\[
m_W=m_Z\cos\theta_W;\quad
e=g\sin\theta_W;\quad
S=n_{\mathcal R}\ln2;\quad
w_{D^+}=-1 .
\]

\textbf{Audit note.} Every entry of this list is re-verified from first principles in exact arithmetic: \texttt{checks/\allowbreak{}verify\_\allowbreak{}merge\_\allowbreak{}tierB.py} §F (totals and cross-references; the identity-level checks live in \texttt{verify\_\allowbreak{}B.py}–\texttt{verify\_\allowbreak{}G.py}).

\phantomsection\label{tag:P3}
\textbf{P3.} Named combinations of $\mathfrak h$ (external comparison only; $\neg$ inputs):

%dsp-level:0
\[
G_N=\frac{\mathfrak h^2}{4\ln2};\qquad
\kappa^2=8\pi G_N=\frac{2\pi\,\mathfrak h^2}{\ln2}\qquad[\mathrm{CT2.9}].
\]

\phantomsection\label{tag:P4}
\textbf{P4.} Free internal configuration parameters ($\neg$ constants of the theory). Independent core: $g,\ g';\ h_0;\ c_W;\ \Lambda;\ \varkappa,\ \beta,\ \alpha_\Omega$. Functionally dependent members (properties of a solution, $\neg$ further freedom): $\theta_W(g,g')$ [EW.24]; $e(g,g')$ [EW.35]; $\mu_\star^2(c_W,\varkappa,\beta,h_0)$ [EW.33]; $m_W,\ m_Z$ [EW.20, EW.23]; $m_H$ [EW.27]. Predictions $:=$ dimensionless ratios among derived quantities at fixed core data.

\phantomsection\label{tag:P5}
\textbf{P5.} $\neg\exists$ quantity identified with a measured external constant. Comparison procedure: fix $\mathfrak h$ once; read P2 $\wedge$ P3; result $:=$ test of A9, $\neg$ input to A1–A8. \textbf{Era-extension note.} The P4 core lists the bosonic-era free internal parameters (Part 2 loci). The fermion-era parameters (the Janus rate ratio $r=x_+/x_-$; the defect residues $\hat\mu^2$, $\hat\mu'^2$; the Ansatz-B formal temperature $\alpha$; the channel scales $\mu$, $\mu_\nu$; the formal fugacities) are registered in Part 3 §Cal-F under this policy; the baryomesonic era introduces none (Part 4 §Cal-BM).

\phantomsection
\section[§F — Foundations: Bookkeeping, Pair Algebra, Ledger Dynamics]{\texorpdfstring{§F — Foundations: Bookkeeping, Pair Algebra, Ledger Dynamics}{SF - Foundations: Bookkeeping, Pair Algebra, Ledger Dynamics}}

\textit{Dep.: §A–§E of this part. Notation: as §A–§C. Machine: \texttt{checks/\allowbreak{}verify\_\allowbreak{}B.py}, \texttt{checks/\allowbreak{}verify\_\allowbreak{}L\_\allowbreak{}lyapunov.py}, \texttt{checks/\allowbreak{}verify\_\allowbreak{}L0\_\allowbreak{}ledger\_\allowbreak{}tilt.py}, \texttt{checks/\allowbreak{}verify\_\allowbreak{}merge\_\allowbreak{}tierB.py}, \texttt{checks/\allowbreak{}verify\_\allowbreak{}promotion\_\allowbreak{}B\_\allowbreak{}rung.py}.}

\phantomsection
\section[Tier 0 — Bookkeeping, Reality, and Capacity]{\texorpdfstring{Tier 0 — Bookkeeping, Reality, and Capacity}{Tier 0 - Bookkeeping, Reality, and Capacity}}

\begin{theorem}[{\texorpdfstring{CT0.1 — Real/virtual dichotomy}{CT0.1 - Real/virtual dichotomy}}]\label{tag:CT0-1}
\textbf{Thm.} $\ \forall(p,q)\in\mathcal Q\times\mathcal Q:\quad\bigl[(p,q)\in\mathcal R\bigr]\ \underline\vee\ \bigl[(p,q)\in\mathcal V\bigr].$

\textbf{Req.} A1, A3, A4; D0.1, D0.2, D0.4.

\textbf{Pf.}

%dsp-level:0
\[
\bigl|\Pi_{\mathcal V}\rho_{pq}\bigr|=\mathfrak T\bigl(\Pi_{\mathcal V}\rho_{pq}\bigr)\in[0,1]\subset\mathbb R
\]

%dsp-level:0
\[
\quad [\mathrm{D0.2},\ \mathrm{A1}:\ \rho_{pq}\succeq0\wedge\mathfrak T\rho_{pq}=1;\ \mathrm{A3}:\ \Pi_{\mathcal V}=\Pi_{\wedge^2}]
\]

%dsp-level:0
\[
x\in\mathbb R\implies\bigl[(x\le\varepsilon_q\ \vee\ x\ge1-\varepsilon_q)\ \underline\vee\ (\varepsilon_q<x<1-\varepsilon_q)\bigr]\quad[\text{total order of }\mathbb R];
\]

%dsp-level:0
\[
\qquad \varepsilon_q\in[0,1]\subset\mathbb R\ [\mathrm{D0.4}]
\]

%dsp-level:0
\[
\bigl[(p,q)\in\mathcal R\bigr]\iff\bigl[\bigl|\Pi_{\mathcal V}\rho_{pq}\bigr|\le\varepsilon_q\ \vee\ \bigl|\Pi_{\mathcal V}\rho_{pq}\bigr|\ge1-\varepsilon_q\bigr]\ [\mathrm{A4}];
\]

%dsp-level:0
\[
\qquad \mathcal V=(\mathcal Q\times\mathcal Q)\setminus\mathcal R\ [\mathrm{D0.1}]
\]

%dsp-level:0
\[
\Longrightarrow\quad\bigl[(p,q)\in\mathcal R\bigr]\ \underline\vee\ \bigl[(p,q)\in\mathcal V\bigr].\ \blacksquare
\]

\textbf{Phys.} Every qubit pair has exactly one status — crystallized (real) or uncrystallized (virtual) — decided by A4's two-branch threshold test on its antisymmetric weight. This dichotomy is the bookkeeping ground of the corpus: real pairs become edges of the emergent graph, virtual pairs remain vacuum.

\end{theorem}

\begin{theorem}[{\texorpdfstring{CT0.2 — Global vacuum}{CT0.2 - Global vacuum}}]\label{tag:CT0-2}
\textbf{Thm.} $\ n_{\mathcal R}=0\ \implies\ \forall(p,q):\ \bigl[(p,q)\in\mathcal V\ \wedge\ \rho_{pq}=\rho_\circ=\tfrac14 I\bigr]\ \wedge\ S_{\mathrm{raw}}(\rho_\circ)=\ln4=\max_{\rho}S_{\mathrm{raw}}(\rho),$ with $\arg\max=\{\tfrac14 I\}$.

\textbf{Req.} A1, A4, A5; D0.1, D0.3; CT0.1; Gibbs lemma: $\ S_{\mathrm{raw}}(\rho)\le\ln d$, $\ S_{\mathrm{raw}}(\rho)=\ln d\iff\rho=I/d$ [linear algebra].

\textbf{Pf.}

%dsp-level:0
\[
n_{\mathcal R}=0\implies\neg\,\exists(p,q)\in\mathcal R\ [\mathrm{D0.1}]\implies\forall(p,q):\ (p,q)\in\mathcal V\ [\mathrm{CT0.1}]
\]

%dsp-level:0
\[
\rho_{pq}=\rho_\circ=\tfrac14 I\quad[\mathrm{A5}]
\]

%dsp-level:0
\[
S_{\mathrm{raw}}(\rho_\circ)=-\mathfrak T\bigl(\tfrac14 I\ln\tfrac14 I\bigr)=-4\cdot\tfrac14\ln\tfrac14=\ln4\quad[\mathrm{D0.3}]
\]

%dsp-level:0
\[
\mathfrak H(\mathcal H_p\otimes\mathcal H_q)\cong\mathfrak H(\mathbb C^4),\ d=4:\quad S_{\mathrm{raw}}(\rho)\le\ln4,\quad S_{\mathrm{raw}}(\rho)=\ln4\iff\rho=\tfrac14 I\quad[\mathrm{A1},\ \text{Gibbs}].\ \blacksquare
\]

\textbf{Phys.} The empty configuration (no real edges) is all-virtual, with every pair at the maximally mixed state $I/4$ and entropy maximal at $\ln 4$: the vacuum of this theory is a state of maximal raw information. All structure — edges, particles, geometry — is a lowering of this entropy by crystallization.

\end{theorem}

\begin{theorem}[{\texorpdfstring{CT0.3 — Seal increments}{CT0.3 - Seal increments}}]\label{tag:CT0-3}
\textbf{Thm.} $\ J^{+}:\ (\Delta n_{\mathcal R},\ \delta h,\ \delta\tau)=(+1,+1,+1);\qquad J^{-}:\ (\Delta n_{\mathcal R},\ \delta h,\ \delta\tau)=(-1,-1,-1).$

\textbf{Req.} A4, A7; D0.1, D0.6, D0.7, D1.2, D1.9.

\textbf{Pf.}

%dsp-level:0
\[
J^\pm\ [\mathrm{D1.2}]:\ \ \exists!\,(p,q):\ \Delta\bigl[(p,q)\in\mathcal R\bigr]=\pm1\ [\mathrm{A4}]\implies\Delta n_{\mathcal R}=\pm1\ [\mathrm{D0.1}]
\]

%dsp-level:0
\[
\Delta\bigl|E_{\mathcal R}\bigr|=\Delta n_{\mathcal R}\ [\mathrm{D0.1}]\implies\delta h=\pm1\ [\mathrm{D0.6}:\ \delta h=\#\{e\in E_{\mathcal R}\ \text{traversed}\}]
\]

%dsp-level:0
\[
\text{seal}=\text{retained update along a real edge}\ [\mathrm{D0.7}]\implies\delta\tau=\pm1\ \text{(same sign)}
\]

%dsp-level:0
\[
\mathrm{D1.9}:\ \operatorname{dom}(\mathrm{tick})\subseteq\{v\}\cup\{e\ni v\}\implies\Delta(\text{all further counters})=0.\ \blacksquare
\]

\textbf{Phys.} Crystallization ($J^+$) and dissolution ($J^-$) are the atomic ledger events: each adds or removes exactly one real edge, one traversal unit, and one tick of ledger time. Every extensive quantity of the corpus moves in integer steps of these events.

\end{theorem}

\begin{theorem}[{\texorpdfstring{CT0.4 — Real outgoing edge registration}{CT0.4 - Real outgoing edge registration}}]\label{tag:CT0-4}
\textbf{Thm.} $\ e=(p,q):\quad e\in\mathcal R_{\rightarrow}(p)\iff e\longmapsto(\mathcal G_n,\phi_e).$

\textbf{Req.} A7; D0.5, D0.8, D1.1; CT0.1.

\textbf{Pf.}

%dsp-level:0
\[
(\Rightarrow):\ \ e\in E_{\mathcal R}\implies\exists\,\phi_e\ [\mathrm{A7}];\qquad q\in V_n\ \text{for}\ n=n(q)\ [\mathrm{D0.8}]\implies\exists\,\mathcal G_n;
\]

%dsp-level:0
\[
\mathrm{D1.1}:\ \ e\longmapsto(\mathcal G_n,\phi_e).\qquad
(\Leftarrow):\ \ \mathrm{D1.1}:\ \neg\bigl[e\longmapsto(\mathcal G_n,\phi_e)\bigr]\implies e\in\mathcal V;
\]

%dsp-level:0
\[
e\longmapsto(\mathcal G_n,\phi_e)\implies e\notin\mathcal V\implies e\in\mathcal R\quad[\mathrm{CT0.1}].\ \blacksquare
\]

\textbf{Phys.} “Real” and “registered in the ledger” are the same property of an edge: real edges are exactly those carrying ledger data (target rung class and phase); unregistered edges are virtual. Existence, in this corpus, is recordability.

\end{theorem}

\begin{theorem}[{\texorpdfstring{CT0.5 — Tick load at rung $n$}{CT0.5 - Tick load at rung}}]\label{tag:CT0-5}
\textbf{Thm.} $\ D_{\mathrm{tick}}(n)=2^{\,n+1}.$

\textbf{Req.} A7; C10; D0.7, D0.8, D0.9, D1.9; CT0.9.

\textbf{Pf.}

%dsp-level:0
\[
v\in V_n\implies\bigl|\{e\in E_{\mathcal R}:\ e\ni v\}\bigr|=n\quad[\mathrm{D0.8}]
\]

%dsp-level:0
\[
\phi_e\in\Phi,\ |\Phi|=2\ [\mathrm{CT0.9}]\implies(\phi_e)_{e\ni v}\in\Phi^{\,n},\quad\bigl|\Phi^{\,n}\bigr|=2^n
\]

%dsp-level:0
\[
\mathrm{D1.9}:\ \ \mathrm{tick\ datum}=(\mathrm{pre},\mathrm{post})\in\Phi^{\,n}\times\Phi^{\,n},\quad \mathrm{pre}\perp\mathrm{post}
\]

%dsp-level:0
\[
\mathrm{C10}/\mathrm{D0.9}\ (\text{register accounting}):\ \ \text{registration of}\ X\ \text{at single-shot linear readout}
\]

%dsp-level:0
\[
= |X|\ \text{orthogonal register directions}
\]

%dsp-level:0
\[
\text{Shannon content}=\log_2|X|=n\ \text{bits}\ \ne\ \text{register cost}\ 2^n\quad(\text{source availability}\ \ne\ \text{throughput})
\]

%dsp-level:0
\[
\Longrightarrow\quad D_{\mathrm{tick}}(n)=2^n+2^n=2^{\,n+1}.\ \blacksquare
\]

\textbf{Phys.} A vertex with $n$ real edges must make $2^{n+1}$ classical alternatives linearly distinguishable per tick — $2^n$ pre-tick and $2^n$ post-tick phase configurations (C10 register accounting: one orthogonal direction per admissible value, not $\log_2$ bit encoding). This exponential growth of source availability is what caps accessible complexity at rung 3 (CT0.11).

\end{theorem}

\begin{theorem}[{\texorpdfstring{CT0.6 — Channel capacity and oversubscription}{CT0.6 - Channel capacity and oversubscription}}]\label{tag:CT0-6}
\textbf{Thm.}

%dsp-level:0
\[
C_{\mathrm{cap}}=\dim_{\mathbb R}\mathfrak H(\mathbb C^4)=16;\qquad
O(n)=\max\bigl(0,\ 2^{\,n+1}-16\bigr);
\]

%dsp-level:0
\[
O(0)=O(1)=O(2)=O(3)=0;\qquad O(4)=16;\qquad O(n)=2^{\,n+1}-16\quad(n\ge4).
\]

\textbf{Req.} A1, A2; D0.9; CT0.5; $\dim_{\mathbb R}\mathfrak H(\mathbb C^d)=d^2$ [linear algebra].

\textbf{Pf.}

%dsp-level:0
\[
\mathcal H_p\otimes\mathcal H_q\cong\mathbb C^2\otimes\mathbb C^2\cong\mathbb C^4\ [\mathrm{A1}];\qquad
\dim_{\mathbb R}\mathfrak H(\mathbb C^d)=d^2\ [\text{linear algebra}]
\]

%dsp-level:0
\[
C_{\mathrm{cap}}=\dim_{\mathbb R}\mathfrak H(\mathbb C^4)=4^2=16\ [\mathrm{D0.9}];\qquad C(p,q)=1\,\mathfrak q\ \ \forall(p,q)\ [\mathrm{A2}]
\]

%dsp-level:0
\[
O(n)=\max\bigl(0,\ D_{\mathrm{tick}}(n)-C_{\mathrm{cap}}\bigr)=\max\bigl(0,\ 2^{\,n+1}-16\bigr)\quad[\mathrm{D0.9},\ \mathrm{CT0.5}]
\]

%dsp-level:0
\[
2^{\,n+1}\le16\iff n\le3\implies O(n)=0\ \ (n\le3);
\]

%dsp-level:0
\[
\qquad 2^{5}-16=16\implies O(4)=16;
\]

%dsp-level:0
\[
\qquad n\ge4\implies2^{\,n+1}\ge32>16\implies O(n)=2^{\,n+1}-16.\ \blacksquare
\]

\textbf{Phys.} One qubit–qubit channel carries exactly 16 real registers per tick — the dimension of the two-qubit Hermitian operator space. Rungs $0$–$3$ fit the budget, rung $3$ exactly saturates it, and rung $4$ doubles the demand to $32$, exceeding it by exactly one channel quantum. All finiteness of the accessible world starts here.

\end{theorem}

\begin{corollary}[{\texorpdfstring{CT0.6a — transit-hop register bound}{CT0.6a - transit-hop register bound}}]\label{tag:CT0-6a}
Let $\gamma=(e_1,\dots,e_{\delta h})$ be a transit chain [D1.5, D1.4] of distinct directed real edges with hop count $\delta h(\gamma)$ [D0.6] whose entire edge-register content is registered through one pair channel over $\delta\tau(\gamma)$ ledger ticks [A2, D0.7]. Then:

%dsp-level:0
\[
\mathsf E(\gamma)=\delta h(\gamma)\,\ln2\ \ \text{nats}\qquad(\delta h\ \ \text{independent binary registers}\ [\mathrm{CT0.9},\ \mathrm{C9}]);
\]

%dsp-level:0
\[
\text{single-channel supply over}\ \delta\tau\ \text{ticks}:\ \log_2C_{\mathrm{cap}}\cdot\delta\tau=4\,\delta\tau\ \ \text{bits}
\]

%dsp-level:0
\[
\qquad [\mathrm{CT0.6}:\ C_{\mathrm{cap}}=16;\ \mathrm{C10}:\ 16\ \text{directions}\iff2^{4}\ \text{alternatives}]
\]

%dsp-level:0
\[
\implies\quad\delta h(\gamma)\ \le\ 4\,\delta\tau(\gamma);\qquad
\mathsf E(\gamma)\ \le\ \min\bigl(\delta h(\gamma),\ 4\,\delta\tau(\gamma)\bigr)\,\ln2\ \ \text{nats}.
\]

A per-tick payload beyond $4$ hop-bits presents $2^{5}=32>16=C_{\mathrm{cap}}$ alternatives at the crossing — oversubscribed [CT0.6]; this is the payload counterpart of CT0.11's vertex-rung admissibility. The bound is sharp: for $\delta h\le4\,\delta\tau$ a schedule of $\delta\tau$ consecutive groups of $\le4$ hops presents $\le2^{4}=16=C_{\mathrm{cap}}$ alternatives per tick — at capacity, no oversubscription [CT0.6]. (Ledger accounting of registered content vs channel supply over a tick window [D0.7]; $\neg$a tick-to-tick transition law.)

\textbf{Req (Cor).} A2; C9, C10; D0.6, D0.7, D1.4, D1.5; CT0.6, CT0.9; exact binary capacity arithmetic [math; machine: \texttt{checks/\allowbreak{}verify\_\allowbreak{}promotion\_\allowbreak{}B\_\allowbreak{}rung.py} §B].

\textbf{Pf (Cor).}

%dsp-level:0
\[
\mathrm{CT0.9}:\ \sigma_e\in\{\pm1\}:\ 1\ \text{real edge}=1\ \text{bit}=\ln2\ \text{nats};
\]

%dsp-level:0
\[
\qquad \text{distinct edges}=\text{independent registers}\implies\mathsf E(\gamma)=\delta h\ln2\ [\mathrm{C9}]
\]

%dsp-level:0
\[
\mathrm{CT0.6}:\ C_{\mathrm{cap}}=16\ \text{register directions per channel crossing};
\]

%dsp-level:0
\[
\qquad \mathrm{C10}:\ |X|\ \text{alternatives}
\]

%dsp-level:0
\[
\mapsto |X|\ \text{directions},
\]

%dsp-level:0
\[
\ \text{Shannon}\ \log_2|X|:\ 16\ \text{directions}
\]

%dsp-level:0
\[
\iff |X|
\]

%dsp-level:0
\[
\le 16\iff
\]

%dsp-level:0
\[
\le 4\ \text{bits per tick}
\]

%dsp-level:0
\[
\delta\tau\ \text{ticks}:\ \text{supply}\le4\,\delta\tau\ \text{bits};
\]

%dsp-level:0
\[
\qquad \text{registered content}\le\text{supply}\implies\delta h\le4\,\delta\tau;
\]

%dsp-level:0
\[
\qquad \mathsf E\le\min(\delta h,\,4\,\delta\tau)\ln2
\]

%dsp-level:0
\[
>4\ \text{bits in one tick}:\ 2^{5}=32>16=C_{\mathrm{cap}}:\ \text{oversubscribed}\ [\mathrm{CT0.6}];
\]

%dsp-level:0
\[
\qquad \text{groups of}\ \le4:\ 2^{4}=16=C_{\mathrm{cap}}:\ \text{at capacity}\ [\mathrm{CT0.6}]\ \blacksquare
\]

\end{corollary}

\begin{theorem}[{\texorpdfstring{CT0.7 — Reality is the swap grading}{CT0.7 - Reality is the swap grading}}]\label{tag:CT0-7}
\textbf{Thm.} Let $\mathbb S\in\mathcal U(\mathcal H_p\otimes\mathcal H_q):\ |a\rangle\otimes|b\rangle\mapsto|b\rangle\otimes|a\rangle$. Then

%dsp-level:0
\[
(1)\quad \mathbb S=\tfrac12\bigl(E_0+\Sigma\bigr);
\]

%dsp-level:0
\[
\qquad \Pi_{\mathcal R}:=\Pi_{\vee^2}=\tfrac14\bigl(3E_0+\Sigma\bigr);
\]

%dsp-level:0
\[
\qquad \Pi_{\mathcal V}:=\Pi_{\wedge^2}=\tfrac14\bigl(E_0-\Sigma\bigr);
\]

%dsp-level:0
\[
\qquad \mathfrak T\Pi_{\mathcal R}=3,\ \ \mathfrak T\Pi_{\mathcal V}=1 .
\]

%dsp-level:0
\[
(2)\quad \forall X:\ S(X)=\mathbb S X\mathbb S .
\]

%dsp-level:0
\[
(3)\quad SX=X\implies[\mathbb S,X]=0\implies X(\vee^2)\subseteq\vee^2\ \wedge\ X(\wedge^2)\subseteq\wedge^2;
\]

%dsp-level:0
\[
SX=-X\implies\{\mathbb S,X\}=0\implies X(\vee^2)\subseteq\wedge^2\ \wedge\ X(\wedge^2)\subseteq\vee^2;
\]

%dsp-level:0
\[
\underbrace{\dim_{\mathbb R}\mathfrak H(\vee^2)+\dim_{\mathbb R}\mathfrak H(\wedge^2)}_{9+1=10}\ \oplus\ \underbrace{2\dim_{\mathbb C}\operatorname{Hom}(\wedge^2,\vee^2)}_{2\cdot3=6}\ =\ 16 .
\]

%dsp-level:0
\[
(4)\quad \mathbf F_{(+,\mathcal R)}=\mathcal E_{++};\quad
\mathbf F_{(-,\mathcal R)}=\mathcal E_{+-};\quad
\mathbf F_{(+,\mathcal V)}=\mathcal E_{-+};\quad
\mathbf F_{(-,\mathcal V)}=\mathcal E_{--}
\]

(first label of $\mathbf F_{(\pm,\cdot)}$: $T$-parity; second label: $S$-parity, $\mathcal R\leftrightarrow+$, $\mathcal V\leftrightarrow-$); $\ F_{J^+}=\mathcal E_{-+}=\operatorname{span}\{J_i\}$, $\ F_{J^-}=\mathcal E_{--}=\operatorname{span}\{K_i\}$: the $S$-odd blocks.

\textbf{Req.} A3, A8; C3, C8; D1.2, D3.10; CT0.6, CT1.1, CT1.2.

\textbf{Pf.}

%dsp-level:0
\[
(1)\ \ \mathbb S=\tfrac12\Bigl(I\otimes I+\textstyle\sum_i\sigma_i\otimes\sigma_i\Bigr)=\tfrac12\bigl(E_0+\Sigma\bigr)\quad[\mathrm{CT1.2}];
\]

basis check: $\ a=b:\ \tfrac12(1+1)|aa\rangle=|aa\rangle$; $\ a\ne b:\ \sum_i\sigma_i\otimes\sigma_i|ab\rangle=2|ba\rangle-|ab\rangle\implies\tfrac12(|ab\rangle+2|ba\rangle-|ab\rangle)=|ba\rangle$.

%dsp-level:0
\[
\mathbb S^2=\mathbb1\ \wedge\ \mathbb S^\dagger=\mathbb S\implies\operatorname{spec}(\mathbb S)=\{\pm1\},\ \ \text{eigenspaces}\ \vee^2,\ \wedge^2\quad[\mathrm{A3}]
\]

%dsp-level:0
\[
\Pi_{\mathcal R}=\mathbb1-\Pi_{\mathcal V}=\tfrac12(\mathbb1+\mathbb S)=\tfrac14\bigl(3E_0+\Sigma\bigr);
\]

%dsp-level:0
\[
\qquad \Pi_{\mathcal V}=\tfrac12(\mathbb1-\mathbb S)=\tfrac14\bigl(E_0-\Sigma\bigr)\quad[\mathrm{A3},\mathrm{A4}]
\]

%dsp-level:0
\[
\mathfrak T E_0=4,\ \mathfrak T\Sigma=0\implies\mathfrak T\Pi_{\mathcal R}=3=\dim\vee^2,\quad\mathfrak T\Pi_{\mathcal V}=1=\dim\wedge^2\quad[\mathrm{CT1.1}]
\]

%dsp-level:0
\[
(2)\ \ \mathbb S(\sigma_\mu\otimes\sigma_\nu)\mathbb S=\sigma_\nu\otimes\sigma_\mu=S(\sigma_\mu\otimes\sigma_\nu)\ [\mathrm{A3}];\quad\text{linearity}\implies S(X)=\mathbb S X\mathbb S\ \ \forall X
\]

%dsp-level:0
\[
(3)\ \ SX=X\Rightarrow\mathbb S X\mathbb S=X\Rightarrow\mathbb S X=X\mathbb S\ (\times\mathbb S\ \text{right});\qquad SX=-X\Rightarrow\mathbb S X=-X\mathbb S
\]

%dsp-level:0
\[
X=X^\dagger\ \wedge\ [\mathbb S,X]=0\implies X\ \text{preserves each}\ \mathbb S\text{-eigenspace};\qquad
\{\mathbb S,X\}=0\implies X:\ \vee^2\leftrightarrow\wedge^2
\]

%dsp-level:0
\[
\dim:\ \ \mathfrak H(\vee^2)\oplus\mathfrak H(\wedge^2)\cong\mathfrak H(\mathbb C^3)\oplus\mathfrak H(\mathbb C^1):\ 9+1=10;
\]

%dsp-level:0
\[
\qquad \operatorname{Hom}(\wedge^2,\vee^2)\cong\mathbb C^{3}:\ 2\cdot3=6;
\]

%dsp-level:0
\[
\qquad 10+6=16\ [\mathrm{CT0.6}]
\]

%dsp-level:0
\[
(4)\ \ S\text{-odd generators}=\{J_i\}\cup\{K_i\}\ (6)\ [\mathrm{CT1.2}];\quad S\text{-even}=\{E_0,\Sigma,\hat E^{ij},W_i\}\ (10)
\]

%dsp-level:0
\[
\mathrm{D1.2}:\ \text{seal moves A4-weight }\wedge^2\!\leftrightarrow\!\vee^2\implies\text{seal generators}
\]

%dsp-level:0
\[
= S\text{-odd}\ [(3)]
\]

%dsp-level:0
\[
\implies F_{J^\pm}
\]

%dsp-level:0
\[
= S\text{-odd blocks}\ [\mathrm{A8}]
\]

%dsp-level:0
\[
T\text{-parity}:\ J_i:+,\ K_i:-\ [\mathrm{CT1.2}]\implies F_{J^+}=\mathcal E_{-+},\ F_{J^-}=\mathcal E_{--}
\]

%dsp-level:0
\[
\mathcal E_{++}\supseteq\{E_0,\Sigma,\hat E^{ij}\}\mapsto\{F_{D^+},F_H,F_G\};\qquad
\mathcal E_{+-}=\operatorname{span}\{W_i\}\mapsto\{F_{W^+},F_{W^-},F_Z\}\quad[\mathrm{A8},\ \mathrm{C8},\ \mathrm{D3.10}].\ \blacksquare
\]

\end{theorem}

\begin{corollary}[{\texorpdfstring{CT0.7a — spin content}{CT0.7a - spin content}}]\label{tag:CT0-7a}
$\ \sum_iW_i^2=\tfrac32E_0+\tfrac12\Sigma$; $\ \sum_iW_i^2\big|_{\vee^2}=2=1(1{+}1)$; $\ \sum_iW_i^2\big|_{\wedge^2}=0=0(0{+}1)$; i.e. $\vee^2$: spin $1$, $\wedge^2$: spin $0$.

\textbf{Req (Cor).} CT0.7, CT1.2.

\textbf{Pf (Cor).}

%dsp-level:0
\[
\textstyle\sum_iW_i^2=\tfrac14\sum_i(\sigma_i\otimes\sigma_0+\sigma_0\otimes\sigma_i)^2=\tfrac14\bigl(3E_0+3E_0+2\Sigma\bigr)=\tfrac32E_0+\tfrac12\Sigma
\]

%dsp-level:0
\[
\Sigma=2\mathbb S-E_0\ [\mathrm{CT0.7}(1)];
\]

%dsp-level:0
\[
\qquad \vee^2:\ \mathbb S=+1\implies\Sigma=+1\implies\tfrac32+\tfrac12=2;
\]

%dsp-level:0
\[
\qquad \wedge^2:\ \mathbb S=-1\implies\Sigma=-3\implies\tfrac32-\tfrac32=0.\ \blacksquare
\]

\textbf{Phys.} The real/virtual split of A4 \textit{is} the exchange (swap) grading of the pair: real $=$ symmetric $=$ spin-1 triplet, virtual $=$ antisymmetric $=$ spin-0 singlet. The A8 force labels $(\pm,\mathcal R/\mathcal V)$ are precisely the two parities $(T,S)$; the Janus seals are exactly the operators that exchange the two grades, i.e. that create and destroy reality. “Existence” in this corpus is Bose exchange symmetry of a qubit pair.

\end{corollary}

\begin{theorem}[{\texorpdfstring{CT0.8 — Admissible range of the crystallization threshold}{CT0.8 - Admissible range of the crystallization threshold}}]\label{tag:CT0-8}
\textbf{Thm.} $\ \bigl[\mathrm{A4}\wedge\mathrm{A5}\wedge\mathrm{D0.1}\ \text{consistent}\bigr]\implies 0\le\varepsilon_q<\tfrac14;$ $\ \ $in particular $\varepsilon_q<1-\varepsilon_q$ (A4's two branches are disjoint); $\qquad\varepsilon_q=0\implies\bigl[(p,q)\in\mathcal R\iff\operatorname{supp}\rho_{pq}\subseteq\vee^2\ \ \vee\ \ \rho_{pq}=\Pi_{\mathcal V}\bigr]$ (branch reads $\alpha_{pq}=0\ /\ \pi$ [A4]); and $\forall$ theorems of the corpus: A4 enters only via CT0.1 (valid $\forall\varepsilon_q\in[0,1]$) and CT0.2 (valid $\iff\varepsilon_q\in[0,\tfrac14)$).

\textbf{Req.} A1, A4, A5; D0.1, D0.2, D0.4; CT0.2, CT0.7.

\textbf{Pf.}

%dsp-level:0
\[
\bigl|\Pi_{\mathcal V}\rho_\circ\bigr|=\mathfrak T\bigl(\Pi_{\mathcal V}\cdot\tfrac14 I\bigr)=\tfrac14\,\mathfrak T\Pi_{\mathcal V}=\tfrac14\quad[\mathrm{A5},\ \mathrm{CT0.7}(1):\ \mathfrak T\Pi_{\mathcal V}=1]
\]

%dsp-level:0
\[
\mathrm{D0.1}:\ \text{vacuum}\iff n_{\mathcal R}=0\implies\rho_\circ\notin\mathcal R\ [\mathrm{CT0.2}]\implies\tfrac14>\varepsilon_q\quad[\mathrm{A4}]
\]

%dsp-level:0
\[
\varepsilon_q\ge0\ [\mathrm{D0.4}];\qquad\bigl|\Pi_{\mathcal V}\rho\bigr|\ge0\ [\mathrm{D0.2}]\implies\varepsilon_q=0\ \text{admissible}
\]

%dsp-level:0
\[
\text{branch 2 at the vacuum}:\ \tfrac14
\]

%dsp-level:0
\[
\ge 1-\varepsilon_q\iff\varepsilon_q\ge\tfrac34:\ \text{excluded by}\ \varepsilon_q<\tfrac14\implies\rho_\circ\ \text{virtual in both branches}
\]

%dsp-level:0
\[
\text{branch disjointness}:\ 0\le\varepsilon_q<\tfrac14\implies\varepsilon_q<\tfrac12<1-\varepsilon_q\implies[0,\varepsilon_q]\cap[1-\varepsilon_q,1]=\varnothing
\]

%dsp-level:0
\[
\varepsilon_q=0:\ \ \mathfrak T\bigl(\Pi_{\mathcal V}\rho_{pq}\bigr)=0\ \wedge\ \Pi_{\mathcal V}\succeq0\ \wedge\ \rho_{pq}\succeq0\iff\Pi_{\mathcal V}\rho_{pq}\Pi_{\mathcal V}=0\iff\operatorname{supp}\rho_{pq}\subseteq\vee^2\quad[\mathrm{A1}]
\]

%dsp-level:0
\[
\varepsilon_q
\]

%dsp-level:0
\[
= 0,
\]

%dsp-level:0
\[
\ \text{branch 2}:\ \ \mathfrak T\bigl(\Pi_{\mathcal V}\rho_{pq}\bigr)
\]

%dsp-level:0
\[
= 1\ \wedge\ \mathfrak T\rho_{pq}
\]

%dsp-level:0
\[
= 1\ \wedge\ \rho_{pq}\succeq0\ \wedge\ \operatorname{rk}\Pi_{\mathcal V}
\]

%dsp-level:0
\[
= 1\ \ [\mathrm{CT0.7}(1)]\iff\rho_{pq}\qquad = \Pi_{\mathcal V}
\]

%dsp-level:0
\[
\text{A4-dependence of all later tags}=\{\mathrm{CT0.1},\ \mathrm{CT0.2}\}\implies\text{range}\ [0,\tfrac14)\ \text{suffices; sharp point}\ \varepsilon_q=0.\ \blacksquare
\]

\textbf{Phys.} The crystallization threshold is not a free constant: demanding that the vacuum be all-virtual confines it to $0\le\varepsilon_q<\tfrac14$ — which simultaneously keeps A4's two branches disjoint ($\varepsilon_q<1-\varepsilon_q$). The sharp choice $\varepsilon_q=0$ says a pair is real exactly when its antisymmetric weight is exactly zero (aligned branch, $\alpha_{pq}=0$) or the pair is the pure singlet (antialigned branch, $\alpha_{pq}=\pi$). No result depends on where inside the range the tolerance sits — the theory is tolerance-independent.

\end{theorem}

\begin{theorem}[{\texorpdfstring{CT0.9 — The edge phase alphabet is binary}{CT0.9 - The edge phase alphabet is binary}}]\label{tag:CT0-9}
\textbf{Thm.} $\ \phi_e\in\Phi=\{0,\pi\}\ (\mathrm{mod}\ 2\pi);\qquad\sigma_e=e^{i\phi_e}=(-1)^{\phi_e/\pi}\in\{+1,-1\};\qquad 1\ \text{real edge}=1\ \text{bit}=\ln2\ \text{nats}.$

\textbf{Req.} A1, A3, A4, A7; C9, C10; D0.5, D1.1; $\dim\wedge^2=\binom{d}{2}$ [linear algebra].

\textbf{Pf.}

%dsp-level:0
\[
e\in E_{\mathcal R}\mapsto\phi_e\ [\mathrm{A7}];
\]

%dsp-level:0
\[
\qquad \mathrm{D1.1}:\ \phi_e:=\text{phase of the registration transport's action on the reality line}\ \wedge^2
\]

%dsp-level:0
\[
\mathrm{A1}:\ \dim\mathcal H_q=d=2\implies\dim\wedge^2=\tbinom{2}{2}=1\quad(\text{the reality line is 1-dimensional}\iff\text{qubits are}\ \mathbb C^2)
\]

%dsp-level:0
\[
\mathrm{A4}:\ \text{reality test}
\]

%dsp-level:0
\[
= \text{extremality of the}\ \wedge^2\ \text{weight (either branch)}\implies\text{registration transport preserves the}\ (S,T)\text{-grading}\ [\mathrm{A3},\ \mathrm{D1.1}]
\]

%dsp-level:0
\[
\mathrm{D1.1}/\mathrm{C10}\ \text{reality principle}:\ \text{transport preserves the real frame structure}\implies\text{action on}\ \wedge^2_{\mathbb R}\ (1\text{-dim real line})
\]

%dsp-level:0
\[
= \text{multiplication by}\ c\in\mathbb R
\]

%dsp-level:0
\[
\text{unitarity}\implies|c|=1\implies c\in\{+1,-1\}\implies e^{i\phi_e}\in\{\pm1\}\implies\phi_e\in\{0,\pi\}\ (\mathrm{mod}\ 2\pi);
\]

%dsp-level:0
\[
\qquad \sigma_e=(-1)^{\phi_e/\pi}\ [\mathrm{D0.5}]
\]

%dsp-level:0
\[
|\Phi|=2\implies\log_2|\Phi|=1\ \text{bit}=\ln2\ \text{nats}\quad[\mathrm{C9}].\ \blacksquare
\]

\textbf{Rmk (repair).} The prior proof argued from the centralizer of $\{S,T\}$ alone — insufficient, since the complex centralizer of the swap is $U(3)\times U(1)$ and contains continuous phases. The binary alphabet follows once the reality principle of D1.1/C10 (registrations are real, grading-preserving frame data — quantum data classicalizes at registration) is explicit. The continuous counterexamples are complex-transport artifacts excluded by that principle.

\textbf{Phys.} The only internal datum a real edge can carry is a single sign — and this is \textit{qubit-specific}: the reality line $\wedge^2$ is one-dimensional exactly because the carriers are $\mathbb C^2$ ($\dim\wedge^2=\binom d2=1\iff d=2$); for higher-dimensional carriers continuous grading-preserving phases would survive. Every edge of the emergent graph stores exactly one bit; all later charge and phase structure is built from these signs.

\end{theorem}

\begin{theorem}[{\texorpdfstring{CT0.10 — The two sign structures are independent}{CT0.10 - The two sign structures are independent}}]\label{tag:CT0-10}
\textbf{Thm.} With $\ P(C):=\prod_{e\in E(C)}\sigma_e\ $ and $\ \pi(C):=(-1)^{\aleph(\sigma,C)}$ :

%dsp-level:0
\[
\neg\bigl[P(C)\ \text{determines}\ \pi(C)\bigr]\ \wedge\ \neg\bigl[\pi(C)\ \text{determines}\ P(C)\bigr];\qquad
L\ge4:\ \ (P,\pi)\twoheadrightarrow\{\pm1\}^2;
\]

%dsp-level:0
\[
\Longrightarrow\quad w(C)=\mathcal J(C)\prod_{e\in E(C)}G(e)\cdot(-1)^{\aleph(\sigma,C)}\ \ [\mathrm{A8}]\ \text{double-counts no edge-sign data}.
\]

\textbf{Req.} A8; D0.5, D3.3.

\textbf{Pf.}

%dsp-level:0
\[
\aleph(\sigma,C)=\#\{i<L:\ \sigma_i\ne\sigma_{i+1}\}\quad[\mathrm{D3.3}]\ :\ \text{function of consecutive differences only};
\]

%dsp-level:0
\[
\qquad P=\prod_{i}\sigma_i\ :\ \text{function of the signs}
\]

%dsp-level:0
\[
\text{witness}_1\ [\pi\ \text{equal},\ P\ \text{opposite}]:\quad(+,-,+):\ \aleph=2,\ P=-1;\qquad(-,+,-):\ \aleph=2,\ P=+1
\]

%dsp-level:0
\[
\text{witness}_2\ [P\ \text{equal},\ \pi\ \text{opposite}]:\quad(+,+,+,+):\ P=+1,\ \aleph=0;\qquad(+,+,-,-):\ P=+1,\ \aleph=1
\]

%dsp-level:0
\[
\Longrightarrow\quad P\not\Rightarrow\pi\ \wedge\ \pi\not\Rightarrow P;
\]

%dsp-level:0
\[
\qquad L\ge4:\ \{(P,\pi)\}=\{(+,+),(+,-),(-,+),(-,-)\}\ \text{attained}\implies(P,\pi)\ \text{surjective}
\]

%dsp-level:0
\[
\Longrightarrow\quad\text{the two factors of}\ w(C)\ \text{carry independent information}.\ \blacksquare
\]

\textbf{Phys.} The product of edge signs and the transition parity measure different things — net sign content versus the number of domain walls between signs — and neither fixes the other. The A8 weight is therefore free of double counting: bulk sign and wall-count contribute independently to every configuration amplitude.

\end{theorem}

\begin{theorem}[{\texorpdfstring{CT0.11 — Maximal accessible rung}{CT0.11 - Maximal accessible rung}}]\label{tag:CT0-11}
\textbf{Thm.} $\ \text{channel-admissible}(v)\iff D_{\mathrm{tick}}\bigl(n(v)\bigr)\le C_{\mathrm{cap}}\iff n(v)\le3;$ accessible vertex sets $=\{V_0,V_1,V_2,V_3\}$ (types $\{\Xi_0,\dots,\Xi_3\}$); $D_{\mathrm{tick}}(3)=C_{\mathrm{cap}}$ (saturation); $\forall v:\ n(v)\ge4\implies O\bigl(n(v)\bigr)>0\implies F_{D^-}$ [D4.1].

\textbf{Req.} D0.8, D0.9, D4.1; CT0.5, CT0.6.

\textbf{Pf.}

%dsp-level:0
\[
\text{admissible}\iff D_{\mathrm{tick}}(n)\le C_{\mathrm{cap}}\quad[\mathrm{D0.9}]
\]

%dsp-level:0
\[
2^{\,n+1}\le16\iff n+1\le4\iff n\le3\quad[\mathrm{CT0.5},\ \mathrm{CT0.6}]
\]

%dsp-level:0
\[
n=3:\ \ 2^{4}=16=C_{\mathrm{cap}}\ (\text{tight});\qquad
n\ge4:\ \ O(n)>0\ [\mathrm{CT0.6}]\implies\text{configuration}\in F_{D^-}\ [\mathrm{D4.1}].\ \blacksquare
\]

\end{theorem}

\begin{corollary}[{\texorpdfstring{CT0.11a — three slots}{CT0.11a - three slots}}]\label{tag:CT0-11a}
$\ \max\{\mathcal H^{(n)}_{\mathrm{slot}}:\ \text{admissible}\}=\mathcal H^{(3)}_{\mathrm{slot}}\cong\mathbb C^3;\qquad 3=\lfloor\log_2 C_{\mathrm{cap}}\rfloor-1=\lfloor\log_216\rfloor-1;$ i.e. $N_{\text{colour}}=3$ is derived, not chosen.

\textbf{Req (Cor).} D3.1; CT0.6, CT0.11.

\textbf{Pf (Cor).} $\mathcal H^{(n)}_{\mathrm{slot}}\cong\mathbb C^n$ [D3.1]; $n\le3$ [CT0.11] $\Rightarrow\max=\mathbb C^3$; $\lfloor\log_2 16\rfloor-1=3$. $\blacksquare$

\end{corollary}

\begin{corollary}[{\texorpdfstring{CT0.11b — directed rung of bidirectional skeletons}{CT0.11b - directed rung of bidirectional skeletons}}]\label{tag:CT0-11b}
A channel $\{p,q\}$, $p\ne q$, is \textit{bidirectional} when both orientations are real: $(p,q)\in\mathcal R\ \wedge\ (q,p)\in\mathcal R$ [A4]. Let $S$ be a finite real-edge subgraph whose every used channel is bidirectional, and $d_S(v):=\#\{\text{channels of}\ S\ \text{incident on}\ v\}$. Then:

%dsp-level:0
\[
n(v)\ \ge\ 2\,d_S(v)\qquad[\text{each bidirectional channel adds}\ 1\ \text{to}\ \partial^{\leftarrow}\ \text{and}\ 1\ \text{to}\ \partial^{\rightarrow}\ \text{at each endpoint}]
\]

%dsp-level:0
\[
S\ \text{exhausts the real edges at}\ v\implies n(v)=2\,d_S(v);
\]

%dsp-level:0
\[
\qquad d_S(v)\ge2\implies n(v)\ge4\implies D_{\mathrm{tick}}(n)=2^{n+1}\ge32>16\implies O\bigl(n(v)\bigr)>0\implies F_{D^-}\ [\mathrm{D4.1}]
\]

%dsp-level:0
\[
\text{unoversubscribed}\iff n(v)\le3\ [\mathrm{CT0.6}]:
\]

%dsp-level:0
\[
\qquad \text{bidirectional}\ S:\ d_S(v)\le1\ (n\in\{0,2\});
\]

%dsp-level:0
\[
\qquad \text{mixed}:\ n(v)=2\,d^{\mathrm{bi}}_S(v)+d^{\mathrm{uni}}_S(v)\le3 .
\]

\textbf{Req (Cor).} A4; D0.8, D4.1; CT0.5, CT0.6, CT0.11; directed-degree enumeration [math; machine: \texttt{checks/\allowbreak{}verify\_\allowbreak{}promotion\_\allowbreak{}B\_\allowbreak{}rung.py} §A].

\textbf{Pf (Cor).}

%dsp-level:0
\[
\mathrm{A4}:\ \mathcal R\text{-membership per ordered pair};
\]

%dsp-level:0
\[
\quad\text{bidirectional channel}
\]

%dsp-level:0
\[
= \{e_{pq},e_{qp}\}\subseteq\mathcal R:\ \partial^{\leftarrow}\!+1,
\]

%dsp-level:0
\[
\ \partial^{\rightarrow}\!+1\ \text{at each endpoint}
\]

%dsp-level:0
\[
\implies n(v)
\]

%dsp-level:0
\[
\ge 2\,
\]

%dsp-level:0
\[
d_S(v)\ [\mathrm{D0.8}]
\]

%dsp-level:0
\[
\text{equality}\iff S\ \text{exhausts}\ v\text{'s incident real edges};
\]

%dsp-level:0
\[
\qquad d_S\ge2:\ n\ge4:\ \mathrm{CT0.5}:\ D_{\mathrm{tick}}(n)=2^{n+1}\ge32;\quad\mathrm{CT0.6}:\ O(n)=2^{n+1}-16>0
\]

%dsp-level:0
\[
\mathrm{CT0.11}:\ n\ge4\implies F_{D^-}\ [\mathrm{D4.1}];
\]

%dsp-level:0
\[
\qquad O(n)=0\iff n\le3\ [\mathrm{CT0.6}];
\]

%dsp-level:0
\[
\qquad \text{mixed count}:\ \text{additivity of}\ \partial^{\leftarrow},\partial^{\rightarrow}\ [\mathrm{D0.8}]\ \blacksquare
\]

\textbf{Phys.} The 16-register channel budget admits vertices of degree at most three: rungs $0$–$3$ are accessible, rung $3$ exactly exhausts the channel, and everything of degree $\ge4$ is oversubscribed and falls to the dark sector. The number three — hence the threefold slot (“colour”) structure of the theory — is forced by capacity, not assumed.

\textbf{Note.} At rungs $n\ge4$ no aspect realization $\iota_n$ exists: the unique derived read of rung-$\ge4$ slot content is $D^-$-absorption at the ledger level [CT4.4a (Part 2) — forward pointer, non-load-bearing].

\end{corollary}

\begin{theorem}[{\texorpdfstring{CT0.12 — Aspect–capacity match}{CT0.12 - Aspect-capacity match}}]\label{tag:CT0-12}
\textbf{Thm.}

%dsp-level:0
\[
\begin{array}{c|c|c}
\text{aspect} & \text{block} & \dim_{\mathbb R}\\ \hline
F_{J^+} & \mathcal E_{-+}=\operatorname{span}\{J_i\} & 3\\
F_{J^-} & \mathcal E_{--}=\operatorname{span}\{K_i\} & 3\\
F_G & \operatorname{span}\{\hat E^{ij}\}\ (\text{traceless}) & 5\\
F_W & \mathcal E_{+-}=\operatorname{span}\{W_i\} & 3\\
F_H & \operatorname{span}\{\Sigma\} & 1\\
F_{D^+} & \operatorname{span}\{E_0\} & 1\\ \hline
\text{total} & \mathfrak H(\mathbb C^4) & 16=C_{\mathrm{cap}}
\end{array}
\]

\textbf{Req.} A8; C8; CT0.6, CT0.7(4), CT1.1, CT1.2.

\textbf{Pf.}

%dsp-level:0
\[
\text{block assignments}\quad[\mathrm{CT1.2},\ \mathrm{CT0.7}(4)];\qquad
3+3+5+3+1+1=16=\dim_{\mathbb R}\mathfrak H(\mathbb C^4)\quad[\mathrm{CT1.1}]
\]

%dsp-level:0
\[
16=C_{\mathrm{cap}}\quad[\mathrm{CT0.6}].\ \blacksquare
\]

\textbf{Rmk.} $\mathbf F_{\mathrm{st}}=F_{\mathrm{EM}}\oplus F_{\mathrm{strong}}\oplus F_{D^-}$ acts on vertex/slot registers [D3.1], not on the pair channel; the CT0.6 bound constrains the pair channel only.

\textbf{Phys.} The six pair-sector force aspects of A8 tile the 16 dimensions of the two-qubit operator space exactly: no part of the channel is forceless and no force overruns it. The stack aspects (electromagnetic, colour, dark-saturation) live on separate vertex registers and are not bounded by this count.

\end{theorem}

\begin{theorem}[{\texorpdfstring{CT0.13 — Aspect activation and the active sets}{CT0.13 - Aspect activation and the active sets}}]\label{tag:CT0-13}
\textbf{Thm.}

%dsp-level:0
\[
\mathcal A(0)=\{J\};\qquad
\mathcal A(1)=\{J,\,G,\,W,\,H,\,D^+\};\qquad
\mathcal A(2)=\mathcal A(1)\cup\{\mathrm{EM}\};
\]

%dsp-level:0
\[
\mathcal A(3)=\mathcal A(2)\cup\{\mathrm{strong}\};\qquad
\mathcal A(n\ge4)=\mathcal A(3)\cup\{D^-\};
\]

with $J:=\{J^+,J^-\}$ and $W:=F_W=F_{W^+}\oplus F_{W^-}\oplus F_Z$ [C8]; $\ \neg\exists$ independent aspect named $Z$ at any rung: neutral seed $W_3\in F_W$; physical $Z$: D3.10, EW.24.

\textbf{Req.} A7, A8; C8; D0.8, D1.12, D2.12, D3.1, D3.2, D3.3, D4.1, D4.3; CT0.1, CT0.7, CT0.11, CT1.2, CT1.4, CT1.10; $\bigl[L^\dagger L=I_n\iff L\in U(n);\ \ \mathfrak u(n)=\mathfrak{su}(n)\oplus\mathfrak u(1)_{\mathrm{centre}}$ — linear algebra; cf. CT3.4, CT3.5$\bigr]$.

\textbf{Pf.}

%dsp-level:0
\[
n=0:\ \ v\in V_0\implies\neg\,\exists e\in E_{\mathcal R}:\ e\ni v\ [\mathrm{D0.8}]\implies\forall\ \text{incident pairs}:\ \in\mathcal V\ [\mathrm{CT0.1}]
\]

%dsp-level:0
\[
\mathrm{CT0.7}:\ \text{operators on the virtual block}=S\text{-odd}=\{F_{J^+},F_{J^-}\}\ [\mathrm{A8}]\implies\mathcal A(0)=\{J\}
\]

%dsp-level:0
\[
n=1:\ \ \bigl|\{e\in E_{\mathcal R}:e\ni v\}\bigr|=1\implies\vee^2\ \text{populated};\qquad
\mathrm{CT0.7}:\ S\text{-even act within}\ \vee^2
\]

%dsp-level:0
\[
\mathrm{CT1.2},\ \mathrm{CT1.10}:\ S\text{-even}=\{E_0,\ \Sigma,\ \hat E^{ij}_{(5)},\ W_i\}\xrightarrow[\mathrm{A8}]{}\{F_{D^+},\ F_H,\ F_G,\ F_W\}
\]

%dsp-level:0
\[
\implies\mathcal A(1)=\{J,G,W,H,D^+\};\qquad
\dim F_W=3\ \text{single block}\implies W_3\ \neg\text{independent}\implies\neg\exists\ \text{aspect}\ Z
\]

%dsp-level:0
\[
\text{slot algebras}:\ \ \mathcal H^{(n)}_{\mathrm{slot}}\cong\mathbb C^n\ [\mathrm{D3.1}];\quad
\text{probability-preserving}\iff L^\dagger L=I_n\iff L\in U(n)
\]

%dsp-level:0
\[
X=\Bigl(X-\tfrac{\mathfrak T X}{n}I_n\Bigr)+\tfrac{\mathfrak T X}{n}I_n\ \text{(unique trace split)}\implies\mathfrak u(n)=\mathfrak{su}(n)\oplus\mathfrak u(1)_{\mathrm{centre}}
\]

%dsp-level:0
\[
\quad [\text{linear algebra; }n=3:\ \mathrm{CT3.4},\mathrm{CT3.5}]
\]

%dsp-level:0
\[
n = 1\ \text{slot}:\ \ \mathfrak u(1)\ \text{on one slot}
\]

%dsp-level:0
\[
= \text{global phase};
\]

%dsp-level:0
\[
\ \ \neg\exists\ \text{second slot}\implies\neg\text{configuration-dependent}
\]

%dsp-level:0
\[
\implies
\]

%dsp-level:0
\[
= F_{D^+}
\]

%dsp-level:0
\[
= \Lambda\mathbb I\ [\mathrm{D4.3}]\ \text{(counted)};
\]

%dsp-level:0
\[
\quad\neg\text{new}\ [\mathrm{D3.2}]
\]

%dsp-level:0
\[
n=2:\ \ \mathfrak u(2)=\mathfrak{su}(2)\oplus\mathfrak u(1):
\]

%dsp-level:0
\[
\quad \iota_2\bigl(\mathfrak{su}(2)\bigr)\cong\mathfrak{su}(2)\cong\operatorname{span}\{W_i\}\ [\mathrm{CT1.4}]:\ \text{same abstract class realized at rung}\ 1
\]

%dsp-level:0
\[
\implies\neg\text{strictly new}\ [\mathrm{D2.12}\ (\iota_n\ \text{bridge}),\ \mathrm{D3.2}]
\]

%dsp-level:0
\[
\text{centre measurable}:\ \ V_2\ni\text{type}(1,1)\ [\mathrm{A7}]\implies\text{relative in/out phase defined}=f_v\not\equiv0\ [\mathrm{D3.3}]\implies\text{strictly new}=F_{\mathrm{EM}}\ [\mathrm{A8}]
\]

%dsp-level:0
\[
\implies\mathcal A(2)=\mathcal A(1)\cup\{\mathrm{EM}\};
\]

%dsp-level:0
\[
\qquad f_v\ \text{requires}\ \partial^{\leftarrow}\!\ge1\ \wedge\ \partial^{\rightarrow}\!\ge1\implies n\ge2\ [\mathrm{D3.3},\mathrm{D0.8}]\implies\mathrm{EM}\notin\mathcal A(0)\cup\mathcal A(1)
\]

%dsp-level:0
\[
n=3:\ \ \mathfrak u(3)=\mathfrak{su}(3)\oplus\mathfrak u(1):\quad
\mathfrak u(1)\ \text{duplicates}\ F_{\mathrm{EM}}\ (n=2)\implies\neg\text{new};
\]

%dsp-level:0
\[
\qquad \mathfrak{su}(3)\ \text{strictly new}=F_{\mathrm{strong}}\ [\mathrm{A8}]
\]

%dsp-level:0
\[
\implies\mathcal A(3)=\mathcal A(2)\cup\{\mathrm{strong}\}
\]

%dsp-level:0
\[
n\ge4:\ \ \text{oversubscribed}\ [\mathrm{CT0.11}]\implies F_{D^-}\ \text{active there}\wedge\text{only there}\ [\mathrm{D4.1}];
\]

%dsp-level:0
\[
\qquad \text{rung-3 structure persists}\implies\mathcal A(n\ge4)=\mathcal A(3)\cup\{D^-\}.\ \blacksquare
\]

\textbf{Phys.} Forces switch on one at a time as vertices acquire edges. Rung $0$ has only the Janus seals that create and destroy reality; rung $1$ lights up gravity ($G$), weak ($W$), Higgs ($H$) and uniform dark ($D^+$); at rung $2$ the local phase along a pass-through becomes measurable — electromagnetism; at rung $3$ the traceless three-slot algebra appears — colour; beyond, configurations lock into dark saturation ($D^-$). The gauge sequence of the corpus is channel bookkeeping, and no rung ever carries an independent $Z$ aspect — only the seed $W_3$ inside $F_W$.

\end{theorem}

\begin{theorem}[{\texorpdfstring{CT0.14 — Mask flag and block dimensions}{CT0.14 - Mask flag and block dimensions}}]\label{tag:CT0-14}
\textbf{Thm.}

%dsp-level:0
\[
\mathcal K_0\subset\mathcal K_1\subset\mathcal K_2\subset\mathcal K_3\subset\mathcal K_{n\ge4}\quad(\text{strict});\qquad
\dim\mathcal K_0=6,\ \ \dim\mathcal K_1=16,\ \ \dim\mathcal K_2=17,\ \ \dim\mathcal K_3=25;
\]

%dsp-level:0
\[
\bigl(\mathcal P_M^{(n)}\bigr)^2=\mathcal P_M^{(n)}=\bigl(\mathcal P_M^{(n)}\bigr)^\dagger;\qquad
\mathfrak T\,\mathcal P_M^{(n)}=\dim\mathcal K_n;\qquad
\dim\mathcal K^{(\mathrm{EM})}=1,\ \ \dim\mathcal K^{(\mathrm{strong})}=8 .
\]

\textbf{Req.} A8; D2.3; CT0.12, CT0.13; CT3.8 (\textbf{F} — forward pointer: the block-dimension row re-verifies the gluon count).

\textbf{Pf.}

%dsp-level:0
\[
\mathrm{CT0.13}:\ \ \mathcal A(0)\subset\mathcal A(1)\subset\mathcal A(2)\subset\mathcal A(3)\subset\mathcal A(n\ge4)\ \ (\text{strict})
\]

%dsp-level:0
\[
\mathrm{D2.3}:\ \ \mathcal K_n=\bigoplus_{A\in\mathcal A(n)}\mathcal K^{(A)}\implies\mathcal K_0\subset\mathcal K_1\subset\mathcal K_2\subset\mathcal K_3\subset\mathcal K_{n\ge4}\ \ (\text{strict})
\]

%dsp-level:0
\[
\dim\mathcal K_0=3+3=6\ [\mathrm{CT0.12}];\qquad
\dim\mathcal K_1=6+(5+3+1+1)=16\ [\mathrm{CT0.12}]
\]

%dsp-level:0
\[
\dim\mathcal K^{(\mathrm{EM})}=\dim Z\bigl(\mathfrak u(2)\bigr)=1\ [\mathrm{CT0.13}\ \text{step}\ n{=}2:\ \text{centre spanned by}\ I_2]\implies\dim\mathcal K_2=17
\]

%dsp-level:0
\[
\dim\mathcal K^{(\mathrm{strong})}=\dim\mathfrak{su}(3)=8\ [\mathrm{CT3.8}]\implies\dim\mathcal K_3=25
\]

%dsp-level:0
\[
\mathrm{D2.3}:\ \mathcal P_M^{(n)}\ \text{orthogonal projection}\implies P^2=P=P^\dagger,\quad\mathfrak T P=\dim\operatorname{Im}P=\dim\mathcal K_n.\ \blacksquare
\]

\textbf{Rmk.} $\dim\mathcal K_1=16=C_{\mathrm{cap}}$: rung $1$ already exhausts the pair channel; growth at rungs $2,3$ is carried entirely by stack registers [CT0.12 Rmk].

\textbf{Phys.} The rung masks form a strict flag $6\subset16\subset17\subset25$: each new rung adds exactly the block of its new force — one dimension for electromagnetism at rung $2$, eight for colour at rung $3$. “What the universe can do at rung $n$” is literally a nested subspace of force space, and the masks are the orthogonal projections onto it.

\end{theorem}

\phantomsection
\section[Tier 1A — The Pair Operator Algebra]{\texorpdfstring{Tier 1A — The Pair Operator Algebra}{Tier 1A - The Pair Operator Algebra}}

\begin{theorem}[{\texorpdfstring{CT1.0 — Signature classification of the $T$-grading}{CT1.0 - Signature classification of the -grading}}]\label{tag:CT1-0}
\textbf{Thm.} Let $\eta\in\{\pm1\}^4$ have exactly one odd entry [A3]. Put $\eta^{(\mathfrak o)}:=(+,\dots,-_{\mathfrak o},\dots,+)$, $\zeta_i:=\eta^0\eta^i$, $\zeta_0:=1$, $\mathsf D_\zeta:\sigma_\mu\mapsto\zeta_\mu\sigma_\mu$, $\mathrm{sf}(X):=\sigma_2X^{\mathrm T}\sigma_2$. Then:

$(1)$ $\eta=\pm\eta^{(\mathfrak o)}$ for exactly one $\mathfrak o\in\{0,1,2,3\}$ and exactly one sign.

$(2)$ $T[\eta]=T[-\eta]$, and $T=\mathsf D_\zeta\otimes\mathsf D_\zeta$.

$(3)$ $[S,T]=0$ and $(\dim\mathcal E_{++},\dim\mathcal E_{+-},\dim\mathcal E_{-+},\dim\mathcal E_{--})=(7,3,3,3)$ for every $\mathfrak o$.

$(4)$ $\mathfrak o=0$: $T=\mathrm{sf}\otimes\mathrm{sf}$ (inversion of both Bloch vectors). $\mathfrak o=k\ge1$: $\mathsf D_\zeta$ reflects the Bloch axis $k$. The three maps $\mathfrak o\in\{1,2,3\}$ are conjugate under $\mathrm{Ad}_{U\otimes U}$, $U\in\mathrm{SU}(2)$; the map $\mathfrak o=0$ is conjugate to none of them.

$(5)$ $T\,\mathrm{Ad}_{U\otimes U}=\mathrm{Ad}_{U\otimes U}\,T$ for all $U\in\mathrm{SU}(2)$ $\iff\mathfrak o=0$.

\textbf{Req.} A3; C1; $[\text{math}]$: $\mathrm{Ad}_U$ fixes $\sigma_0$ and acts on $\operatorname{span}\{\sigma_i\}$ by $R(U)\in SO(3)$, every $R\in SO(3)$ arising; $\sigma_2U^{\mathrm T}\sigma_2=U^\dagger$ for $U\in\mathrm{SU}(2)$.

\textbf{Pf.}

%dsp-level:0
\[
\text{(1)}:\ \text{odd entry}=-1\implies\eta=\eta^{(\mathfrak o)};\quad\text{odd entry}=+1\implies\eta=-\eta^{(\mathfrak o)};\quad\mathfrak o\ \text{is the unique odd slot}\quad[\mathrm{A3}]
\]

%dsp-level:0
\[
\text{(2)}:\ \ \eta^\mu\eta^\nu=(-\eta^\mu)(-\eta^\nu)\implies T[\eta]=T[-\eta];
\]

%dsp-level:0
\[
\qquad \zeta_\mu\zeta_\nu=(\eta^0\eta^\mu)(\eta^0\eta^\nu)=\eta^\mu\eta^\nu\ \ [(\eta^0)^2=1;\ \zeta_0=\eta^0\eta^0]
\]

%dsp-level:0
\[
\implies(\mathsf D_\zeta\otimes\mathsf D_\zeta)(\sigma_\mu\otimes\sigma_\nu)=\zeta_\mu\zeta_\nu\,\sigma_\mu\otimes\sigma_\nu=T(\sigma_\mu\otimes\sigma_\nu)\ \Rightarrow\ T=\mathsf D_\zeta\otimes\mathsf D_\zeta
\]

%dsp-level:0
\[
\text{(3)}:\ \ ST(\sigma_\mu\otimes\sigma_\nu)=\eta^\mu\eta^\nu\,\sigma_\nu\otimes\sigma_\mu=TS(\sigma_\mu\otimes\sigma_\nu)\implies[S,T]=0;
\]

%dsp-level:0
\[
\qquad \sigma_\mu\otimes\sigma_\mu:\ (S,T)=(+,+);\quad
\sigma_\mu\otimes\sigma_\nu\pm\sigma_\nu\otimes\sigma_\mu\ (\mu<\nu):\ (S,T)=(\pm,\eta^\mu\eta^\nu)
\]

%dsp-level:0
\[
n_-:=\#\{\mu<\nu:\eta^\mu\eta^\nu=-1\}=\#\bigl\{\{\mathfrak o,\nu\}:\nu\ne\mathfrak o\bigr\}=3;\qquad n_+=6-n_-=3
\]

%dsp-level:0
\[
\dim\mathcal E_{++}=4+n_+=7,\quad\dim\mathcal E_{+-}=n_-=3,\quad\dim\mathcal E_{-+}=n_+=3,\quad\dim\mathcal E_{--}=n_-=3
\]

%dsp-level:0
\[
\text{(4)}:\ \ \mathfrak o=0:\ \zeta_i=\eta^0\eta^i=-1\ \forall i;
\]

%dsp-level:0
\[
\qquad \mathrm{sf}(\sigma_0)=\sigma_0,\ \ \mathrm{sf}(\sigma_i)=\sigma_2\sigma_i^{\mathrm T}\sigma_2=-\sigma_i\ (i=1,2,3)\implies\mathsf D_\zeta=\mathrm{sf}
\]

%dsp-level:0
\[
\mathfrak o=k\ge1:\ \zeta_k=\eta^0\eta^k=-1,\ \ \zeta_i=+1\ (i\ne k)\implies\mathsf D_\zeta=\text{reflection of Bloch axis }k
\]

%dsp-level:0
\[
\mathrm{Ad}_U\,\mathsf D_\zeta\,\mathrm{Ad}_U^{-1}=\mathsf D_{\zeta'}\ \text{with}\ \operatorname{diag}(\zeta')=R(U)\operatorname{diag}(\zeta)R(U)^{-1}\quad[\text{Req}]
\]

%dsp-level:0
\[
\text{axis reflections}:\ R\ \text{permuting the axes}\implies T^{(1)}\sim T^{(2)}\sim T^{(3)};
\]

%dsp-level:0
\[
\qquad \operatorname{spec}(-\mathbb1_3)=\{-1,-1,-1\}\ne\{-1,+1,+1\}\implies T^{(0)}\not\sim T^{(k)}
\]

%dsp-level:0
\[
\text{(5)}\ (\Leftarrow):\ \ \mathrm{sf}(UXU^\dagger)=\sigma_2\bar UX^{\mathrm T}U^{\mathrm T}\sigma_2=U\,\sigma_2X^{\mathrm T}\sigma_2\,U^\dagger=U\,\mathrm{sf}(X)\,U^\dagger
\]

%dsp-level:0
\[
\quad [\sigma_2U^{\mathrm T}\sigma_2=U^\dagger\Rightarrow\sigma_2\bar U\sigma_2=U]
\]

%dsp-level:0
\[
\implies\mathrm{sf}\otimes\mathrm{sf}\ \text{commutes with}\ \mathrm{Ad}_{U\otimes U};
\]

%dsp-level:0
\[
\qquad \text{(5)}\ (\Rightarrow):\ \ U_c:=\tfrac12\bigl(\mathbb1-i(\sigma_1+\sigma_2+\sigma_3)\bigr):\ \ \mathrm{Ad}_{U_c}:\ \sigma_1\mapsto\sigma_2\mapsto\sigma_3\mapsto\sigma_1
\]

%dsp-level:0
\[
T\,\mathrm{Ad}_{U_c\otimes U_c}(\sigma_0\otimes\sigma_i)=\eta^0\eta^{i+1}\,\sigma_0\otimes\sigma_{i+1};
\]

%dsp-level:0
\[
\qquad \mathrm{Ad}_{U_c\otimes U_c}\,T(\sigma_0\otimes\sigma_i)=\eta^0\eta^{i}\,\sigma_0\otimes\sigma_{i+1}\quad(i\bmod3)
\]

%dsp-level:0
\[
\implies\eta^1=\eta^2=\eta^3\ \ [\eta^0=\pm1];\qquad\text{exactly one odd entry}\implies\mathfrak o=0.\ \blacksquare
\]

\textbf{Phys.} A3 now says only that the sign pattern is $1+3$. The overall sign is the mostly-plus / mostly-minus convention; which spatial slot is odd is a relabelling. What remains is one selection: the odd slot is the identity slot, which is the unique choice making $T$ rotation-covariant, namely the spin flip $\mathrm{sf}\otimes\mathrm{sf}$ (time reversal on the pair). Any other slot gives a reflection that singles out a spatial direction.

\textbf{Rmk (what $\mathfrak o\ge1$ would do).} The sector dimensions $(7,3,3,3)$ do not change, but CT1.2's table does. For $\mathfrak o=1$: $W_2,W_3\in\mathcal E_{++}$ (not $\mathcal E_{+-}$); $J_2,J_3\in\mathcal E_{--}$ and $K_2,K_3\in\mathcal E_{-+}$ (swapped); $E^{12},E^{13}\in\mathcal E_{+-}$ (not $\mathcal E_{++}$). Only $W_1,J_1,K_1$ sit where CT1.2 puts them. Symmetric statements hold for $\mathfrak o=2,3$. (Numerical check over all four $\mathfrak o$ on the 16-element basis; machine: \texttt{checks/\allowbreak{}verify\_\allowbreak{}CT10\_\allowbreak{}signature.py}.) So the selection $\mathfrak o=0$ in C13 is load-bearing for CT1.2 and everything built on the $W_i$ triplet, not cosmetic.

\end{theorem}

\begin{theorem}[{\texorpdfstring{CT1.1 — Joint $(S,T)$ decomposition of $\mathfrak H(\mathbb C^4)$}{CT1.1 - Joint decomposition of}}]\label{tag:CT1-1}
\textbf{Thm.} $\ [S,T]=0;$ $\qquad\mathfrak H(\mathbb C^4)=\mathcal E_{++}\oplus\mathcal E_{+-}\oplus\mathcal E_{-+}\oplus\mathcal E_{--};$ $\qquad(\dim\mathcal E_{++},\dim\mathcal E_{+-},\dim\mathcal E_{-+},\dim\mathcal E_{--})=(7,3,3,3);\qquad7+3+3+3=16.$

\textbf{Req.} A3, C13; C3; CT1.0; Pauli basis facts: $\ \mathfrak T(\sigma_\mu\sigma_\alpha)=2\delta_{\mu\alpha}$, $\sigma_\mu^\dagger=\sigma_\mu$ [linear algebra].

\textbf{Pf.}

%dsp-level:0
\[
\bigl\{\sigma_\mu\otimes\sigma_\nu\bigr\}_{\mu,\nu=0}^{3}\subset\mathfrak H(\mathbb C^4);\qquad
\mathfrak T\bigl((\sigma_\mu\otimes\sigma_\nu)(\sigma_\alpha\otimes\sigma_\beta)\bigr)=4\,\delta_{\mu\alpha}\delta_{\nu\beta}
\]

%dsp-level:0
\[
\Longrightarrow\ \text{orthogonal}\implies\text{basis};\qquad\dim_{\mathbb R}\mathfrak H(\mathbb C^4)=16
\]

%dsp-level:0
\[
\mathrm{A3}:\ \ S:\ \sigma_\mu\otimes\sigma_\nu\mapsto\sigma_\nu\otimes\sigma_\mu\ (\text{permutation});
\]

%dsp-level:0
\[
\qquad T:\ \sigma_\mu\otimes\sigma_\nu\mapsto\eta^\mu\eta^\nu\,\sigma_\mu\otimes\sigma_\nu,\ \ \eta^\mu\eta^\nu\in\{\pm1\}
\]

%dsp-level:0
\[
\implies S^2=T^2=\mathbb1,\ \ \operatorname{spec}\subseteq\{\pm1\};
\]

%dsp-level:0
\[
\qquad ST(\sigma_\mu\otimes\sigma_\nu)=\eta^\mu\eta^\nu\,\sigma_\nu\otimes\sigma_\mu=TS(\sigma_\mu\otimes\sigma_\nu)\implies[S,T]=0
\]

%dsp-level:0
\[
[S,T]=0\ \wedge\ S^2=T^2=\mathbb1\implies\text{simultaneous eigenspaces}\ \mathcal E_{st}\ \text{span}\ \mathfrak H(\mathbb C^4)\quad[\mathrm{C3}]
\]

%dsp-level:0
\[
\eta:=(-,+,+,+)\ [\mathrm{C13}]:\quad T=+1\iff(\mu,\nu)\in\{0,0\}\cup\{i,j\};\qquad T=-1\iff(\mu,\nu)\in\{0,i\}\cup\{i,0\}
\]

%dsp-level:0
\[
\mathcal E_{++}:\ \{\sigma_\mu\otimes\sigma_\mu\}_{\mu=0}^{3}\ (4)\ \cup\ \{\sigma_i\otimes\sigma_j+\sigma_j\otimes\sigma_i\}_{i<j}\ (3)=7
\]

%dsp-level:0
\[
\mathcal E_{+-}:\ \{\sigma_0\otimes\sigma_i+\sigma_i\otimes\sigma_0\}_{i=1}^{3}=3;\qquad
\mathcal E_{-+}:\ \{\sigma_i\otimes\sigma_j-\sigma_j\otimes\sigma_i\}_{i<j}=3
\]

%dsp-level:0
\[
\mathcal E_{--}:\ \{\sigma_0\otimes\sigma_i-\sigma_i\otimes\sigma_0\}_{i=1}^{3}=3;\qquad
7+3+3+3=16\implies\mathfrak H(\mathbb C^4)=\bigoplus_{s,t}\mathcal E_{st}.\ \blacksquare
\]

\textbf{Phys.} The 16-dimensional space of two-qubit observables splits under exchange symmetry ($S$) and the sign grading $T$ ($\eta^\mu\eta^\nu$; metric-sign reading by C13) into blocks $7+3+3+3$. This grading is the skeleton of the whole force structure: the $7$-block supplies gravity, Higgs and uniform dark; the $3$-blocks supply weak, Janus-seal, and the remaining virtual sector.

\end{theorem}

\begin{theorem}[{\texorpdfstring{CT1.2 — Generator basis and sector assignment}{CT1.2 - Generator basis and sector assignment}}]\label{tag:CT1-2}
\textbf{Thm.} Define

%dsp-level:0
\[
E_0:=\sigma_0\otimes\sigma_0;\qquad
\Sigma:=\textstyle\sum_{i=1}^{3}\sigma_i\otimes\sigma_i;\qquad
W_i:=\tfrac12\bigl(\sigma_0\otimes\sigma_i+\sigma_i\otimes\sigma_0\bigr);
\]

%dsp-level:0
\[
J_i:=\tfrac12\bigl(\sigma_j\otimes\sigma_k-\sigma_k\otimes\sigma_j\bigr)\ [(i,j,k)\ \text{cyclic}];
\]

%dsp-level:0
\[
\qquad K_i:=\tfrac12\bigl(\sigma_0\otimes\sigma_i-\sigma_i\otimes\sigma_0\bigr);
\]

%dsp-level:0
\[
\qquad E^{ij}:=\tfrac12\bigl(\sigma_i\otimes\sigma_j+\sigma_j\otimes\sigma_i\bigr).
\]

Then

%dsp-level:0
\[
\begin{array}{c|c|c|c|c}
\text{generator} & \# & S & T & \text{sector}\\ \hline
E_0 & 1 & + & + & \mathcal E_{++}\\
\Sigma & 1 & + & + & \mathcal E_{++}\\
\hat E^{ij}\ (\text{traceless}) & 5 & + & + & \mathcal E_{++}\\
W_i & 3 & + & - & \mathcal E_{+-}\\
J_i & 3 & - & + & \mathcal E_{-+}\\
K_i & 3 & - & - & \mathcal E_{--}
\end{array}
\qquad 1+1+5+3+3+3=16\implies\text{basis of}\ \mathfrak H(\mathbb C^4);
\]

%dsp-level:0
\[
W_i=\tfrac12\bigl(\sigma_i^{(p)}+\sigma_i^{(q)}\bigr)=\text{total-spin generator}.
\]

\textbf{Req.} A3, C1, C3, C13; CT1.0, CT1.1.

\textbf{Pf.}

%dsp-level:0
\[
S E_0=E_0;\ \ \eta^0\eta^0=+1\implies E_0\in\mathcal E_{++}\quad[\mathrm{A3}]
\]

%dsp-level:0
\[
S E^{ij}=E^{ij}\ (\text{by construction});\ \ \eta^i\eta^j=+1\ [\mathrm{C13}]\implies E^{ij}\in\mathcal E_{++};\qquad\#\{E^{ij}\}_{i\le j}=6
\]

%dsp-level:0
\[
\delta_{ij}E^{ij}=\textstyle\sum_iE^{ii}=\Sigma\ (1\ \text{component})\ \oplus\ \hat E^{ij}\ (5);\qquad S,T\ \text{diagonal on index pairs}\implies\text{both}\subseteq\mathcal E_{++}
\]

%dsp-level:0
\[
W_i:\ S=+;\ \ \eta^0\eta^i=-1\ [\mathrm{C13}]\implies T=-\implies W_i\in\mathcal E_{+-}
\]

%dsp-level:0
\[
J_i:\ S=-;\ \ \eta^j\eta^k=+1\ [\mathrm{C13}]\implies T=+\implies J_i\in\mathcal E_{-+}
\]

%dsp-level:0
\[
K_i:\ S=-;\ \ \eta^0\eta^i=-1\ [\mathrm{C13}]\implies T=-\implies K_i\in\mathcal E_{--}
\]

%dsp-level:0
\[
\text{counts}\ (1,1,5,3,3,3)\ \text{match}\ (7{=}1{+}1{+}5,\ 3,\ 3,\ 3)\ [\mathrm{CT1.1}];
\]

%dsp-level:0
\[
\qquad \text{distinct orthogonal combinations}\implies\text{independent}\implies\text{basis}
\]

%dsp-level:0
\[
\sigma_0\otimes\sigma_i\equiv\sigma_i^{(p)},\ \ \sigma_i\otimes\sigma_0\equiv\sigma_i^{(q)}\implies W_i=\tfrac12\bigl(\sigma_i^{(p)}+\sigma_i^{(q)}\bigr).\ \blacksquare
\]

\textbf{Phys.} This fixes the named basis used throughout the corpus: identity $E_0$, exchange scalar $\Sigma$, total spin $W_i$ (weak sector), the antisymmetric spatial and temporal generators $J_i,K_i$ (Janus/seal sectors), and the symmetric spatial tensors $E^{ij}$ (gravity). Each generator’s $(S,T)$ parity decides which force aspect of A8 it belongs to.

\end{theorem}

\begin{theorem}[{\texorpdfstring{CT1.3 — Centrality of the identity}{CT1.3 - Centrality of the identity}}]\label{tag:CT1-3}
\textbf{Thm.} $\ \forall X\in\mathfrak H(\mathbb C^4):\quad[E_0,X]=0.$

\textbf{Req.} C1; CT1.2.

\textbf{Pf.}

%dsp-level:0
\[
E_0=\sigma_0\otimes\sigma_0=I_2\otimes I_2=I_4\quad[\mathrm{C1},\ \mathrm{CT1.2}];\qquad
[I_4,X]=I_4X-XI_4=0.\ \blacksquare
\]

\textbf{Phys.} The identity direction commutes with everything, so the aspect it carries — the uniform dark term $F_{D^+}=\Lambda\mathbb I$ — acts identically on every configuration: a cosmological-constant-type force.

\end{theorem}

\begin{theorem}[{\texorpdfstring{CT1.4 — The $W_i$ close $\mathfrak{su}(2)$}{CT1.4 - The close}}]\label{tag:CT1-4}
\textbf{Thm.} $\ [W_i,W_j]=i\,\epsilon_{ijk}W_k;\qquad\operatorname{span}_{\mathbb R}\{W_i\}\cong\mathfrak{su}(2)\ [\mathrm{C2}].$

\textbf{Req.} C2; CT1.2; $\ [\sigma_i,\sigma_j]=2i\,\epsilon_{ijk}\sigma_k$ [Pauli algebra].

\textbf{Pf.}

%dsp-level:0
\[
\text{definitions cyclically symmetric}\implies\text{suffices}:\ [W_1,W_2]=iW_3
\]

%dsp-level:0
\[
[W_1,W_2]=\tfrac14\bigl[\sigma_0\otimes\sigma_1+\sigma_1\otimes\sigma_0,\ \ \sigma_0\otimes\sigma_2+\sigma_2\otimes\sigma_0\bigr]\quad[\mathrm{CT1.2}]
\]

%dsp-level:0
\[
[\sigma_0\otimes\sigma_1,\ \sigma_0\otimes\sigma_2]=\sigma_0^2\otimes[\sigma_1,\sigma_2]=2i\,(\sigma_0\otimes\sigma_3)
\]

%dsp-level:0
\[
[\sigma_0\otimes\sigma_1,\ \sigma_2\otimes\sigma_0]=\sigma_2\otimes\sigma_1-\sigma_2\otimes\sigma_1=0;\qquad
[\sigma_1\otimes\sigma_0,\ \sigma_0\otimes\sigma_2]=0\quad[\text{distinct factors}]
\]

%dsp-level:0
\[
[\sigma_1\otimes\sigma_0,\ \sigma_2\otimes\sigma_0]=[\sigma_1,\sigma_2]\otimes\sigma_0^2=2i\,(\sigma_3\otimes\sigma_0)
\]

%dsp-level:0
\[
[W_1,W_2]=\tfrac14\bigl(2i\,\sigma_0\otimes\sigma_3+2i\,\sigma_3\otimes\sigma_0\bigr)=\tfrac i2\bigl(\sigma_0\otimes\sigma_3+\sigma_3\otimes\sigma_0\bigr)=iW_3.\ \blacksquare
\]

\textbf{Phys.} The three symmetric mixed generators close a perfect $\mathfrak{su}(2)$: they are the total spin of the qubit pair. This $\mathfrak{su}(2)$ becomes the weak-isospin gauge algebra of the corpus — the weak force is the pair’s rotation algebra, not an added postulate.

\end{theorem}

\begin{theorem}[{\texorpdfstring{CT1.5 — Commutators of $\Sigma$ with $J_i$, $K_i$, $W_i$}{CT1.5 - Commutators of with , ,}}]\label{tag:CT1-5}
\textbf{Thm.} $\ [\Sigma,J_i]=4i\,K_i;\qquad[\Sigma,K_i]=-4i\,J_i;\qquad[\Sigma,W_i]=0.$

\textbf{Req.} CT0.7(1),(3) $\ [\Sigma=2\mathbb S-E_0;\ \ S\text{-even}\Rightarrow[\mathbb S,\cdot]=0]$, CT1.2, CT1.3; Pauli algebra.

\textbf{Pf.}

%dsp-level:0
\[
\text{cyclicity}\implies\text{suffices}:\ i=1
\]

%dsp-level:0
\[
[\sigma_1\otimes\sigma_1,\ J_1]=0:
\]

%dsp-level:0
\[
\quad [\sigma_1\otimes\sigma_1,\ \sigma_2\otimes\sigma_3]=(i\sigma_3)\otimes(-i\sigma_2)-(-i\sigma_3)\otimes(i\sigma_2)=0;\quad\text{likewise}\ \sigma_3\otimes\sigma_2
\]

%dsp-level:0
\[
[\sigma_2\otimes\sigma_2,\ \sigma_2\otimes\sigma_3]=\sigma_0\otimes[\sigma_2,\sigma_3]=2i\,(\sigma_0\otimes\sigma_1);\qquad
[\sigma_2\otimes\sigma_2,\ \sigma_3\otimes\sigma_2]=2i\,(\sigma_1\otimes\sigma_0)
\]

%dsp-level:0
\[
\implies[\sigma_2\otimes\sigma_2,\ J_1]=\tfrac12\bigl(2i\,\sigma_0\otimes\sigma_1-2i\,\sigma_1\otimes\sigma_0\bigr)=2i\,K_1
\]

%dsp-level:0
\[
[\sigma_3\otimes\sigma_3,\ \sigma_2\otimes\sigma_3]=-2i\,(\sigma_1\otimes\sigma_0);
\]

%dsp-level:0
\[
\qquad [\sigma_3\otimes\sigma_3,\ \sigma_3\otimes\sigma_2]=-2i\,(\sigma_0\otimes\sigma_1)\implies[\sigma_3\otimes\sigma_3,\ J_1]=2i\,K_1
\]

%dsp-level:0
\[
[\Sigma,J_1]=0+2iK_1+2iK_1=4i\,K_1
\]

%dsp-level:0
\[
[\sigma_1\otimes\sigma_1,\ K_1]=0\ (\text{as above});
\]

%dsp-level:0
\[
\qquad [\sigma_2\otimes\sigma_2,\ \sigma_0\otimes\sigma_1]=-2i\,(\sigma_2\otimes\sigma_3);
\]

%dsp-level:0
\[
\qquad [\sigma_2\otimes\sigma_2,\ \sigma_1\otimes\sigma_0]=-2i\,(\sigma_3\otimes\sigma_2)
\]

%dsp-level:0
\[
\implies[\sigma_2\otimes\sigma_2,\ K_1]=\tfrac12\bigl(-2i\,\sigma_2\otimes\sigma_3+2i\,\sigma_3\otimes\sigma_2\bigr)=-2i\,J_1
\]

%dsp-level:0
\[
[\sigma_3\otimes\sigma_3,\ \sigma_0\otimes\sigma_1]=2i\,(\sigma_3\otimes\sigma_2);
\]

%dsp-level:0
\[
\qquad [\sigma_3\otimes\sigma_3,\ \sigma_1\otimes\sigma_0]=2i\,(\sigma_2\otimes\sigma_3)\implies[\sigma_3\otimes\sigma_3,\ K_1]=-2i\,J_1
\]

%dsp-level:0
\[
[\Sigma,K_1]=-4i\,J_1
\]

%dsp-level:0
\[
\Sigma=2\mathbb S-E_0\ [\mathrm{CT0.7}(1)];\qquad
SW_i=W_i\ (S\text{-even})\implies[\mathbb S,W_i]=0\ [\mathrm{CT0.7}(3)];\qquad
[E_0,W_i]=0\ [\mathrm{CT1.3}]
\]

%dsp-level:0
\[
\implies[\Sigma,W_i]=2[\mathbb S,W_i]-[E_0,W_i]=0.\ \blacksquare
\]

\textbf{Phys.} The exchange scalar $\Sigma$ rotates the two seal sectors into one another ($J\leftrightarrow K$, a $90^\circ$ rotation inside the virtual sector) while leaving total spin untouched. The volume/trace mode therefore couples to Janus (edge creation–annihilation) dynamics but not to spin.

\end{theorem}

\begin{theorem}[{\texorpdfstring{CT1.6 — Commutators $[J_i,J_j]$ and $[K_i,K_j]$}{CT1.6 - Commutators and}}]\label{tag:CT1-6}
\textbf{Thm.} $\ [J_i,J_j]=i\,\epsilon_{ijk}W_k;\qquad[K_i,K_j]=i\,\epsilon_{ijk}W_k.$

\textbf{Req.} CT1.2; Pauli algebra.

\textbf{Pf.}

%dsp-level:0
\[
\text{suffices}:\ [J_1,J_2]=iW_3=[K_1,K_2]
\]

%dsp-level:0
\[
[J_1,J_2]=\tfrac14\bigl[\sigma_2\otimes\sigma_3-\sigma_3\otimes\sigma_2,\ \ \sigma_3\otimes\sigma_1-\sigma_1\otimes\sigma_3\bigr]
\]

%dsp-level:0
\[
[\sigma_2\otimes\sigma_3,\ \sigma_3\otimes\sigma_1]=(i\sigma_1)\otimes(i\sigma_2)-(-i\sigma_1)\otimes(-i\sigma_2)=0
\]

%dsp-level:0
\[
[\sigma_2\otimes\sigma_3,\ -\sigma_1\otimes\sigma_3]=-[\sigma_2,\sigma_1]\otimes\sigma_3^2=2i\,(\sigma_3\otimes\sigma_0)
\]

%dsp-level:0
\[
[-\sigma_3\otimes\sigma_2,\ \sigma_3\otimes\sigma_1]=-\sigma_0\otimes[\sigma_2,\sigma_1]=2i\,(\sigma_0\otimes\sigma_3)
\]

%dsp-level:0
\[
[-\sigma_3\otimes\sigma_2,\ -\sigma_1\otimes\sigma_3]=(i\sigma_2)\otimes(i\sigma_1)-(-i\sigma_2)\otimes(-i\sigma_1)=0
\]

%dsp-level:0
\[
\implies[J_1,J_2]=\tfrac14\bigl(2i\,\sigma_3\otimes\sigma_0+2i\,\sigma_0\otimes\sigma_3\bigr)=iW_3
\]

%dsp-level:0
\[
[K_1,K_2]:\ \ \text{cross terms}=0\ [\text{distinct factors}];
\]

%dsp-level:0
\[
\qquad [K_1,K_2]=\tfrac14\Bigl([\sigma_0\otimes\sigma_1,\sigma_0\otimes\sigma_2]+[\sigma_1\otimes\sigma_0,\sigma_2\otimes\sigma_0]\Bigr)
\]

%dsp-level:0
\[
=\tfrac14\bigl(2i\,\sigma_0\otimes\sigma_3+2i\,\sigma_3\otimes\sigma_0\bigr)=iW_3.\ \blacksquare
\]

\textbf{Phys.} Neither seal sector closes on itself: brackets of two $J$’s or two $K$’s land in the total-spin (weak) sector. Reality-flipping operations source the rotational/weak structure of the pair — the first algebraic link between the Janus dynamics and gauge structure.

\end{theorem}

\begin{theorem}[{\texorpdfstring{CT1.7 — Mixed commutators $[J_i,K_j]$}{CT1.7 - Mixed commutators}}]\label{tag:CT1-7}
\textbf{Thm.}

%dsp-level:0
\[
[J_i,K_j]=i\bigl(\delta_{ij}\Sigma-E^{ij}\bigr)=\tfrac i2\Bigl(2\delta_{ij}\Sigma-\sigma_i\otimes\sigma_j-\sigma_j\otimes\sigma_i\Bigr);
\]

%dsp-level:0
\[
[J_i,K_i]=i\!\!\sum_{k\ne i}\!\!\sigma_k\otimes\sigma_k;\qquad
[J_i,K_j]=-iE^{ij}=-\tfrac i2\bigl(\sigma_i\otimes\sigma_j+\sigma_j\otimes\sigma_i\bigr)\ \ (i\ne j);
\]

%dsp-level:0
\[
[J_i,K_j]\in(\mathcal E_{++})_{\mathrm{sym\ spatial}};\qquad [J_i,K_j]\notin\operatorname{span}\{W_k\}.
\]

\textbf{Req.} CT1.1, CT1.2; $\epsilon$-identities: $\ \epsilon_{iab}\epsilon_{bjm}=\delta_{ij}\delta_{am}-\delta_{im}\delta_{aj}$, $\ \epsilon_{iab}\epsilon_{ajm}=\delta_{bj}\delta_{im}-\delta_{bm}\delta_{ij}$.

\textbf{Pf.}

%dsp-level:0
\[
J_i=\tfrac12\,\epsilon_{iab}\,\sigma_a\otimes\sigma_b\ (a,b\in\{1,2,3\})\quad[\mathrm{CT1.2}];\qquad
K_j=\tfrac12\bigl(\sigma_0\otimes\sigma_j-\sigma_j\otimes\sigma_0\bigr)
\]

%dsp-level:0
\[
[\sigma_a\otimes\sigma_b,\ \sigma_0\otimes\sigma_j]=\sigma_a\otimes[\sigma_b,\sigma_j]=2i\,\epsilon_{bjm}\,\sigma_a\otimes\sigma_m
\]

%dsp-level:0
\[
[\sigma_a\otimes\sigma_b,\ \sigma_j\otimes\sigma_0]=[\sigma_a,\sigma_j]\otimes\sigma_b=2i\,\epsilon_{ajm}\,\sigma_m\otimes\sigma_b
\]

%dsp-level:0
\[
\implies[J_i,K_j]=\tfrac i2\,\epsilon_{iab}\bigl(\epsilon_{bjm}\,\sigma_a\otimes\sigma_m-\epsilon_{ajm}\,\sigma_m\otimes\sigma_b\bigr)
\]

%dsp-level:0
\[
\epsilon_{iab}\epsilon_{bjm}\,\sigma_a\otimes\sigma_m=\delta_{ij}\Sigma-\sigma_j\otimes\sigma_i;\qquad
\epsilon_{iab}\epsilon_{ajm}\,\sigma_m\otimes\sigma_b=\sigma_i\otimes\sigma_j-\delta_{ij}\Sigma
\]

%dsp-level:0
\[
[J_i,K_j]
\]

%dsp-level:0
\[
= \tfrac i2\Bigl(\delta_{ij}\Sigma-\sigma_j\otimes\sigma_i-\sigma_i\otimes\sigma_j+\delta_{ij}\Sigma\Bigr)
\]

%dsp-level:0
\[
= \tfrac i2\bigl(2\delta_{ij}\Sigma-\sigma_i\otimes\sigma_j-\sigma_j\otimes\sigma_i\bigr)
\]

%dsp-level:0
\[
= i\bigl(\delta_{ij}\Sigma-E^{ij}\bigr)
\]

%dsp-level:0
\[
i=j:\ \ i\bigl(\Sigma-\sigma_i\otimes\sigma_i\bigr)=i\!\!\sum_{k\ne i}\!\!\sigma_k\otimes\sigma_k;\qquad
i\ne j:\ \ \delta_{ij}=0\implies-iE^{ij}
\]

%dsp-level:0
\[
\Sigma,\ E^{ij}\in\mathcal E_{++};\ \ W_k\in\mathcal E_{+-};\ \ \text{direct sum}\ [\mathrm{CT1.1},\ \mathrm{CT1.2}]\implies[J_i,K_j]\notin\operatorname{span}\{W_k\}.\ \blacksquare
\]

\textbf{Phys.} Bracketing a spatial seal with a temporal seal produces symmetric spatial — metric-type — data, never weak data. Mixed virtual processes feed the gravitational sector: the algebraic seed of “quantum fluctuations sourcing geometry.”

\end{theorem}

\begin{theorem}[{\texorpdfstring{CT1.8 — Action of $W_i$ on $J_j$ and $K_j$}{CT1.8 - Action of on and}}]\label{tag:CT1-8}
\textbf{Thm.} $\ [W_i,J_j]=i\,\epsilon_{ijk}J_k;\qquad[W_i,K_j]=i\,\epsilon_{ijk}K_k;$ i.e. $\{J_i\}$, $\{K_i\}$: $SO(3)$ vectors under $\operatorname{span}\{W_i\}$.

\textbf{Req.} CT1.2; Pauli algebra.

\textbf{Pf.}

%dsp-level:0
\[
\text{suffices}:\ [W_3,J_1]=iJ_2,\quad[W_3,K_1]=iK_2
\]

%dsp-level:0
\[
[W_3,J_1]=\tfrac14\bigl[\sigma_0\otimes\sigma_3+\sigma_3\otimes\sigma_0,\ \ \sigma_2\otimes\sigma_3-\sigma_3\otimes\sigma_2\bigr]
\]

%dsp-level:0
\[
[\sigma_0\otimes\sigma_3,\ \sigma_2\otimes\sigma_3]=\sigma_2\otimes[\sigma_3,\sigma_3]=0;
\]

%dsp-level:0
\[
\qquad [\sigma_0\otimes\sigma_3,\ -\sigma_3\otimes\sigma_2]=-\sigma_3\otimes[\sigma_3,\sigma_2]=2i\,(\sigma_3\otimes\sigma_1)
\]

%dsp-level:0
\[
[\sigma_3\otimes\sigma_0,\ \sigma_2\otimes\sigma_3]=[\sigma_3,\sigma_2]\otimes\sigma_3=-2i\,(\sigma_1\otimes\sigma_3);\qquad
[\sigma_3\otimes\sigma_0,\ -\sigma_3\otimes\sigma_2]=0
\]

%dsp-level:0
\[
\implies[W_3,J_1]=\tfrac i2\bigl(\sigma_3\otimes\sigma_1-\sigma_1\otimes\sigma_3\bigr)=iJ_2
\]

%dsp-level:0
\[
[W_3,K_1]=\tfrac14\bigl[\sigma_0\otimes\sigma_3+\sigma_3\otimes\sigma_0,\ \ \sigma_0\otimes\sigma_1-\sigma_1\otimes\sigma_0\bigr]
\]

%dsp-level:0
\[
[\sigma_0\otimes\sigma_3,\ \sigma_0\otimes\sigma_1]=2i\,(\sigma_0\otimes\sigma_2);\qquad
[\sigma_0\otimes\sigma_3,\ -\sigma_1\otimes\sigma_0]=0=[\sigma_3\otimes\sigma_0,\ \sigma_0\otimes\sigma_1]
\]

%dsp-level:0
\[
[\sigma_3\otimes\sigma_0,\ -\sigma_1\otimes\sigma_0]=-[\sigma_3,\sigma_1]\otimes\sigma_0=-2i\,(\sigma_2\otimes\sigma_0)
\]

%dsp-level:0
\[
\implies[W_3,K_1]=\tfrac14\bigl(2i\,\sigma_0\otimes\sigma_2-2i\,\sigma_2\otimes\sigma_0\bigr)=iK_2;\qquad
\text{cyclicity}\implies\text{general form}.\ \blacksquare
\]

\textbf{Phys.} Total spin rotates both virtual sectors as vectors: the seal generators carry spin $1$ under the weak $\mathfrak{su}(2)$. The reality/virtuality bookkeeping is rotation-covariant — angular momentum is built into the ledger before any dynamics is introduced.

\end{theorem}

\begin{theorem}[{\texorpdfstring{CT1.9 — Sector closure table and the global algebra}{CT1.9 - Sector closure table and the global algebra}}]\label{tag:CT1-9}
\textbf{Thm.}

%dsp-level:0
\[
\begin{array}{c|c}
\text{bracket} & \text{closes in}\\ \hline
[W_i,W_j] & \operatorname{span}\{W_k\}=\mathcal E_{+-}\\
{}[W_i,J_j] & \operatorname{span}\{J_k\}=\mathcal E_{-+}\\
{}[W_i,K_j] & \operatorname{span}\{K_k\}=\mathcal E_{--}\\
{}[J_i,J_j] & \operatorname{span}\{W_k\}\\
{}[K_i,K_j] & \operatorname{span}\{W_k\}\\
{}[J_i,K_j] & (\mathcal E_{++})_{\mathrm{sym\ spatial}}\\
{}[\Sigma,J_i] & \operatorname{span}\{K_i\}\\
{}[\Sigma,K_i] & \operatorname{span}\{J_i\}\\
{}[\Sigma,W_i] & 0\\
{}[E_0,X] & 0\ \ \forall X
\end{array}
\]

%dsp-level:0
\[
\bigl(\mathfrak H(\mathbb C^4),\ (X,Y)\mapsto-i[X,Y]\bigr)\cong\mathfrak u(4)\ [\mathrm{C2}];\qquad
Z\bigl(\mathfrak u(4)\bigr)=\operatorname{span}_{\mathbb R}\{E_0\};
\]

%dsp-level:0
\[
\mathfrak{su}(4)=\operatorname{span}_{\mathbb R}\bigl\{\Sigma,\ \hat E^{ij},\ W_i,\ J_i,\ K_i\bigr\};\qquad
\dim_{\mathbb R}\mathfrak{su}(4)=1+5+3+3+3=15 .
\]

\textbf{Req.} C2; CT1.1–CT1.8; $\ \mathfrak T(\sigma_\mu\otimes\sigma_\nu)=4\,\delta_{\mu0}\delta_{\nu0}$.

\textbf{Pf.}

%dsp-level:0
\[
\text{rows }1\text{–}5:\ [\mathrm{CT1.4},\ \mathrm{CT1.8},\ \mathrm{CT1.6}];\qquad
\text{row }6:\ [\mathrm{CT1.7}];\qquad
\text{rows }7\text{–}9:\ [\mathrm{CT1.5}];\qquad
\text{row }10:\ [\mathrm{CT1.3}]
\]

%dsp-level:0
\[
\text{sector labels}:\ [\mathrm{CT1.2}]
\]

%dsp-level:0
\[
X,Y\in\mathfrak H(\mathbb C^4)\implies(-i[X,Y])^\dagger=-i[X,Y];
\]

%dsp-level:0
\[
\qquad \text{bilinear}\ \wedge\ \text{antisymmetric}\ \wedge\ \text{Jacobi}\ [\text{associative commutator}]
\]

%dsp-level:0
\[
\dim_{\mathbb R}\mathfrak H(\mathbb C^4)=16\ [\mathrm{CT1.1}]\implies\bigl(\mathfrak H(\mathbb C^4),-i[\cdot,\cdot]\bigr)=\mathfrak u(4)\quad[\mathrm{C2}]
\]

%dsp-level:0
\[
X\in Z\iff[X,Y]=0\ \ \forall Y\in\mathfrak H(\mathbb C^4)\implies X\ \text{scalar}\implies Z=\operatorname{span}_{\mathbb R}\{E_0\}\quad[\mathrm{CT1.3}]
\]

%dsp-level:0
\[
\mathfrak T(\sigma_\mu\otimes\sigma_\nu)=4\,\delta_{\mu0}\delta_{\nu0}\implies E_0\ \text{alone has trace}\implies\text{traceless complement}=\mathfrak{su}(4),\ \dim=16-1=15
\]

%dsp-level:0
\[
\mathfrak{su}(4)=\operatorname{span}_{\mathbb R}\{\Sigma,\hat E^{ij},W_i,J_i,K_i\}:\quad1+5+3+3+3=15.\ \blacksquare
\]

\textbf{Phys.} The full commutator structure closes: the $16$ pair observables form $\mathfrak u(4)=\mathfrak u(1)\oplus\mathfrak{su}(4)$, with the central $\mathfrak u(1)$ ($E_0$) as the uniform dark aspect and the $15$ traceless directions carrying gravity ($5$), weak ($3$), the two seal sectors ($3{+}3$) and the exchange scalar ($1$). Every force algebra later in the corpus is a subalgebra of this one $16$-dimensional algebra.

\end{theorem}

\begin{theorem}[{\texorpdfstring{CT1.10 — $SO(3)$ decomposition of $\mathcal E_{++}$}{CT1.10 - decomposition of}}]\label{tag:CT1-10}
\textbf{Thm.} Under $\operatorname{span}\{W_i\}$:

%dsp-level:0
\[
\mathcal E_{++}\ \longrightarrow\ \mathbf 1\oplus\mathbf 5\oplus\mathbf 1\qquad(7=1+5+1);
\]

%dsp-level:0
\[
E_0:\ \mathbf 1;\qquad\Sigma:\ \mathbf 1;\qquad
\hat E^{ij}:=E^{ij}-\tfrac13\delta^{ij}\Sigma:\ \mathbf 5\ (\text{spin }2),\quad\delta_{ij}\hat E^{ij}=0 .
\]

\textbf{Req.} CT1.1, CT1.3, CT1.5, CT1.8; symmetric traceless rank-$2$ tensor $=\ell=2$ representation [$SO(3)$ representation theory].

\textbf{Pf.}

%dsp-level:0
\[
\dim\mathcal E_{++}=7\quad[\mathrm{CT1.1}]
\]

%dsp-level:0
\[
[E_0,W_i]=0\ [\mathrm{CT1.3}];\ \ E_0\ \text{index-free}\implies E_0:\ \mathbf 1
\]

%dsp-level:0
\[
E^{ij}=\underbrace{\tfrac13\delta^{ij}\Sigma}_{\text{trace}}+\underbrace{\hat E^{ij}}_{\text{traceless}}\ \ (\text{unique split});\qquad
\delta_{ij}\hat E^{ij}=\Sigma-\tfrac13\cdot3\cdot\Sigma=0
\]

%dsp-level:0
\[
[\Sigma,W_i]=0\ [\mathrm{CT1.5}]\implies\Sigma:\ \mathbf 1
\]

%dsp-level:0
\[
[W_a,\sigma_i\otimes\sigma_j]=i\,\epsilon_{aik}\,\sigma_k\otimes\sigma_j+i\,\epsilon_{ajk}\,\sigma_i\otimes\sigma_k\quad[\mathrm{CT1.2},\ \text{Pauli algebra}]
\]

%dsp-level:0
\[
\implies[W_a,E^{ij}]=i\bigl(\epsilon_{aik}E^{kj}+\epsilon_{ajk}E^{ik}\bigr);
\]

%dsp-level:0
\[
\qquad [W_a,\hat E^{ij}]=[W_a,E^{ij}]-\tfrac13\delta^{ij}\underbrace{[W_a,\Sigma]}_{=0\ [\mathrm{CT1.5}]}=i\bigl(\epsilon_{aik}\hat E^{kj}+\epsilon_{ajk}\hat E^{ik}\bigr)
\]

%dsp-level:0
\[
\implies\hat E^{ij}:\ \text{traceless symmetric rank-}2\ \text{tensor}=\ell=2:\ \mathbf 5\quad[SO(3)\ \text{representation theory}]
\]

%dsp-level:0
\[
1+1+5=7=\dim\mathcal E_{++}\implies\mathcal E_{++}=\mathbf 1\oplus\mathbf 5\oplus\mathbf 1.\ \blacksquare
\]

\textbf{Phys.} Seen by rotations, the real sector contains exactly two scalars and one five-component spin-$2$ object. The spin-$2$ quintet is the graviton field of the corpus; with the three $W_i$ as shift, the $(7,3)$ real blocks supply the ten components of a Lorentzian metric (CT1.12): conformal ($1$) $+$ traceless metric ($5$) $+$ volume ($1$) $+$ shift ($3$).

\end{theorem}

\begin{theorem}[{\texorpdfstring{CT1.11 — The two scalars are distinct}{CT1.11 - The two scalars are distinct}}]\label{tag:CT1-11}
\textbf{Thm.} $\ E_0\in Z\bigl(\mathfrak H(\mathbb C^4)\bigr);\qquad\Sigma\notin Z\bigl(\mathfrak H(\mathbb C^4)\bigr);\qquad\forall c\in\mathbb R\setminus\{0\}:\ E_0\ne c\,\Sigma;$ hence the A8 assignment $E_0\mapsto F_{D^+}$, $\Sigma\mapsto F_H$ [CT0.12] is well-defined and non-degenerate.

\textbf{Req.} A8; CT0.12, CT1.3, CT1.5.

\textbf{Pf.}

%dsp-level:0
\[
[E_0,X]=0\ \ \forall X\implies E_0\in Z\quad[\mathrm{CT1.3}]
\]

%dsp-level:0
\[
[\Sigma,J_1]=4i\,K_1\ne0\ \ (K_1\ne0)\implies\Sigma\notin Z\quad[\mathrm{CT1.5}]
\]

%dsp-level:0
\[
c\ne0:\ \ X\in Z\iff cX\in Z\implies\bigl[E_0\in Z\ \wedge\ \Sigma\notin Z\bigr]\implies E_0\ne c\,\Sigma\ \ \forall c\ne0
\]

%dsp-level:0
\[
\implies E_0,\ \Sigma:\ \text{independent directions of}\ \mathcal E_{++};\qquad
E_0\mapsto F_{D^+}=\Lambda\,\mathbb I,\ \ \Sigma\mapsto F_H\quad[\mathrm{A8},\ \mathrm{CT0.12}].\ \blacksquare
\]

\textbf{Phys.} The two rotation scalars are provably different objects: $E_0$ is central — it cannot distinguish configurations — and supplies the uniform dark aspect $\Lambda\mathbb I$; $\Sigma$ is non-central — it couples to the seals and to geometry — and supplies the scalar (Higgs-type) aspect. The corpus’s two scalar directions are established as distinct before either receives a physical name.

\end{theorem}

\begin{theorem}[{\texorpdfstring{CT1.11a — The involution table of the four sectors}{CT1.11a - The involution table of the four sectors}}]\label{tag:CT1-11a}
\textbf{Thm.} $(0)$ \textbf{Dictionary.} A8's four sectors are CT1.1's blocks under the identification

%dsp-level:0
\[
\mathbf F_{(+,\mathcal R)}=F_G\oplus F_H\oplus F_{D^+}\ (5{+}1{+}1)\ \leftrightarrow\ \mathcal E_{++};
\]

%dsp-level:0
\[
\qquad \mathbf F_{(-,\mathcal R)}=F_{W^+}\oplus F_{W^-}\oplus F_Z\ (3)\ \leftrightarrow\ \mathcal E_{+-};
\]

%dsp-level:0
\[
\mathbf F_{(+,\mathcal V)}=F_{J^+}\ (3)\ \leftrightarrow\ \mathcal E_{-+};\qquad
\mathbf F_{(-,\mathcal V)}=F_{J^-}\ (3)\ \leftrightarrow\ \mathcal E_{--}
\]

— first label $=$ the $T$-grading ($\eta^\mu\eta^\nu$ [A3, C13]), $\mathcal R/\mathcal V=$ the $S$-grading (reality is the swap grading [CT0.7]); machine-verified on CT1.2's generator table (17 generators: $E_0,\Sigma,\hat E^{ij}{\times}5,W_i,J_i,K_i$). $(1)$ \textbf{Status of the two gradings.} $S=\operatorname{Ad}(\mathbb S)$ with $\mathbb S=\tfrac12(E_0+\Sigma)$ [CT0.7(1)] — an \textit{inner algebra automorphism} ($\mathbb S^2=\mathbb S$? no: $\mathbb S$ unitary involution, machine). $T$ is a \textit{pure grading}: the diagonal $\pm1$-map $\sigma_\mu\otimes\sigma_\nu\mapsto\eta^\mu\eta^\nu\,\sigma_\mu\otimes\sigma_\nu$ [A3] — $\neg$ multiplicative (witness: $T(\sigma_0\otimes\sigma_1)\,T(\sigma_0\otimes\sigma_2)=(-1)(-1)\ [\mathrm{C13}]\ \sigma_0\otimes(\sigma_1\sigma_2)\ne T(\sigma_0\otimes\sigma_1\sigma_2)$-type failure, machine), hence $\neg$ inner. $[S,T]=0$ on all 16 basis elements [CT1.1]; the grading group on $\mathfrak H(\mathbb C^4)$ is $\mathbb Z_2\times\mathbb Z_2$ with simultaneous eigenspaces $\mathcal E_{st}$, $\dim=(7,3,3,3)$. $(2)$ \textbf{Table of the declared reversals.} With $\omega$ [D2.6] and $\Theta_{\mathfrak L}$ [D2.7]:

%dsp-level:0
\[
\begin{array}{c|cccc|c}
& \mathbf F_{(+,\mathcal R)} & \mathbf F_{(-,\mathcal R)} & \mathbf F_{(+,\mathcal V)} & \mathbf F_{(-,\mathcal V)} & \text{edge/chain data}\\ \hline
S=\operatorname{Ad}(\mathbb S) & + & + & - & - & \text{reality grading [CT0.7]}\\
T\ (\text{grading}) & + & - & + & - & \eta^\mu\eta^\nu\ [\mathrm{A3},\mathrm{C13}]\\
\omega & \text{fixed} & W^+\!\leftrightarrow\!W^-;\ Z\ \text{fixed} & \neg\text{declared} & \neg\text{declared} & s(c)\leftrightarrow t(c);\ \text{charge}\mapsto-\text{charge [D2.6, EM.9];}\ \omega^2=1;\ R\mapsto R^{\mathrm T};\ \omega\notin\operatorname{Aut}(\Gamma),\ |V|\ge3\\
\Theta_{\mathfrak L} & \text{fixed} & \text{fixed} & \leftrightarrow & \leftrightarrow & \text{ledger order}\mapsto\text{reverse};\ J^+\leftrightarrow J^-;\ G(e)\mapsto G(e)^*\ (\alpha_e\widehat{ds}^{\,2}(e)\mapsto-\alpha_e\widehat{ds}^{\,2}(e),\ \sigma_e\ \text{fixed [D3.5]});\ \mathcal J\mapsto\mathcal J^*\ \text{[D3.6]};\ w\mapsto w^*\ \text{[CT2.8(1)]};\ \Theta^2=1\\
\omega\Theta_{\mathfrak L} & \text{fixed} & W^+\!\leftrightarrow\!W^- & \leftrightarrow & \leftrightarrow & \text{the fourth class}
\end{array}
\]

$\omega$ acts on the edge set, $\Theta_{\mathfrak L}$ on phase/ledger order — disjoint data $\Rightarrow$ they commute, and $\langle\omega,\Theta_{\mathfrak L}\rangle\cong\mathbb Z_2\times\mathbb Z_2$ with classes $(1,\omega,\Theta_{\mathfrak L},\omega\Theta_{\mathfrak L})$; $\omega$ is charge-odd [D2.6; covariance of the closed-core charge read: EM.9a], $\Theta_{\mathfrak L}$ phase-odd and charge-even [D2.7, D3.5, D3.6]. $(3)$ \textbf{Edge-class answer (the sector question at edge level).} Real edges split into $\sigma_e=\pm1$ [CT0.9: $\Phi=\{0,\pi\}$]; virtual edges carry $\neg$registration [A5, A7, D0.5: $\phi_e$ defined on $E_{\mathcal R}$ only] — a proposed $(\pm,\mathcal V)$ split has no read: the virtual class is single and ungraded (weight $1$, no phase). The two sign structures ($\sigma$ vs $(-1)^{\aleph}$) are independent [CT0.10]. $(4)$ \textbf{[READING] rows (recorded, $\neg$adopted).} External CPT vocabulary maps as: $C\leftrightarrow\omega$ (the charge-odd reversal [D2.6]; sign-level realization $\sigma_e\mapsto-\sigma_e$: EM.24(1) — \textbf{E}, forward pointer); $T\leftrightarrow\Theta_{\mathfrak L}$ (ledger-time reversal [D2.7] — $\neg$ anti-unitary: no anti-unitary operator is declared; conjugation acts on weights/phases [CT2.8(1), D3.6]); $P\leftrightarrow$ the alignment-pattern reversal $\mathcal P\mapsto E_{\mathcal R}\setminus\mathcal P$ (equivalently $\alpha_e\mapsto\pi-\alpha_e$) with the mirror classification of Def-EM.21/Def-EM.22, EM.23 — \textbf{E}, forward pointers: a pattern-level reversal on configuration data (spacetime parity as a Lorentzian field-algebra read has count $0$ in Parts 1–2; the orientation read $\chi=\operatorname{sign}(k)$ remains [CANDIDATE], not adopted); "CPT" $\leftrightarrow$ the derived composition $C\circ P\circ T$ on $\mathfrak C(\mathcal R)$ with unconditional $|w|^2$-invariance [EM.24(4)–(5) — \textbf{E}, forward pointer]; at the involution-table level its reversal part is the $\omega\Theta_{\mathfrak L}$ class of $(2)$. This is the exact amount of reversal structure the axioms declare — $\neg$ a Lorentzian CPT theorem (no Lorentzian field algebra, no spinor bundle in the inventory).

\textbf{Req.} A3, A5, A7, A8; C3; D0.5, D2.6, D2.7, D3.5, D3.6; CT0.7, CT0.7a, CT0.9, CT0.10, CT0.12, CT1.1, CT1.2, CT1.6, CT1.9, CT1.11 (Part 1); CT2.8(1) (\textbf{C} — forward pointer, re-verification); EM.9 (\textbf{E} — forward pointer, re-verification); exact $\mathbb Q(i)$ matrix arithmetic, exhaustive pattern enumeration, non-multiplicativity witness [math; machine: \texttt{checks/\allowbreak{}verify\_\allowbreak{}merge\_\allowbreak{}tierA.py} §E].

\textbf{Pf.}

%dsp-level:0
\[
(0):\ \text{CT1.2's table verbatim: }\ E_0,
\]

%dsp-level:0
\[
\Sigma,
\]

%dsp-level:0
\[
\hat E^{ij}\in\mathcal E_{++};
\]

%dsp-level:0
\[
\ W_i\in\mathcal E_{+-};
\]

%dsp-level:0
\[
\ J_i\in\mathcal E_{-+};
\]

%dsp-level:0
\[
\ K_i\in\mathcal E_{--}\ \wedge\ \text{CT0.12's aspect rows}\ [\text{machine}]
\]

%dsp-level:0
\[
(1):\ \mathbb S=\tfrac12(E_0+\Sigma)\ [\mathrm{CT0.7}(1)];\quad \operatorname{Ad}(\mathbb S)(X)=\mathbb S X\mathbb S=S(X)\ [\mathrm{CT0.7}(2)];
\]

%dsp-level:0
\[
\qquad T^2=\mathbb 1,\ S^2=\mathbb 1,\ [S,T]=0\ [\mathrm{A3},\ \mathrm{CT1.1}]
\]

%dsp-level:0
\[
T\ \neg\text{multiplicative}:\ T(\sigma_0\otimes\sigma_1)\,
\]

%dsp-level:0
\[
T(\sigma_0\otimes\sigma_2)
\]

%dsp-level:0
\[
= (-\sigma_0\otimes\sigma_1)(-\sigma_0\otimes\sigma_2)
\]

%dsp-level:0
\[
= \sigma_0\otimes i\sigma_3
\]

%dsp-level:0
\[
\ne -\,
\]

%dsp-level:0
\[
i\,
\]

%dsp-level:0
\[
\sigma_0\otimes\sigma_3
\]

%dsp-level:0
\[
= T(\sigma_0\otimes i\sigma_3)\ [\text{machine witness}]
\]

%dsp-level:0
\[
(2):\ D2.6/D2.7\ \text{verbatim on data rows};
\]

%dsp-level:0
\[
\quad J^\pm\ \text{blocks}:\ \Theta_{\mathfrak L}:\ J^+\leftrightarrow J^-\ [\mathrm{D2.7}]\ \wedge\ \mathbf F_{(+,\mathcal V)}
\]

%dsp-level:0
\[
= F_{J^+},
\]

%dsp-level:0
\[
\ \mathbf F_{(-,\mathcal V)}
\]

%dsp-level:0
\[
= F_{J^-}\ [\mathrm{A8}]\implies\text{sector exchange}
\]

%dsp-level:0
\[
\omega:\ \text{charge}
\]

%dsp-level:0
\[
\mapsto -\text{charge}\ [\mathrm{D2.6}]
\]

%dsp-level:0
\[
\implies W^+\leftrightarrow W^-\ \text{inside}\ \mathbf F_{(-,\mathcal R)}\ [\text{charge eigen-decomposition of the weak block: }F_{W^+}\oplus F_{W^-}\oplus F_Z,\ \mathrm{A8}];
\]

%dsp-level:0
\[
\quad \omega\ \text{on}\ J^\pm:\ \neg\text{declared}
\]

%dsp-level:0
\[
\omega^2=1:\ (R^{\mathrm T})^{\mathrm T}=R;
\]

%dsp-level:0
\[
\qquad \omega\notin\operatorname{Aut}(\Gamma):\ p(a)
\]

%dsp-level:0
\[
= b\ \wedge\ p(b)
\]

%dsp-level:0
\[
= a\ \forall(a,b)
\]

%dsp-level:0
\[
\in E:\ \bot\ \text{for}\ |V|
\]

%dsp-level:0
\[
\ge 3\ [\text{machine: exhaustive on the tetrahedron}]
\]

%dsp-level:0
\[
\text{disjoint supports}\ (\text{edge set vs phase/order})\implies[\omega,\Theta_{\mathfrak L}]=0;\qquad \omega^2=\Theta_{\mathfrak L}^2=1\implies\mathbb Z_2\times\mathbb Z_2.\ \blacksquare
\]

\textbf{Phys.} The corpus declares exactly two reversals and they commute: one flips orientation and charge on the edge set, the other flips ledger order and conjugates phases — four classes in total. The four force sectors of A8 are the joint grading atlas of the pair register; the Janus pair is the only sector pair a declared reversal exchanges, and spacetime parity has no force-sector house here — its pattern-level read lives on the alignment data (EM.23, EM.24 — \textbf{E}). The table is the complete reversal inventory at sector level: nothing further is derivable without a new declaration.

\end{theorem}

\begin{theorem}[{\texorpdfstring{CT1.11b — The native Clifford frames of the pair register and the weak-covariance pincer}{CT1.11b - The native Clifford frames of the pair register and the weak-covariance pincer}}]\label{tag:CT1-11b}
\textbf{Thm.} On $\mathfrak H(\mathbb C^4)$ with CT1.2's named basis and CT0.7's $\mathbb S$, $\Pi_{\mathcal R}=\tfrac14(3E_0+\Sigma)$, $\Pi_{\mathcal V}=\tfrac14(E_0-\Sigma)$. Call the $15$ non-identity $\sigma_\mu\otimes\sigma_\nu$ [A3, CT1.1] the \textit{Pauli products} — each Hermitian, squaring to $+I_4$, traceless. A \textit{frame} is a $4$-subset of pairwise-anticommuting Pauli products; for a frame $\mathcal F=(e_1,e_2,e_3,e_4)$ let $\gamma_{\mathcal F}:=e_1e_2e_3e_4$ (its \textit{pseudoscalar}; frame- and order-dependent: every transposition flips its sign). $(1)$ \textbf{Native frames.} There are exactly $30$ frames (unordered; $720$ as ordered tuples) among the $\binom{15}{4}=1365$ candidates. For every frame: $\gamma_{\mathcal F}$ is Hermitian, $\gamma_{\mathcal F}^2=I_4$, $\mathfrak T\gamma_{\mathcal F}=0$, $\{\gamma_{\mathcal F},e_i\}=0$; and the $16$ real products $\{e_{i_1}\cdots e_{i_r}\}_{r=0..4}$ are $\mathbb R$-linearly independent — the frame generates a real subalgebra isomorphic to $Cl(4,0)$ acting on $\mathbb C^4\cong\mathbb R^8$, \textbf{with no new axiom}: A3's Pauli facts and CT1.1's basis suffice. The signature is Euclidean ($e_i^2=+I_4$); the representative $\eta=(-,+,+,+)$ [C13] of A3's sign vector labels the $T$-grading only [CT1.1] and enters no Clifford relation. No canonical pseudoscalar exists: a distinguished "$\gamma$-element" would be a choice, $\neg$a derivation. $(2)$ \textbf{The pseudoscalar catalog (exhaustive).} All $15$ Pauli products occur as $\pm\gamma_{\mathcal F}$, \textbf{each from exactly $2$ frames}, in three species:

%dsp-level:0
\[
\begin{array}{c|c|c|c}
\text{species} & \gamma_{\mathcal F}\ (\text{up to sign}) & [\gamma_{\mathcal F},\mathbb S] & \text{commuting}\ W_i\\ \hline
\text{mixed}\ (6) & \sigma_0\otimes\sigma_i\ /\ \sigma_i\otimes\sigma_0 & \pm2iJ_i & W_i\ \text{only}\\
\text{diagonal}\ (3) & \sigma_i\otimes\sigma_i\in\mathcal E_{++} & 0 & W_i\ \text{only (anticomm.}\ W_j,W_k)\\
\text{off-diag. spatial}\ (6) & \sigma_i\otimes\sigma_j\ (i\ne j) & \mp2iK_k,\ k=6{-}i{-}j & \text{none (anticomm.}\ W_k)
\end{array}
\]

No pseudoscalar commutes with more than one $W_i$; none anticommutes with more than two. Witness frame $\mathcal F_0=(\sigma_1\otimes\sigma_2,\ \sigma_2\otimes\sigma_2,\ \sigma_3\otimes\sigma_2,\ \sigma_0\otimes\sigma_1)$:

%dsp-level:0
\[
\gamma_{\mathcal F_0}=\sigma_0\otimes\sigma_3=\operatorname{diag}(1,-1,1,-1);
\]

%dsp-level:0
\[
\qquad P^{(\mathcal F_0)}_\pm:=\tfrac12(I_4\mp\gamma_{\mathcal F_0})=\operatorname{diag}(0,1,0,1)\ /\ \operatorname{diag}(1,0,1,0)
\]

— both idempotent, orthogonal, summing to $I_4$, trace $2$, rank $2$ (exact forms, machine). $(3)$ \textbf{The weak-covariance pincer.} No grading involution of the pair register can be weak-covariant, from either side: \textbf{(a) commuting side.} $\{M:[M,W_i]=0\ \forall i\}=\operatorname{span}_{\mathbb C}\{\Pi_{\mathcal R},\Pi_{\mathcal V}\}$ exactly (complex dimension $2$ — Schur on the $3\oplus1$ decomposition [CT1.10, CT0.7]; machine nullspace solve). Its involutions are $a\Pi_{\mathcal R}+b\Pi_{\mathcal V}$ with $a^2=b^2=1$, traces $3a+b\in\{\pm4,\pm2\}$; a pseudoscalar needs trace $0$ (cyclic flip: $\mathfrak T(e_1e_2e_3e_4)=-\mathfrak T(e_2e_3e_4e_1)$) — $3a+b=0\ \wedge\ a^2=b^2=1$ has \textbf{no solution}. \textbf{(b) anticommuting side.} $\{M:\{M,W_i\}=0\ \forall i\}=\mathbb C\cdot\Pi_{\mathcal V}$ exactly (complex dimension $1$; machine nullspace solve): rank $1$, trace $1$ — singular, never an involution. Structural reason: an invertible $M$ anticommuting with every $W_i$ would make $\operatorname{Ad}(M)$ a Lie automorphism of $\operatorname{span}\{W_i\}$ equal to $-\mathrm{id}$ — but $-\mathrm{id}$ reverses the bracket ($[W_1,W_2]=iW_3\mapsto-iW_3\ne iW_3$) [CT1.4]: contradiction. \textbf{(c) exhaustive.} Every one of the $15$ attained pseudoscalars fails (a) and (b) explicitly, entrywise: for each, $[\gamma,W_i]\ne0$ for some $i$ and $\{\gamma,W_i\}\ne0$ for some $i$. \textbf{(d) carrier side.} On any $2$-dim carrier with weak action $W_i\mapsto\tfrac12\sigma_i$ (the weak seed locus: EW.32): $\{m:[m,\sigma_i]=0\ \forall i\}=\mathbb C\cdot I_2$ and $\{m:\{m,\sigma_i\}=0\ \forall i\}=\{0\}$ — the weak seed representation is irreducible: no nontrivial projector commutes with the weak action, and no grade anticommutes with it. $(4)$ \textbf{The native selectivity (what the weak force does grade).} The weak force grades the pair arena by \textbf{$S$-parity}: $W_i\Pi_{\mathcal V}=0$ exactly $\forall i$; $\sum_iW_i^2=2\Pi_{\mathcal R}$ [Cor CT0.7a]; $\bigcap_i\ker W_i=\operatorname{im}\Pi_{\mathcal V}$ exactly (the $1$-dim singlet line; machine) — a $3$-active/$1$-inactive split. This is a force-arena selectivity; a matter grade would have to live on a carrier on which the gauge action is trivial — the opposite pattern. $(5)$ \textbf{Non-alignment.} $[\gamma_{\mathcal F},\mathbb S]=0\iff\gamma_{\mathcal F}=\pm\,\sigma_i\otimes\sigma_i$: the $S$-definite pseudoscalars are exactly the three diagonal ones — they live in $\mathcal E_{++}$, the gravity/Higgs block [CT1.1, CT1.2], and are still $\neg$weak-covariant (each anticommutes with the two non-own $W$: e.g.\ $[\sigma_3\otimes\sigma_3,\ W_1]=2iE^{23}$ exactly — the weak rotation turns an $S$-definite pseudoscalar into gravitational shear). Under the swap, off-diagonal pseudoscalar grades rotate into the Janus seal sector: $[\sigma_0\otimes\sigma_3,\ \mathbb S]=2iJ_3$. The frame splits and the $S$-split provably do not align: $\operatorname{rank}(P^{(\mathcal F_0)}_\pm\Pi_{\mathcal R})=2$, $\operatorname{rank}(P^{(\mathcal F_0)}_\pm\Pi_{\mathcal V})=1$ — cross-ranks $(2,1,2,1)$ against the $S$-split's $(3,1)$.

\textbf{Req.} A3; C3; CT0.7, CT0.7a, CT1.1, CT1.2, CT1.4, CT1.10, CT1.11a (Part 1); EW.32 (\textbf{D}; seed-carrier locus — forward pointer, corroboration only: clause (3d) is self-contained by Schur [math]); exhaustive frame enumeration, exact nullspace solves over $\mathbb Q(i)$, Schur's lemma [math; machine: \texttt{checks/\allowbreak{}verify\_\allowbreak{}merge\_\allowbreak{}tierC.py}].

\textbf{Pf.}

%dsp-level:0
\[
(1):\ \gamma_{\mathcal F}^2=(e_1e_2e_3e_4)(e_1e_2e_3e_4)=(-1)^{\binom42}e_1^2e_2^2e_3^2e_4^2=(-1)^6I_4=I_4;
\]

%dsp-level:0
\[
\qquad \mathfrak T\gamma_{\mathcal F}=-\mathfrak T(e_2e_3e_4e_1)=-\mathfrak T\gamma_{\mathcal F}\implies\mathfrak T\gamma_{\mathcal F}=0\ [\text{cyclic}]
\]

%dsp-level:0
\[
\{\gamma_{\mathcal F},e_i\}=0:\ e_i\ \text{anticommutes with the three other factors},\ \text{commutes with itself};
\]

%dsp-level:0
\[
\qquad 30\ \text{frames}:\ \text{exhaustive enumeration of}\ 1365\ \text{candidates}\ [\text{machine}]
\]

%dsp-level:0
\[
\mathbb R\text{-independence of the}\ 16\ \text{products}:\ \text{exact elimination on}\ \mathbb Q\text{-flattened components, all}\ 30\ \text{frames}\ [\text{machine}]
\]

%dsp-level:0
\[
(2):\ \text{canonical-sign census over the}\ 30\ \text{frames}\ [\text{machine}];
\]

%dsp-level:0
\[
\quad \text{species}=\text{index patterns};
\]

%dsp-level:0
\[
\quad [\sigma_0\otimes\sigma_3,\mathbb S]=2iJ_3,\ [\sigma_3\otimes\sigma_3,W_1]=2iE^{23}:\ \text{direct products}\ [\text{machine}]
\]

%dsp-level:0
\[
\mathcal F_0:\ \text{direct multiplication}:\ \gamma_{\mathcal F_0}=\sigma_0\otimes\sigma_3;
\]

%dsp-level:0
\[
\quad P_\pm^2=P_\pm,\ P_+P_-=0,\ P_++P_-=I_4,\ \mathfrak T P_\pm=2\ [\text{machine}]
\]

%dsp-level:0
\[
(3a):\ \text{nullspace of the}\ 3\times16\ \text{complex linear system}\ [X,W_i]
\]

%dsp-level:0
\[
= 0\ \text{in the Pauli basis}\
\]

%dsp-level:0
\[
= \ \operatorname{span}\{\Pi_{\mathcal R},\Pi_{\mathcal V}\}\ [\text{machine, exact}];
\]

%dsp-level:0
\[
\quad \text{Schur on}\ 3\oplus1\ [\mathrm{CT1.10};\ \text{math}]
\]

%dsp-level:0
\[
\text{involutions in the commutant}:\ (a\Pi_{\mathcal R}+b\Pi_{\mathcal V})^2=I_4\iff a^2=b^2=1;
\]

%dsp-level:0
\[
\quad \mathfrak T=3a+b\in\{4,2,-2,-4\}:\ \neg\exists\ \text{trace-0 involution}
\]

%dsp-level:0
\[
(3b):\ \text{nullspace of}\ \{X,W_i\}=0\ =\ \mathbb C\cdot\Pi_{\mathcal V}\ [\text{machine}];\qquad W_i\Pi_{\mathcal V}=0\ [(4)]\implies\operatorname{rank}\Pi_{\mathcal V}=1:\ \text{singular}
\]

%dsp-level:0
\[
\text{invertible case}:\ \operatorname{Ad}(M)W_i
\]

%dsp-level:0
\[
= -W_i\implies\operatorname{Ad}(M)[W_1,W_2]
\]

%dsp-level:0
\[
= [-W_1,-W_2]
\]

%dsp-level:0
\[
= [W_1,W_2]
\]

%dsp-level:0
\[
= iW_3,
\]

%dsp-level:0
\[
\ \text{but}\ \operatorname{Ad}(M)(iW_3)
\]

%dsp-level:0
\[
= -iW_3:\ \bot\ [\mathrm{CT1.4}]
\]

%dsp-level:0
\[
(3c):\ \text{entrywise for all}\ 15\ [\text{machine}];\qquad (3d):\ 2\times2\ \text{nullspace solves}\ [\text{machine};\ \text{Schur}]
\]

%dsp-level:0
\[
(4):\ W_i\ \text{preserves}\ \vee^2,\wedge^2\ [\mathrm{CT0.7}(3)];
\]

%dsp-level:0
\[
\quad W_i|_{\wedge^2}=0\ (\text{singlet});
\]

%dsp-level:0
\[
\quad \sum_iW_i^2=\tfrac32E_0+\tfrac12\Sigma=2\Pi_{\mathcal R}\ [\text{Cor CT0.7a}];
\]

%dsp-level:0
\[
\quad \bigcap\ker W_i:\ \text{machine}
\]

%dsp-level:0
\[
(5):\ [\sigma_\mu\otimes\sigma_\nu,\mathbb S]=0\iff\sigma_\mu\otimes\sigma_\nu=\sigma_\nu\otimes\sigma_\mu\iff\mu=\nu\ [\text{machine over all}\ 15];
\]

%dsp-level:0
\[
\quad \text{ranks by exact elimination}.\ \blacksquare
\]

\textbf{Phys.} The pair register carries a full Euclidean Clifford structure natively — thirty frames, fifteen pseudoscalars, all of it derived from A3's Pauli facts with no declaration anywhere — and the weak force selects exactly one split inside it: the $3$-active/$1$-inactive $S$-parity grading. But the pincer is exact: from the commuting side the only weak-covariant candidates are the two $S$-projections, whose involutions all have nonzero trace; from the anticommuting side the candidate set collapses to a singular singlet line. No pseudoscalar grade of this register can ever behave as a weak-covariant matter grade — the obstruction is not a shortage of matrices but the absence of a spacetime-side carrier on which the gauge action is trivial, and that absence is a declaration the corpus has not made.

\textbf{Rmk (deficit relocation record).} The completeness audit's first deficit is hereby \textit{relocated inside Parts 1–2}: native frames exist (clause (1)); their grades cannot be weak-covariant (clause (3)); a Dirac-type operator would additionally require a per-edge frame assignment $e\mapsto\mathcal F_e$ coupled to the EW.32 transport — an enumerated would-be declaration of the audit layer, recorded here as \textbf{not adopted, no tag issued}. The reversal structure declared by the corpus remains exactly $\langle\omega,\Theta_{\mathfrak L}\rangle\cong\mathbb Z_2\times\mathbb Z_2$ [CT1.11a].

\end{theorem}

\phantomsection
\section[Tier 2 — Lyapunov Stability of the Pair State]{\texorpdfstring{Tier 2 — Lyapunov Stability of the Pair State}{Tier 2 - Lyapunov Stability of the Pair State}}

\begin{theorem}[{\texorpdfstring{L0 — The ledger–tilt coupling}{L0 - The ledger-tilt coupling}}]\label{tag:L0}
\textbf{Thm.} \textit{Dynamical scope (declared at use — the crystallization–tilt coupling; CT4.5a precedent: A4–A7 declare the criterion, the tick and the ledger, $\neg$ a state update at the seal):} every A4 crystallization event applies exactly one Gibbs tilt step [D0.3] to the pair state of that event. Under A5 ($\rho_\circ=\tfrac14I$), A6 ($\mathcal T(\rho)=U\rho U^\dagger$, $U\sim\mu_\star$; $|\mathfrak L(\tau_{n+1})|=|\mathfrak L(\tau_n)|+n_{\mathrm{real\_out}}(\tau_n)$), A7 ($\mathfrak L=\bigcup_{n\ge0}\mathfrak L(\tau_n)$, $|\mathfrak L(\tau_n)|=\sum_{k<n}n_{\mathrm{real\_out}}(\tau_k)$), D0.3/D0.4 ($\gamma:=1-4\varepsilon_q$), CT0.8: $(1)$ \textbf{Tick preserves $V$.} $\ V\bigl(\tilde K(\mathcal T(\rho))\bigr)=V\bigl(\tilde K(\rho)\bigr)$ for every draw [A6; L1(2)]. $(2)$ \textbf{Events apply the tilt.} At tick $\tau_n$ the A4 criterion selects $n_{\mathrm{real\_out}}(\tau_n)$ edges into the ledger [A6]; under the declared scope each selection maps its pair state

%dsp-level:0
\[
\rho\ \longmapsto\ \tilde{\mathcal T}_\gamma(\rho)=\frac{\rho\,e^{-\gamma\rho}}{\mathfrak T(\rho\,e^{-\gamma\rho})},\qquad
\gamma=1-4\varepsilon_q\in(0,1]\quad[\mathrm{D0.3};\ \mathrm{CT0.8}];
\]

$Z_\gamma>0$ on faithful $\rho$ [Cor L0a] $\implies\tilde{\mathcal T}_\gamma(\rho)$ faithful, $\succeq0$, $\mathfrak T=1$ [D0.3]; $\ \tilde{\mathcal T}_\gamma\ne\mathcal T$ [D0.3]. $(3)$ \textbf{Tilt contracts $V$.} $\ V'\le V$, equality iff $\rho=\tfrac14I$ [A5; L1(5)]. $(4)$ \textbf{Ledger counts tilt steps.} $\#\{\text{tilt steps at }\tau_n\}=n_{\mathrm{real\_out}}(\tau_n)$; cumulative count $|\mathfrak L(\tau_n)|=\sum_{k<n}n_{\mathrm{real\_out}}(\tau_k)$ [A6, A7 — exact integer recursion]. On one pair channel with $k$ events interleaved with arbitrary draws, the $V$-sequence is that of $\tilde{\mathcal T}_\gamma^{k}(\rho_0)$ — the interleaving is spectrally inert by $(1)$; per fixed $k$: $\ \mu_k=\bigl(1-\tfrac\gamma4\bigr)^{k}\mu_0+O(\mu_0^2)$ [L2, iterated]. $(5)$ \textbf{Contraction is the ledger dynamics.} The deficits of L1–L3 are realized once per ledger increment: the deficit fraction per tilt step is $\tfrac{\gamma(8-\gamma)}{16}$ of $V$ to exact leading order [L2], accumulating as $V\mapsto\bigl(1-\tfrac\gamma4\bigr)^{2k}V+O(\delta\rho_0^3)$ after $k$ recorded events; the inertial read $m_i^2(q)=\Delta S(q)$ [L3(4)] is the deficit read per increment.

\textbf{Req.} A4; A5; A6; A7; D0.2; D0.3; D0.4; D1.2; CT0.8; L1; L2 (forward pointers — the derivation order of \texttt{index/\allowbreak{}dependency\_\allowbreak{}tree.md} places this tag after L2); the declared crystallization–tilt dynamical scope (theorem-local; see Status); exact integer recursion, exact Haar covariance over $\mathbb Q(i)$, exact rational $\log/\exp$-enclosures [math; machine: \texttt{checks/\allowbreak{}verify\_\allowbreak{}L0\_\allowbreak{}ledger\_\allowbreak{}tilt.py} §A–§D].

\textbf{Pf.}

%dsp-level:0
\[
(1):\ K(\mathcal T\rho)=UKU^\dagger\ [\text{functional-calculus covariance}];
\]

%dsp-level:0
\[
\qquad \mathfrak T\bigl(K(\mathcal T\rho)\,\mathcal T\rho\bigr)=\mathfrak T(K\rho)\ [\text{cyclicity}]\implies\tilde K(\mathcal T\rho)=U\tilde KU^\dagger\implies V'=V\ [\mathrm{L1}(2)]
\]

%dsp-level:0
\[
(2):\ \varepsilon_q\in[0,\tfrac14)\ [\mathrm{CT0.8}]\implies\gamma=1-4\varepsilon_q\in(0,1]\ [\text{exact endpoints}:\ \varepsilon_q=0\mapsto\gamma=1];
\]

%dsp-level:0
\[
\qquad Z_\gamma\in(0,\,e^{-\gamma/4}]\ [\text{Cor L0a}]\implies\lambda_i'=\lambda_ie^{-\gamma\lambda_i}/Z_\gamma>0\ [\exp>0\ \text{on exact enclosures}]
\]

%dsp-level:0
\[
\implies\tilde{\mathcal T}_\gamma(\rho)\ \text{faithful},\ \succeq0,\ \mathfrak T\bigl(\tilde{\mathcal T}_\gamma(\rho)\bigr)=1\ [\mathrm{D0.3}\ \text{normalization}]
\]

%dsp-level:0
\[
(3):\ \mathrm{L1}(5)\ [\text{Jensen on}\ \phi_\gamma;\ \text{zero locus}\ \rho=\tfrac14I]
\]

%dsp-level:0
\[
(4):\ |\mathfrak L(\tau_0)|
\]

%dsp-level:0
\[
= 0\ [\mathrm{A7}];
\]

%dsp-level:0
\[
\quad\text{recursion}\ [\mathrm{A6}]
\]

%dsp-level:0
\[
\implies |\mathfrak L(\tau_n)|
\]

%dsp-level:0
\[
= \textstyle\sum_{k<n}n_{\mathrm{real\_out}}(\tau_k)\ [\mathrm{A7}]\ \text{by induction on}\ n\ [\text{exact integers; machine}]
\]

%dsp-level:0
\[
\operatorname{spec}\bigl(\mathcal T(\rho)\bigr)
\]

%dsp-level:0
\[
= \operatorname{spec}(\rho):\ \mathfrak T\bigl((U\rho U^\dagger)^m\bigr)
\]

%dsp-level:0
\[
= \mathfrak T(\rho^m),
\]

%dsp-level:0
\[
\ m
\]

%dsp-level:0
\[
= 1..4\ [\text{cyclicity; machine, exact over}\ \mathbb Q(i)]
\]

%dsp-level:0
\[
\implies V\circ\mathcal T\qquad = V\ [(1)]
\]

%dsp-level:0
\[
\implies V\bigl(\tilde{\mathcal T}\,\mathcal T\,\tilde{\mathcal T}\cdots(\rho_0)\bigr)=V\bigl(\tilde{\mathcal T}^{\,k}(\rho_0)\bigr)\ [\text{composition}];
\]

%dsp-level:0
\[
\qquad \mu_k
\]

%dsp-level:0
\[
= \bigl(1-\tfrac\gamma4\bigr)^k\mu_0+O(\mu_0^2):\ \text{induction on L2's}\ \mu'
\]

%dsp-level:0
\[
= \bigl(1-\tfrac\gamma4\bigr)\mu+O(\mu^2)\ [\text{coefficient-wise composition; machine},\ k\le3]
\]

%dsp-level:0
\[
(5):\ \mathrm{L2}\ \text{per step}:\ V'-V=-\tfrac{\gamma(8-\gamma)}{16}V+O(\delta\rho^3);\qquad
\text{accumulation at leading order from}\ (4).\ \blacksquare
\]

\end{theorem}

\begin{corollary}[{\texorpdfstring{L0a — tilt modular update}{L0a - tilt modular update}}]\label{tag:L0a}
Faithful $\rho$, $K:=-\log\rho$ [D0.3], $K':=K\bigl(\tilde{\mathcal T}_\gamma(\rho)\bigr)$, $Z_\gamma:=\mathfrak T(\rho\,e^{-\gamma\rho})$:

%dsp-level:0
\[
K'=K+\gamma\rho+\log Z_\gamma\cdot I;\qquad
Z_\gamma\in(0,\,e^{-\gamma/4}];\qquad
\mathfrak T(K'-K)=\gamma+4\log Z_\gamma\le0,
\]

with equality iff $\rho=\tfrac14I$; $\ \gamma=1:\ Z_1\le e^{-1/4}$, $\ \mathfrak T(K'-K)=1+4\log Z_1\le0$.

\textbf{Req (Cor).} A1; A5; D0.3; CT0.8; Jensen's inequality for $p\,e^{-\gamma p}$, exact rational $\log/\exp$-enclosures [math; machine: \texttt{checks/\allowbreak{}verify\_\allowbreak{}L0\_\allowbreak{}ledger\_\allowbreak{}tilt.py} §B–§C].

\textbf{Pf (Cor).}

%dsp-level:0
\[
\rho,\ \tilde{\mathcal T}_\gamma(\rho)\ \text{simultaneously diagonalizable}\ [\text{functional calculus on a commuting family}];
\]

%dsp-level:0
\[
\qquad \lambda_i'=\lambda_ie^{-\gamma\lambda_i}/Z_\gamma
\]

%dsp-level:0
\[
\implies -\log\lambda_i'
\]

%dsp-level:0
\[
= -\log\lambda_i+\gamma\lambda_i+\log Z_\gamma\ [\text{exact spectral identity; machine at two enclosure tolerances}]
\]

%dsp-level:0
\[
\implies K'
\]

%dsp-level:0
\[
= K+\gamma\rho+\log Z_\gamma\cdot I
\]

%dsp-level:0
\[
g(p):=p\,e^{-\gamma p}:\ \ g''(p)=e^{-\gamma p}\bigl(\gamma^2p-2\gamma\bigr)\le0\ \ \text{on}\ \ [0,1],\ \gamma\in(0,1]
\]

%dsp-level:0
\[
\quad [\gamma p\le1<2;\ \text{exact rational sign}]
\]

%dsp-level:0
\[
\text{Jensen}:\ Z_\gamma
\]

%dsp-level:0
\[
= \textstyle\sum_ig(\lambda_i)
\]

%dsp-level:0
\[
\le 4\,
\]

%dsp-level:0
\[
g\bigl(\tfrac14\bigr)
\]

%dsp-level:0
\[
= e^{-\gamma/4},
\]

%dsp-level:0
\[
\ \ \text{equality iff}\ \lambda_i
\]

%dsp-level:0
\[
= \tfrac14\ \forall i\iff\rho
\]

%dsp-level:0
\[
= \tfrac14I\ [\text{strict concavity}];
\]

%dsp-level:0
\[
\qquad Z_\gamma>0\ [\lambda_i>0]
\]

%dsp-level:0
\[
\mathfrak T(K'-K)=\gamma\,\mathfrak T\rho+4\log Z_\gamma=\gamma+4\log Z_\gamma\le\gamma+4\bigl(-\tfrac\gamma4\bigr)=0;
\]

%dsp-level:0
\[
\qquad \gamma=1:\ Z_1\le e^{-1/4},\ \ \mathfrak T(K'-K)=1+4\log Z_1.\ \blacksquare
\]

\textbf{Phys.} The Haar tick rotates the pair state without touching the Lyapunov function; the ledger grows exactly at the A4 events; under the declared crystallization dynamics every event is one Gibbs tilt step — the only contraction a pair state undergoes. The monotonicity of L1–L3 is thereby the per-increment dynamics of the ledger: ticks rotate, events contract, and $|\mathfrak L(\tau_n)|$ is the tilt-step count. Each tilt step updates the modular Hamiltonian by $K\mapsto K+\gamma\rho+\log Z_\gamma\cdot I$, whose trace deficit $\gamma+4\log Z_\gamma\le0$ vanishes only at the A5 vacuum. No particle identification.

\end{corollary}

\begin{theorem}[{\texorpdfstring{L1 — The Lyapunov function of the pair state}{L1 - The Lyapunov function of the pair state}}]\label{tag:L1}
\textbf{Thm.} Faithful pair states $\rho$ [A1: $\rho\in\mathfrak H(\mathbb C^4)$, $\rho\succeq0$, $\mathfrak T\rho=1$, $\operatorname{spec}\rho\subset(0,1]$]; with

%dsp-level:0
\[
\delta\rho:=\rho-\tfrac14I\ \ [\mathrm{A5}:\ \rho_\circ=\tfrac14I];\qquad
\mu:=\operatorname{spec}(\delta\rho)\subset\bigl[-\tfrac14,\tfrac34\bigr];\qquad
\textstyle\sum_i\mu_i=0\ \ [\mathfrak T\rho=1];
\]

%dsp-level:0
\[
\gamma:=1-4\varepsilon_q\in(0,1]\ \ [\mathrm{D0.3};\ \mathrm{A4},\ \mathrm{D0.4},\ \mathrm{CT0.8}];
\]

%dsp-level:0
\[
K:=-\log\rho;\qquad
\tilde K:=K-\mathfrak T(K\rho)\,I\ \ [\mathrm{D0.3}];\qquad
V(\tilde K):=\mathfrak T\bigl(\tilde K^2\bigr)\ \ge0;
\]

$(1)$ \textbf{Zero locus.} $V=0\iff\tilde K=0\iff\rho=\tfrac14I$ — the A5 vacuum. $(2)$ \textbf{Tick invariance.} $V\bigl(\tilde K(\mathcal T(\rho))\bigr)=V\bigl(\tilde K(\rho)\bigr)$ for every draw [A6: $\mathcal T(\rho)=U\rho U^\dagger$]. $(3)$ \textbf{Decay identity ($\gamma$-dependent).} Under the Gibbs tilt $\rho':=\tilde{\mathcal T}_\gamma(\rho)$ [D0.3], with $V':=V\bigl(\tilde K(\rho')\bigr)$:

%dsp-level:0
\[
V'-V=-\gamma\sum_i\phi_\gamma(\mu_i)+O(\delta\rho^4);\qquad
\phi_\gamma(x):=2x\log(1+4x)-\gamma x^2;
\]

%dsp-level:0
\[
\text{leading}:\ \ V'-V=-\gamma(8-\gamma)\textstyle\sum_i\mu_i^2+O(\delta\rho^3);
\]

%dsp-level:0
\[
\qquad \gamma=1:\ \ V'-V=-\textstyle\sum_i\phi(\mu_i)+O(\delta\rho^4),\ \ \phi:=\phi_1 .
\]

$(4)$ \textbf{$\phi_\gamma$-calculus.} $\phi_\gamma(0)=\phi_\gamma'(0)=0$; $\ \phi_\gamma=\phi_1+(1-\gamma)x^2$;

%dsp-level:0
\[
\phi_\gamma''(x)=\frac{14+16x-32x^2}{(1+4x)^2}+2(1-\gamma)\quad[\text{exact rational identity}];
\]

%dsp-level:0
\[
\phi_\gamma\ \text{convex on}\ W_\gamma\supseteq W:=\Bigl[\tfrac{1-2\sqrt2}{4},\tfrac{1+2\sqrt2}{4}\Bigr]\supseteq\bigl[-\tfrac14,\tfrac34\bigr]\quad\forall\gamma\in(0,1];
\]

%dsp-level:0
\[
\qquad \phi_\gamma>0\ \ \text{on}\ \ \bigl[-\tfrac14,\tfrac34\bigr]\setminus\{0\}.
\]

$(5)$ \textbf{Monotonicity.} Jensen: $\ \sum_i\phi_\gamma(\mu_i)\ \ge\ 4\,\phi_\gamma\bigl(\tfrac14\sum_i\mu_i\bigr)=0\ $; hence, on the faithful simplex, for every $\gamma\in(0,1]$:

%dsp-level:0
\[
V'\le V,\qquad V'=V\iff\tilde K=0\iff\rho=\tfrac14I .
\]

\textbf{Req.} A1; A4; A5; A6; D0.3, D0.4; CT0.8 (Part 1); exact formal-power-series arithmetic over $\mathbb Q$ with rational $\gamma$, exact $\mathbb Q(\sqrt2)$ order, Jensen's inequality, exact rational $\log$-enclosures, continuity–density extension [math; machine: \texttt{checks/\allowbreak{}verify\_\allowbreak{}L\_\allowbreak{}lyapunov.py}].

\textbf{Pf.}

%dsp-level:0
\[
(1):\ \tilde K=0\iff K=\mathfrak T(K\rho)I\iff\rho=e^{-c}I\iff\rho=\tfrac14I\ [\mathfrak T\rho=1];
\]

%dsp-level:0
\[
\qquad V=\mathfrak T(\tilde K^2)\ge0\ [\tilde K\ \text{Hermitian}],\quad V=0\iff\tilde K=0
\]

%dsp-level:0
\[
(2):\ K(\mathcal T\rho)=-\log\bigl(U\rho U^\dagger\bigr)=UKU^\dagger\ [\text{functional-calculus covariance}];
\]

%dsp-level:0
\[
\qquad \mathfrak T\bigl(K(\mathcal T\rho)\,\mathcal T\rho\bigr)=\mathfrak T\bigl(UK\rho U^\dagger\bigr)=\mathfrak T(K\rho)\ [\text{cyclicity}]
\]

%dsp-level:0
\[
\implies\tilde K(\mathcal T\rho)=U\tilde KU^\dagger\implies V'=V\ [\text{cyclicity}]
\]

%dsp-level:0
\[
(3):\ \rho\ \text{faithful}\implies\rho,
\]

%dsp-level:0
\[
\ \tilde{\mathcal T}_\gamma(\rho)\ \text{simultaneously diagonalizable}:\ \rho
\]

%dsp-level:0
\[
= \operatorname{diag}(\lambda),
\]

%dsp-level:0
\[
\ \lambda_i\qquad = \tfrac14+\mu_i;
\]

%dsp-level:0
\[
\qquad \tilde K=\operatorname{diag}\bigl(-\log\lambda_i-S\bigr),\ \ S:=\mathfrak T(K\rho)=-\textstyle\sum_j\lambda_j\log\lambda_j
\]

%dsp-level:0
\[
\lambda_i'=\lambda_ie^{-\gamma\lambda_i}/Z_\gamma,\ \ Z_\gamma=\textstyle\sum_j\lambda_je^{-\gamma\lambda_j};
\]

%dsp-level:0
\[
\qquad \log\lambda_i'=\log\lambda_i-\gamma\lambda_i-\log Z_\gamma;
\]

%dsp-level:0
\[
\qquad e^{-\gamma/4}\ \text{cancels in the ratio}\implies\text{exact}\ \mathbb Q[[t]]\ \text{data}
\]

%dsp-level:0
\[
\text{formal series along}\ \rho(t)
\]

%dsp-level:0
\[
= \tfrac14I+t\,
\]

%dsp-level:0
\[
\delta_0\ (\delta_0\ \text{Hermitian traceless, exact rational data}),
\]

%dsp-level:0
\[
\ \gamma\in\{\tfrac14,\tfrac12,\tfrac34,1\}:
\]

%dsp-level:0
\[
V'-V+\gamma\textstyle\sum_i\phi_\gamma(\mu_i(t))\ \text{has vanishing}\ t^0..t^3\ \text{coefficients}\ [\text{machine, exact}];
\]

%dsp-level:0
\[
\qquad t^4\ \text{coefficient}\ne0\ \text{generically}\implies O(\delta\rho^4)\ \text{sharp}
\]

%dsp-level:0
\[
\phi_\gamma(x)
\]

%dsp-level:0
\[
= 2x\bigl(4x-8x^2+\tfrac{64}3x^3-\dots\bigr)-\gamma x^2
\]

%dsp-level:0
\[
= (8-\gamma)x^2-16x^3+O(x^4)\implies\text{leading}:\ V'-V
\]

%dsp-level:0
\[
= -\gamma(8-\gamma)\textstyle\sum_i\mu_i^2+O(\delta\rho^3)
\]

%dsp-level:0
\[
(4):\ \phi_\gamma'(x)=2\log(1+4x)+\tfrac{8x}{1+4x}-2\gamma x;
\]

%dsp-level:0
\[
\qquad \phi_\gamma''(x)=\tfrac8{1+4x}+\tfrac8{(1+4x)^2}-2\gamma=\tfrac{14+16x-32x^2}{(1+4x)^2}+2(1-\gamma)\ [\text{common denominator}]
\]

%dsp-level:0
\[
\gamma\in(0,1]:\ 2(1-\gamma)\ge0\implies\phi_\gamma''\ge\phi_1''\ \text{pointwise}\implies W_\gamma\supseteq W;
\]

%dsp-level:0
\[
\qquad 16x^2-8x-7=0\iff x=\tfrac{1\pm2\sqrt2}4\ [\mathbb Q(\sqrt2),\ \text{exact}];
\]

%dsp-level:0
\[
\qquad W\supseteq[-\tfrac14,\tfrac34]\ [\sqrt2\ge1]
\]

%dsp-level:0
\[
\phi_\gamma=\phi_1+(1-\gamma)x^2;
\]

%dsp-level:0
\[
\qquad \phi_1\ \text{strictly convex on}\ W,\ \phi_1(0)=\phi_1'(0)=0\implies\phi_1>0\ \text{off}\ 0\implies\phi_\gamma>0\ \text{off}\ 0\ \text{on}\ [-\tfrac14,\tfrac34]
\]

%dsp-level:0
\[
(5):\ \text{Jensen on convex}\ \phi_\gamma:\ \tfrac14\sum_i\phi_\gamma(\mu_i)\ge\phi_\gamma\bigl(\tfrac14\sum_i\mu_i\bigr)=\phi_\gamma(0)=0;
\]

%dsp-level:0
\[
\qquad =\iff\mu=0\ \ [\textstyle\sum\mu=0;\ \text{strict convexity}]
\]

%dsp-level:0
\[
\gamma>0\implies V'-V=-\gamma\textstyle\sum_i\phi_\gamma(\mu_i)+O(\delta\rho^4)\le0\ \text{near the vacuum};
\]

%dsp-level:0
\[
\qquad \text{globally: exact rational}\ \log\text{-enclosure sweep},\ \gamma\in\{\tfrac14,\tfrac12,\tfrac34,1\}\ [\text{machine}]
\]

%dsp-level:0
\[
\text{extension to every}\ \gamma
\]

%dsp-level:0
\[
\in (0,1]:\ V'-V\ \text{continuous in}\ (\operatorname{spec}\rho,\gamma)\ \text{on the faithful simplex}\times(0,1]\ [\log,\exp,\ \text{rational maps}];
\]

%dsp-level:0
\[
\ \
\]

%dsp-level:0
\[
\le 0\ \text{on a dense rational set}\implies
\]

%dsp-level:0
\[
\le 0\ \text{everywhere}\ [\text{math}]
\]

%dsp-level:0
\[
V'=V\iff\tilde K=0:\ (\Leftarrow)\ (1);
\]

%dsp-level:0
\[
\qquad (\Rightarrow)\ \text{near vacuum by}\ (3),(4)\ \text{strictness}\ (\gamma>0);\ \ \text{sweep points strict off-vacuum}\ [\text{machine}].\ \blacksquare
\]

\textbf{Phys.} The pair state carries a natural nonnegative stability functional: the trace-square of its centered surprisal. Ledger ticks rotate it without changing it — the Haar conjugation of A6 is spectral, not contractive. The Gibbs tilt — a separate operation from the tick, determined by $\varepsilon_q$ — strictly lowers it, by an exact per-eigenvalue deficit $\gamma\phi_\gamma$, driving every faithful state toward the A5 vacuum. The tilt strength $\gamma=1-4\varepsilon_q$ is the normalized distance from the crystallization threshold to the vacuum: at $\varepsilon_q=0$ (sharp crystallization) the tilt is at maximum strength; at the vacuum threshold $\varepsilon_q=\tfrac14$ it vanishes. The tilt is grounded in the ledger–tick coupling [L0]: it is applied once per A4 crystallization event, and the ledger counts the steps.

\end{theorem}

\begin{theorem}[{\texorpdfstring{L2 — The temporal decay law}{L2 - The temporal decay law}}]\label{tag:L2}
\textbf{Thm.} Under L1's setting, per tilt step [D0.3; L1(3)], to exact leading order:

%dsp-level:0
\[
V'-V=-\Bigl[1-\bigl(1-\tfrac\gamma4\bigr)^{2}\Bigr]V+O(\delta\rho^3)=-\frac{\gamma(8-\gamma)}{16}\,V+O(\delta\rho^3);\qquad
V'=\bigl(1-\tfrac\gamma4\bigr)^{2}V+O(\delta\rho^3);
\]

%dsp-level:0
\[
\gamma=1:\ \ V'-V=-\tfrac7{16}\,V+O(\delta\rho^3),\ \ V'=\tfrac9{16}V+O(\delta\rho^3);\qquad
\gamma\to0:\ \ V'=V\ \ (\text{identity-tilt limit});
\]

%dsp-level:0
\[
\mu'=\bigl(1-\tfrac\gamma4\bigr)\mu+O(\mu^2)\quad[\text{eigenvalue contraction}];
\]

%dsp-level:0
\[
\qquad \gamma=1:\ \ \mu'=\tfrac34\mu;
\]

%dsp-level:0
\[
\qquad 1-\bigl(1-\tfrac\gamma4\bigr)^2=\frac{\gamma(8-\gamma)}{16}\ \ [\text{exact rationals}].
\]

\textbf{Req.} A1; A4; A5; D0.3, D0.4; CT0.8 (Part 1); L1; exact formal-series arithmetic [math; machine: \texttt{checks/\allowbreak{}verify\_\allowbreak{}L\_\allowbreak{}lyapunov.py} §B].

\textbf{Pf.}

%dsp-level:0
\[
K=-\log\bigl(\tfrac14I+\delta\rho\bigr)=\ln4\,I-4\,\delta\rho+O(\delta\rho^2);\qquad \mathfrak T(K\rho)=\ln4+O(\delta\rho^2)\ \ [\mathfrak T\delta\rho=0]
\]

%dsp-level:0
\[
\implies\tilde K
\]

%dsp-level:0
\[
= K-\mathfrak T(K\rho)I
\]

%dsp-level:0
\[
= -4\,
\]

%dsp-level:0
\[
\delta\rho+O(\delta\rho^2)
\]

%dsp-level:0
\[
\implies V
\]

%dsp-level:0
\[
= 16\,
\]

%dsp-level:0
\[
\mathfrak T(\delta\rho^2)+O(\delta\rho^3)
\]

%dsp-level:0
\[
= 16\textstyle\sum_i\mu_i^2+O(\delta\rho^3)\ \ [\gamma\text{-free}]
\]

%dsp-level:0
\[
\mathrm{L1}(3):\ V'-V=-\gamma\textstyle\sum_i\phi_\gamma(\mu_i)+O(\delta\rho^4);\qquad \gamma\phi_\gamma(x)=\gamma(8-\gamma)x^2+O(x^3)
\]

%dsp-level:0
\[
\implies V'-V
\]

%dsp-level:0
\[
= -\gamma(8-\gamma)\textstyle\sum_i\mu_i^2+O(\delta\rho^3)
\]

%dsp-level:0
\[
= -\frac{\gamma(8-\gamma)}{16}\,
\]

%dsp-level:0
\[
V+O(\delta\rho^3)
\]

%dsp-level:0
\[
= -\Bigl[1-\bigl(1-\tfrac\gamma4\bigr)^2\Bigr]V+O(\delta\rho^3)
\]

%dsp-level:0
\[
\lambda'=\lambda e^{-\gamma\lambda}/Z_\gamma:\quad\text{series}:\ \ \mu'=\bigl(1-\tfrac\gamma4\bigr)\mu+O(\mu^2)\ [\text{exact coefficient; machine}];
\]

%dsp-level:0
\[
\qquad V'=16\bigl(1-\tfrac\gamma4\bigr)^2\textstyle\sum_i\mu_i^2+O(\delta\rho^3)
\]

%dsp-level:0
\[
\gamma=1:\ \ 1-\tfrac14=\tfrac34,\ \ \bigl(\tfrac34\bigr)^2=\tfrac9{16},\ \ 1-\tfrac9{16}=\tfrac7{16},\ \ \gamma(8-\gamma)=7;\qquad
\gamma\to0:\ \ \bigl(1-\tfrac\gamma4\bigr)^2\to1.\ \blacksquare
\]

\textbf{Phys.} Near the vacuum, every tilt step contracts the Lyapunov function by a $\gamma$-dependent exact rational factor: $\bigl(1-\tfrac\gamma4\bigr)^2$ of $V$ survives, $\tfrac{\gamma(8-\gamma)}{16}$ of it is lost. At $\gamma=1$ ($\varepsilon_q=0$, sharp crystallization) the contraction is $\tfrac9{16}$ and the deficit exactly $\tfrac7{16}$ of $V$; at the vacuum threshold ($\gamma\to0$) the contraction is $1$ and the deficit $0$. The $\tfrac7{16}$ is not a stipulation: it is the decay rate at sharp crystallization, derived from $\varepsilon_q=0$ [D0.4]. The tilt is grounded in the ledger–tick coupling [L0]: one step per recorded crystallization event, so the decay law is the per-increment dynamics of the ledger.

\end{theorem}

\begin{theorem}[{\texorpdfstring{L3 — Entropy deficit and the inertial mass reading}{L3 - Entropy deficit and the inertial mass reading}}]\label{tag:L3}
\textbf{Thm.} Under L1, L2: $(1)$ \textbf{Entropy deficit.} With $\delta\rho=\tfrac14\sum_ar_aG_a$ on the Pauli-tensor basis $\{G_a\}=\{\sigma_\mu\otimes\sigma_\nu\}_{(\mu,\nu)\ne(0,0)}$ [C1; $\mathfrak T(G_aG_b)=4\delta_{ab}$: CT1.1], $|\delta\rho|^2:=\mathfrak T(\delta\rho^2)$:

%dsp-level:0
\[
\Delta S_{\mathrm{pair}}:=2\,|\delta\rho|^2=\tfrac12\,|r|^2;
\]

%dsp-level:0
\[
\qquad S\bigl(\rho\,\big\|\,\tfrac14I\bigr)=\Delta S_{\mathrm{pair}}+O(\delta\rho^3)
\]

%dsp-level:0
\[
\qquad \bigl[S(\rho\|\sigma):=\mathfrak T\bigl(\rho(\log\rho-\log\sigma)\bigr)\bigr]
\]

$(2)$ \textbf{Lyapunov identification.} $\ \tilde K=-4\,\delta\rho+O(\delta\rho^2);\qquad V(\tilde K)=16\,|\delta\rho|^2+O(\delta\rho^3)=8\,\Delta S_{\mathrm{pair}}+O(\delta\rho^3).$ $(3)$ \textbf{Deficit decay ($\gamma$-dependent).} $\ V'-V=-\dfrac{\gamma(8-\gamma)}{2}\,\Delta S_{\mathrm{pair}}+O(\delta\rho^3);\qquad\gamma=1:\ \ V'-V=-\tfrac72\,\Delta S_{\mathrm{pair}}+O(\delta\rho^3).$ $(4)$ \textbf{Inertial read.} $\ m_i^2(q):=\Delta S(q)\quad[\Delta S(q):\ \text{the patch deficit, L4};\ \text{dimensionless, nats};\ \neg\text{a particle identification}].$

\textbf{Req.} A1; A5; C1; D0.3; CT1.1 (Part 1); L1, L2; exact formal-series arithmetic, relative-entropy eigenvalue identity [math; machine: \texttt{checks/\allowbreak{}verify\_\allowbreak{}L\_\allowbreak{}lyapunov.py} §B].

\textbf{Pf.}

%dsp-level:0
\[
(1):\ \delta\rho\ \text{Hermitian}\wedge\ \mathfrak T\delta\rho
\]

%dsp-level:0
\[
= 0\implies\delta\rho
\]

%dsp-level:0
\[
= \tfrac14\sum_ar_aG_a,
\]

%dsp-level:0
\[
\ \ r_a
\]

%dsp-level:0
\[
= \mathfrak T(\delta\rho\,G_a)\in\mathbb R\ \ [\mathrm{C1},\ \mathrm{CT1.1}:\ \mathfrak T(G_aG_b)=4\delta_{ab}]
\]

%dsp-level:0
\[
|\delta\rho|^2=\mathfrak T(\delta\rho^2)=\tfrac1{16}\sum_{ab}r_ar_b\,\mathfrak T(G_aG_b)=\tfrac14|r|^2\implies2|\delta\rho|^2=\tfrac12|r|^2
\]

%dsp-level:0
\[
S\bigl(\rho\|\tfrac14I\bigr)
\]

%dsp-level:0
\[
= \mathfrak T(\rho\log\rho)+\ln4
\]

%dsp-level:0
\[
= \textstyle\sum_i\bigl(\tfrac14+\mu_i\bigr)\log(1+4\mu_i)\ [\text{exact eigenvalue identity}]
\]

%dsp-level:0
\[
= 2\textstyle\sum_i\mu_i^2+O(\mu^3)
\]

%dsp-level:0
\[
= \Delta S_{\mathrm{pair}}+O(\delta\rho^3)
\]

%dsp-level:0
\[
(2):\ \mathrm{L2}\ \text{Pf}:\ \tilde K=-4\,\delta\rho+O(\delta\rho^2);\qquad V=16\,\mathfrak T(\delta\rho^2)+O(\delta\rho^3)=8\,\Delta S_{\mathrm{pair}}+O(\delta\rho^3)\ \ [(1)]
\]

%dsp-level:0
\[
(3):\ \mathrm{L2}:\ V'-V
\]

%dsp-level:0
\[
= -\tfrac{\gamma(8-\gamma)}{16}V+O(\delta\rho^3)
\]

%dsp-level:0
\[
= -\tfrac{\gamma(8-\gamma)}{16}\cdot8\,
\]

%dsp-level:0
\[
\Delta S_{\mathrm{pair}}+O(\delta\rho^3)
\]

%dsp-level:0
\[
= -\tfrac{\gamma(8-\gamma)}{2}\Delta S_{\mathrm{pair}}+O(\delta\rho^3)
\]

%dsp-level:0
\[
\gamma=1:\ \ \tfrac{\gamma(8-\gamma)}{2}=\tfrac72
\]

%dsp-level:0
\[
(4):\ \text{definition};\qquad\text{a read},\ \neg\text{an identification}.\ \blacksquare
\]

\textbf{Phys.} The pair state's departure from the vacuum is one object under three exact leading reads: a squared Bloch radius, an entropy deficit against the maximally mixed pair, and the Lyapunov functional itself — locked by exact rational factors that depend on the tilt strength $\gamma=1-4\varepsilon_q$. The inertial read names the deficit $m_i^2(q)=\Delta S(q)$: mass-as-missing-information per channel, dimensionless in nats — a reading, not a particle claim; no pole-mass identification is asserted. The tilt is grounded in the ledger–tick coupling [L0]; the patch-level deficit $\Delta S(q)$ with its integer degree factor is derived in L4.

\end{theorem}

\begin{theorem}[{\texorpdfstring{L4 — The patch entropy deficit}{L4 - The patch entropy deficit}}]\label{tag:L4}
\textbf{Thm.} Vertex $q$ of finite incidence (theorem-local scope); incident pair channels $\{(q,p)\}$ [A1]; $n(q)=\partial^{\leftarrow}(q)+\partial^{\rightarrow}(q)$ [D0.8]; $n_{\mathcal R}(q)$ [A4]; $n_{\mathcal V}(q):=n(q)-n_{\mathcal R}(q)$ [D0.1]. Patch deficit:

%dsp-level:0
\[
\Delta S(q)\ :=\ \sum_{(q,p)\in\mathcal V}S\bigl(\rho_{qp}\,\big\|\,\tfrac14I\bigr)
\]

%dsp-level:0
\[
\qquad [\text{sum over the virtual incident channels of}\ q;\ \ S(\cdot\|\cdot)\ \text{as in L3}(1)]
\]

$(1)$ \textbf{Channel sum.} $\ \Delta S(q)=\sum_{(q,p)\in\mathcal V}\Delta S_{\mathrm{pair}}(pq)+O(\delta\rho^3)=2\sum_{(q,p)\in\mathcal V}|\delta\rho_{pq}|^2+O(\delta\rho^3)$ [L3(1) per channel]. $(2)$ \textbf{Mean form.} For $n_{\mathcal V}(q)\ge1$, with $\langle\Delta S_{\mathrm{pair}}\rangle_q:=\frac1{n_{\mathcal V}(q)}\sum_{(q,p)\in\mathcal V}\Delta S_{\mathrm{pair}}(pq)$:

%dsp-level:0
\[
\Delta S(q)\
\]

%dsp-level:0
\[
= \ n_{\mathcal V}(q)\,
\]

%dsp-level:0
\[
\langle\Delta S_{\mathrm{pair}}\rangle_q+O(\delta\rho^3)\
\]

%dsp-level:0
\[
= \ 2\,
\]

%dsp-level:0
\[
n_{\mathcal V}(q)\,
\]

%dsp-level:0
\[
\langle|\delta\rho|^2\rangle_q+O(\delta\rho^3)\
\]

%dsp-level:0
\[
= \ \tfrac12\,
\]

%dsp-level:0
\[
n_{\mathcal V}(q)\,
\]

%dsp-level:0
\[
\langle|r|^2\rangle_q+O(\delta\rho^3);
\]

$n_{\mathcal V}(q)=0\implies\Delta S(q)=0$ [empty sum]. $(3)$ \textbf{Patch Lyapunov read.} $\ V_{\mathrm{patch}}(q):=\sum_{(q,p)\in\mathcal V}V(\tilde K_{pq})=8\,\Delta S(q)+O(\delta\rho^3)$; under per-channel tilt steps [D0.3]: $\ V_{\mathrm{patch}}'-V_{\mathrm{patch}}=-\frac{\gamma(8-\gamma)}2\,\Delta S(q)+O(\delta\rho^3)$ [L3(2), (3) summed]. $(4)$ \textbf{Inertial read, grounded.} $\ m_i^2(q)=\Delta S(q)$ [L3(4)] $=2\,n_{\mathcal V}(q)\,\langle|\delta\rho|^2\rangle_q+O(\delta\rho^3)$ — an integer multiplicity $n_{\mathcal V}(q)\in\mathbb Z_{\ge0}$ times the mean channel deficit; the read vanishes to leading order iff $n_{\mathcal V}(q)=0$ or every virtual-channel deviation vanishes.

\textbf{Req.} A1; A4; A5; C1; D0.1, D0.3, D0.8; L3; the finite-incidence scope (theorem-local); exact formal-series arithmetic over $\mathbb Q$, exact integer bookkeeping [math; machine: \texttt{checks/\allowbreak{}verify\_\allowbreak{}L0\_\allowbreak{}ledger\_\allowbreak{}tilt.py} §E].

\textbf{Pf.}

%dsp-level:0
\[
(1):\ S\bigl(\rho\,\|\,\tfrac14I\bigr)=\Delta S_{\mathrm{pair}}+O(\delta\rho^3)\ \text{per channel}\ [\mathrm{L3}(1),\ \text{sharp cubic remainder}];
\]

%dsp-level:0
\[
\qquad \text{finite sum over}\ n_{\mathcal V}(q)\ \text{channels preserves the order}\ [\text{linearity}]
\]

%dsp-level:0
\[
(2):\ n_{\mathcal V}\cdot\tfrac1{n_{\mathcal V}}\textstyle\sum=\sum\ [n_{\mathcal V}\ge1];
\]

%dsp-level:0
\[
\qquad |\delta\rho|^2=\tfrac14|r|^2\implies2\cdot\tfrac14=\tfrac12\ [\mathrm{L3}(1);\ \text{exact rationals}];
\]

%dsp-level:0
\[
\qquad n_{\mathcal V}=0:\ \text{empty sum}=0
\]

%dsp-level:0
\[
(3):\ \textstyle\sum_eV(\tilde K_e)=\sum_e\bigl(8\,\Delta S_{\mathrm{pair}}(e)+O(\delta\rho_e^3)\bigr)=8\,\Delta S(q)+O(\delta\rho^3)\ [\mathrm{L3}(2)];
\]

%dsp-level:0
\[
\textstyle\sum_e\bigl(V'(\tilde K_e)-V(\tilde K_e)\bigr)
\]

%dsp-level:0
\[
= -\tfrac{\gamma(8-\gamma)}2\sum_e\Delta S_{\mathrm{pair}}(e)+O(\delta\rho^3)
\]

%dsp-level:0
\[
= -\tfrac{\gamma(8-\gamma)}2\Delta S(q)+O(\delta\rho^3)\ [\mathrm{L3}(3);\ \text{channel-wise tilt, D0.3}]
\]

%dsp-level:0
\[
(4):\ (2)\ \text{into the L3}(4)\ \text{definition};
\]

%dsp-level:0
\[
\qquad n(q),\ n_{\mathcal R}(q)\ \text{integer counts}\ [\mathrm{D0.8},\ \mathrm{A4}]\implies n_{\mathcal V}(q)\in\mathbb Z_{\ge0}.\ \blacksquare
\]

\textbf{Phys.} The patch-level entropy deficit is the pair deficit counted with its multiplicity: every virtual channel incident to $q$ contributes its own missing information against the A5 vacuum, and a fully crystallized patch contributes zero. This grounds the symbol $\Delta S(q)$ of the inertial read $m_i^2(q)=\Delta S(q)$ [L3(4)]: mass-as-missing-information is an integer degree-deficit factor $n_{\mathcal V}(q)=n(q)-n_{\mathcal R}(q)$ times the mean channel deficit at leading order. A reading, not an identification.

\end{theorem}

\begin{theorem}[{\texorpdfstring{L5 — The second-order sector flow of the tilt}{L5 - The second-order sector flow of the tilt}}]\label{tag:L5}
\textbf{Thm.} Under L1's setting, with $c:=1-\tfrac\gamma4$, $\ d:=\tfrac{\gamma(8-\gamma)}8$ [exact rationals in $\gamma$; L2]: $(1)$ \textbf{Quadratic flow.} With $\delta\rho':=\tilde{\mathcal T}_\gamma(\rho)-\tfrac14I$ [D0.3]:

%dsp-level:0
\[
\delta\rho'\ =\ c\,\delta\rho\ -\ d\,\Bigl(\delta\rho^2-\tfrac14|\delta\rho|^2E_0\Bigr)\ +\ O(\delta\rho^3);\qquad
\mathfrak T\Bigl(\delta\rho^2-\tfrac14|\delta\rho|^2E_0\Bigr)=0;\qquad
\delta\rho^2\succeq0 .
\]

$(2)$ \textbf{Single-sector landing.} $\delta\rho\in\mathcal E_{st}$ (a single $(S,T)$ sector [C3; CT1.2]) $\implies S(\delta\rho^2)=\delta\rho^2\ \wedge\ T(\delta\rho^2)=\delta\rho^2:\ \delta\rho^2\in\mathcal E_{++}$; the traceless part lies in $\operatorname{span}\{\Sigma,\hat E^{ij}\}$ [CT1.2, CT1.10]; the square channel of $\mathcal E_{++}$ closes in $\mathcal E_{++}$ — the quadratic transport inside the real block is one-directional. $(3)$ \textbf{Partition function.} $\ \log Z_\gamma=-\tfrac\gamma4-d\,|\delta\rho|^2+O(\delta\rho^3)=-\tfrac\gamma4-\tfrac{\gamma(8-\gamma)}{16}\Delta S_{\mathrm{pair}}+O(\delta\rho^3)$ [Cor L0a; L3(1)]; $\ \gamma=1:\ \log Z_1=-\tfrac14-\tfrac78|\delta\rho|^2=-\tfrac14-\tfrac7{16}\Delta S_{\mathrm{pair}}+O(\delta\rho^3)$. $(4)$ \textbf{Modular trace deficit.} $\ \mathfrak T(K'-K)=\gamma+4\log Z_\gamma$ [Cor L0a] $=-4d\,|\delta\rho|^2+O(\delta\rho^3)=-\tfrac{\gamma(8-\gamma)}4\Delta S_{\mathrm{pair}}+O(\delta\rho^3)\le0$ near the vacuum; $\ \gamma=1:\ -\tfrac74\Delta S_{\mathrm{pair}}+O(\delta\rho^3)$. $(5)$ \textbf{Sharp-crystallization table ($\gamma=1$).} $\ c=\tfrac34;\quad d=\tfrac78;\quad1-c^2=\tfrac7{16}\ [\mathrm{L2}];\quad\tfrac d2=\tfrac7{16};\quad2d=\tfrac72\quad[\text{exact rationals}]$.

\textbf{Req.} A1; A5; C1, C3; D0.3, D0.4; CT0.8; CT1.1, CT1.2, CT1.10; L0 (Cor L0a); L1, L2, L3; exact formal-series arithmetic over $\mathbb Q$, exhaustive square-channel census over $\mathbb Q(i)$ (all $256$ basis pairs) [math; machine: \texttt{checks/\allowbreak{}verify\_\allowbreak{}L0\_\allowbreak{}ledger\_\allowbreak{}tilt.py} §F].

\textbf{Pf.}

%dsp-level:0
\[
(1):\ \rho\,e^{-\gamma\rho}=e^{-\gamma/4}\bigl(\tfrac14E_0+\delta\rho\bigr)e^{-\gamma\,\delta\rho}\ [\text{functional calculus on a commuting family}]
\]

%dsp-level:0
\[
=e^{-\gamma/4}\Bigl[\tfrac14E_0+c\,\delta\rho-\gamma\bigl(1-\tfrac\gamma8\bigr)\delta\rho^2+O(\delta\rho^3)\Bigr],
\]

%dsp-level:0
\[
\qquad \gamma\bigl(1-\tfrac\gamma8\bigr)=\tfrac{\gamma(8-\gamma)}8=d\ [\text{exact rationals}]
\]

%dsp-level:0
\[
Z_\gamma=e^{-\gamma/4}\bigl[1-d\,|\delta\rho|^2+O(\delta\rho^3)\bigr]\ [\mathfrak T\delta\rho=0,\ \mathfrak T\delta\rho^2=|\delta\rho|^2];
\]

%dsp-level:0
\[
\qquad \rho'=\bigl(\tfrac14E_0+c\,\delta\rho-d\,\delta\rho^2\bigr)\bigl(1+d\,|\delta\rho|^2\bigr)+O(\delta\rho^3)
\]

%dsp-level:0
\[
\implies\delta\rho'
\]

%dsp-level:0
\[
= c\,
\]

%dsp-level:0
\[
\delta\rho-d\bigl(\delta\rho^2-\tfrac14|\delta\rho|^2E_0\bigr)+O(\delta\rho^3)\ [\text{coefficient-wise machine:}\ \gamma\in\{\tfrac14,\tfrac12,\tfrac34,1\}\times5\ \text{exact rational families}];
\]

%dsp-level:0
\[
\mathfrak T\bigl(\delta\rho^2-\tfrac14|\delta\rho|^2E_0\bigr)=|\delta\rho|^2-|\delta\rho|^2=0;\qquad
\delta\rho^2\succeq0\ [\text{Hermitian square}]
\]

%dsp-level:0
\[
(2):\ X^2
\]

%dsp-level:0
\[
= \textstyle\sum_ax_a^2G_a^2+\sum_{a<b}x_ax_b\{G_a,G_b\}:\ \text{the commutator channel cancels in squares}\ [\text{ordered double sum}]
\]

%dsp-level:0
\[
\text{square channel}:\ \{G_a,G_b\}
\]

%dsp-level:0
\[
\ne 0
\]

%dsp-level:0
\[
\implies T(G_aG_b)
\]

%dsp-level:0
\[
= \vartheta_a\vartheta_b\,
\]

%dsp-level:0
\[
G_aG_b\ [\text{exhaustive: all}\ 136\ \text{nonvanishing pairs of the}\ 256;\ \text{machine}];
\]

%dsp-level:0
\[
\text{single sector}\implies\vartheta_a\vartheta_b=t^2=+1;\qquad
S=\operatorname{Ad}(\mathbb S)\ \text{inner}\implies S(G_aG_b)=S(G_a)S(G_b)=s^2G_aG_b\ [\mathrm{CT0.7}(2)]
\]

%dsp-level:0
\[
\implies T(X^2)=X^2\ \wedge\ S(X^2)=X^2\implies X^2\in\mathcal E_{++}\ [\mathrm{CT1.1}];
\]

%dsp-level:0
\[
\qquad \mathcal E_{++}=\operatorname{span}\{E_0,\Sigma,\hat E^{ij}\}\ [\mathrm{CT1.2},\ \mathrm{CT1.10}]\implies\text{traceless part}\subseteq\operatorname{span}\{\Sigma,\hat E^{ij}\}
\]

%dsp-level:0
\[
\text{one-directionality}: (+,+)\cdot(+,+)\ \text{square channel}
\]

%dsp-level:0
\[
\mapsto (+,+):\ \text{the same census restricted to}\ \mathcal E_{++}\ [\text{machine}];
\]

%dsp-level:0
\[
\qquad \text{the}\ \neg\text{multiplicativity witness of CT1.11a}(1)\ \text{lies on the commutator channel}\ (\{G_a,G_b\}=0)\ \text{and is reproduced}
\]

%dsp-level:0
\[
(3):\ \log Z_\gamma
\]

%dsp-level:0
\[
= -\tfrac\gamma4+\log\bigl(1-d|\delta\rho|^2+O(\delta\rho^3)\bigr)
\]

%dsp-level:0
\[
= -\tfrac\gamma4-d|\delta\rho|^2+O(\delta\rho^3)\ [\log(1+u)=u+O(u^2)]
\]

%dsp-level:0
\[
\Delta S_{\mathrm{pair}}=2|\delta\rho|^2\ [\mathrm{L3}(1)];\qquad\tfrac d2=\tfrac{\gamma(8-\gamma)}{16};\qquad\gamma=1:\ d=\tfrac78,\ \tfrac d2=\tfrac7{16}
\]

%dsp-level:0
\[
(4):\ \mathfrak T(K'-K)
\]

%dsp-level:0
\[
= \gamma+4\log Z_\gamma\ [\text{Cor L0a}]
\]

%dsp-level:0
\[
= \gamma+4\bigl(-\tfrac\gamma4-d|\delta\rho|^2\bigr)+O(\delta\rho^3)
\]

%dsp-level:0
\[
= -4d|\delta\rho|^2+O(\delta\rho^3)
\]

%dsp-level:0
\[
-4d=-\tfrac{\gamma(8-\gamma)}2;
\]

%dsp-level:0
\[
\qquad -4d|\delta\rho|^2=-2d\,\Delta S_{\mathrm{pair}}=-\tfrac{\gamma(8-\gamma)}4\Delta S_{\mathrm{pair}};
\]

%dsp-level:0
\[
\qquad \gamma=1:\ -\tfrac74\Delta S_{\mathrm{pair}};
\]

%dsp-level:0
\[
\qquad \le0\ \text{near the vacuum}\ [\Delta S_{\mathrm{pair}}\ge0,\ \text{zero only at}\ \tfrac14I:\ \mathrm{L1}(1),\ \mathrm{L3}(1)]
\]

%dsp-level:0
\[
(5):\ \text{exact Fraction table}\ [\text{machine}].\ \blacksquare
\]

\textbf{Phys.} To second order the tilt does two things: it shrinks the deviation uniformly by $c=1-\tfrac\gamma4$, and it transports the traceless square of the deviation with exact weight $d=\tfrac{\gamma(8-\gamma)}8$ into the $S$-even, $T$-even real block — a quadratic sector flow that lands wholly in $\mathcal E_{++}$ for any single-sector input and stays there under further squaring. At sharp crystallization the transport weight is exactly $\tfrac78$, the partition function reads $-\tfrac14-\tfrac7{16}\Delta S_{\mathrm{pair}}$, and the modular trace deficit reads $-\tfrac74\Delta S_{\mathrm{pair}}$ — one exact rational family with the Lyapunov deficit $\tfrac7{16}$ [L2]. A reading, not an identification.

\end{theorem}

\phantomsection
\section[Closure exports of Part 1 §F (consumed by Part 2)]{\texorpdfstring{Closure exports of Part 1 §F (consumed by Part 2)}{Closure exports of Part 1 SF (consumed by Part 2)}}

%dsp-level:0
\[
\begin{array}{lll}
\mathrm{CT0.1}\ \text{dichotomy} & \mathrm{CT0.5}\ D_{\mathrm{tick}}(n)=2^{n+1} & \mathrm{CT1.1}\ (7,3,3,3)\\
\mathrm{CT0.2}\ \rho_\circ,\ \ln4 & \mathrm{CT0.6(a)}\ C_{\mathrm{cap}}=16,\ O(n);\ \delta h\le4\,\delta\tau & \mathrm{CT1.2}\ \text{basis}\ \{E_0,\Sigma,W_i,J_i,K_i,\hat E^{ij}\}\\
\mathrm{CT0.3}\ J^\pm\ \text{increments} & \mathrm{CT0.7}\ \Pi_{\mathcal R},\Pi_{\mathcal V};\ \mathbf F_{(\pm,\cdot)}=\mathcal E_{st} & \mathrm{CT1.3}\text{–}\mathrm{CT1.9}\ \text{closure table};\ \mathfrak u(4)=\mathfrak u(1)\oplus\mathfrak{su}(4)\\
\mathrm{CT0.4}\ \text{real}=\text{registered} & \mathrm{CT0.8}\ 0\le\varepsilon_q<\tfrac14 & \mathrm{CT1.10}\ \mathcal E_{++}=\mathbf1\oplus\mathbf5\oplus\mathbf1\\
\mathrm{CT0.9}\ \Phi=\{0,\pi\};\ 1\ \text{edge}=1\ \text{bit} & \mathrm{CT0.11(a)}\ n\le3;\ \mathbb C^3\ \text{slots};\ \mathrm{CT0.11b}\ n=2d\ \text{(bidirectional)} & \mathrm{CT1.11}\ E_0\nparallel\Sigma\\
\mathrm{CT0.10}\ (P,\pi)\ \text{independent} & \mathrm{CT0.12}\ 3{+}3{+}5{+}3{+}1{+}1=16 & \mathrm{CT0.13}\ \mathcal A(n)\\
\mathrm{CT0.14}\ \dim\mathcal K_n=(6,16,17,25) & &
\end{array}
\]

$\blacksquare$

\phantomsection
\section[§G — Master index of results]{\texorpdfstring{§G — Master index of results}{SG - Master index of results}}

%tbl-level:3
{\scriptsize\setlength{\tabcolsep}{4pt}
\begin{longtable}{>{\raggedright\arraybackslash}p{\dimexpr0.1834\textwidth-2\tabcolsep\relax}|>{\raggedright\arraybackslash}p{\dimexpr0.7540\textwidth-2\tabcolsep\relax}|>{\raggedright\arraybackslash}p{\dimexpr0.0626\textwidth-2\tabcolsep\relax}}
\toprule
Tags & Content & Locus \\
\addlinespace[2pt]
\midrule
\endfirsthead
\toprule
Tags & Content & Locus \\
\addlinespace[2pt]
\midrule
\endhead
\bottomrule
\endlastfoot
A1–A9 · C1–C13 · D0.1–D4.3 · P1–P5 & axioms · conventions · the definition layer · calibration & Part 1 §A–§E \\
CT0.1–CT0.14 (+CT0.6a, CT0.7a, CT0.11a/b) & dichotomy, vacuum, seals, tick load $2^{n+1}$, $C_{\mathrm{cap}}=\allowbreak 16$, swap grading, $\varepsilon_q$ range, $\Phi=\allowbreak \{0,\pi\}$, $n\le\allowbreak 3$, $\mathcal A(n)$, mask blocks & Part 1 §F \\
CT1.0–CT1.11 (+CT1.11a/b) & signature classification, $(S,\allowbreak T)$ sector decomposition, generator basis, $\mathfrak{su}(2)$ closures, commutator tables, global algebra, involution table, native Clifford frames and the weak-covariance pincer & Part 1 §F \\
L0–L5 (+Cor L0a) & ledger–tilt coupling, tilt modular update, Lyapunov function, decay law $\gamma(8-\gamma)/16$, entropy deficit and the inertial read, patch deficit, second-order sector flow & Part 1 §F \\
CT1.12–CT1.29 (+CT1.12a, CT1.19a/b, CT1.29a) · CT2.1–CT2.11 & metric emergence, conformal rescaling, linearized spin-2, continuum limit, effective action, $T_{\mu\nu}$, field equations, conservation, $\Lambda=\allowbreak D^+$, entropy/calibration, one free scale & Part 2 §C \\
EW.1–EW.36 (+EW.5a, EW.36a) & electroweak: masks, $Q=\allowbreak T_3+\allowbreak Y/2$, abelian $\mathfrak u(1)$, scalar trace mode, rung-3 wrapping, $m_W$, neutral mass matrix, $\theta_W$, $m_W=\allowbreak m_Z\cos\theta_W$, $e=\allowbreak g\sin\theta_W$, spectrum & Part 2 §D \\
EM.1–EM.26 (+EM.9a/b, Def-EM.20–22) & electromagnetic: rung-2 $\mathfrak u(1)$, photon masslessness/polarizations/entropy, end-charge, quantization, conservation, $G(e)$, dynamic per-edge phase, alignment patterns and parity, CPT bridge, spectrum & Part 2 §E \\
CT3.1–CT3.22 (+CT3.10a, CT3.17a, CT3.22a) & colour: $N=\allowbreak 3$ forced, $\mathfrak u(3)=\allowbreak \mathfrak{su}(3)\oplus\mathfrak u(1)$, rung-3 algebra, $f_{abc}$, gluons $=\allowbreak 8$, masslessness, baryon/meson singlets, sealing grammar, confinement & Part 2 §F \\
CT4.1–CT4.13 (+CT4.1a, CT4.4a, CT4.5a, CT4.7a) & saturation: $O(n)\to D^-$, absorbing lock, $w_{D^-}$, dark contrast, entropy triangle, bounds, global closure, falsifiers, audit & Part 2 §G \\
Def-H.1–Def-H.43 · FT.1–FT.145 (structural) · FT.153–FT.155 & the fermion layer: measurement primitives, the fermion definition, floors, spin/stability, flicker, spectra, Clifford/branches, path integrals, family channels, carrier/parity, rung one; the interface definitions and their derived interface theorems & Part 3 \\
Def-BM.1–2 · BM.1–BM.6 & colour-sealed composites: baryons/mesons, junction grammar, observability, the forced family group, invariance laws, census consumption, construction specs & Part 4 \\
the 25 Category-B cells & the computed data cells and their consumption block & Census \\
\end{longtable}
}

$\blacksquare$

